% !TeX program = xelatex
\documentclass[11pt,a4paper,fontset=fandol]{ctexart}
\usepackage[margin=18mm]{geometry}
\usepackage{amsmath,amssymb,bm,enumitem,hyperref}
\hypersetup{hidelinks,pdftitle={工科数学分析逐题训练}}
\setlist[enumerate]{label=（\arabic*）,leftmargin=2.6em,itemsep=0.25em}
\newcommand{\dd}{\,\mathrm{d}}
\newenvironment{exercise}[2]{%
  \par\medskip\noindent\textbf{第#1题}\quad\textit{#2}\par\smallskip
}{\par\medskip}
\allowdisplaybreaks
\setcounter{secnumdepth}{-1}
\sloppy
\setlength{\emergencystretch}{4em}
\title{工科数学分析逐题训练\\[0.4em]\large 计算、论证与条件辨析}
\author{}
\date{2026年10月20日}
\begin{document}
\maketitle
\section*{使用说明}
本题册覆盖向量与空间几何、多元微分、重积分、曲线积分、曲面积分及无穷级数，难度按工科数学分析期末考试控制。练习既要求完成计算，也关注条件是否充分、定理能否使用、参数变化对结论的影响。

全册共818道题，按阅读顺序统一编号为1至818。34节各设两道递进题，纳入连续编号；每节先完成基本训练，再练习条件辨析与参数综合。

题目所需函数、区域、参数与方向均在题干中给出，可独立使用。常规版与iPad横屏版共用同一份题目；iPad版每页一道题，下方留白用于草稿。未给出本校考试大纲，具体章目仍以授课要求为准。

本册仅包含改编题目，未附答案；按章节练习使用，不代表某一套试卷的题量或时间安排。证明题可使用工科数学分析课堂已学定理，使用时应说明关键条件。
\tableofcontents
\clearpage
\section{第8章 向量代数与空间解析几何}
\subsection{向量及其线性运算}
\begin{exercise}{1}{线性组合与线性无关}
设 $\vec u=\vec a-\vec b+\vec c$，$\vec v=\vec a+3\vec b-2\vec c$，其中 $\vec a,\vec b,\vec c$ 线性无关。用它们表示 $2\vec u-\vec v$，并判断是否存在实数 $s,t$ 使 $s\vec u+t\vec v=\vec c$，说明理由。
\end{exercise}
\begin{exercise}{2}{坐标运算与正交条件}
设 $\vec a=(0,-1,1)$，$\vec b=(1,2,0)$，$\vec c=(2,0,1)$。求 $\vec p=\vec a-\vec b+\vec c$、$\vec q=\vec a+2\vec b-3\vec c$ 的坐标，并求使 $\vec p+\lambda\vec q$ 与 $\vec c$ 垂直的实数 $\lambda$。
\end{exercise}
\begin{exercise}{3}{三等分点与端点恢复}
线段 $AB$ 的两个三等分点为 $C(2,0,2)$、$D(5,-2,0)$。求无序端点集 $\{A,B\}$ 及线段长度，并说明交换 $C,D$ 的先后次序对结果有何影响。
\end{exercise}
\begin{exercise}{4}{内外分点与参数范围}
设 $A(x_1,y_1,z_1)\ne B(x_2,y_2,z_2)$。对实数 $\lambda$，讨论线段 $AB$ 上满足距离关系 $|AM|=\lambda|MB|$ 的点 $M$ 是否存在、是否唯一；存在时求其坐标。再讨论将 $M$ 的范围改为整条直线 $AB$ 时，正参数 $\lambda$ 对解的个数有何影响。
\end{exercise}
\begin{exercise}{5}{平行条件与投影}
设 $A(1,2,3)$、$B(3,5,t)$，$\vec a=(2,3,-2)$。求使 $\overrightarrow{AB}$ 与 $\vec a$ 平行的 $t$，并求此时 $\overrightarrow{AB}$ 在 $\vec a$ 方向上的标量投影。
\end{exercise}
\begin{exercise}{6}{向量平行的参数判定}
设 $\vec a=(1,2,3)$、$\vec b=(2,k,t)$。求使两向量平行的所有 $(k,t)$，并判断是否可能反向。若仅要求 $\vec b-\vec a$ 与 $\vec a$ 平行，参数条件是否改变？证明你的结论。
\end{exercise}
\begin{exercise}{7}{反向单位向量与定长向量}
设 $A(-1,2,2)$、$B(1,2,0)$。求与 $\overrightarrow{AB}$ 反向的单位向量 $\vec e$，并求满足 $\vec v\parallel\overrightarrow{AB}$、$|\vec v|=3$ 的全部向量。
\end{exercise}
\begin{exercise}{8}{方向角及其约束}
设 $A(1,1,-\sqrt2)$、$B(2,2,0)$。求 $\overrightarrow{AB}$ 的模、方向余弦与方向角，并由本题数据验证方向余弦的平方和为 $1$。能否存在方向角均为锐角且三者相等的非零向量？若能，求共同方向角。
\end{exercise}
\begin{exercise}{9}{标量投影与垂直分解}
向量 $\vec a$ 的模为 $9$，与有向 $u$ 轴的夹角为 $\frac\pi4$。求其在该轴上的标量投影及垂直于该轴的分量的模。若将轴的正方向反转，这两个结果如何改变？
\end{exercise}
\begin{exercise}{10}{等模条件与正交}
设 $\vec a=(3,-5,8)$、$\vec b=(-1,1,p)$。求满足 $|\vec a+\vec b|=|\vec a-\vec b|$ 的 $p$，并证明一般向量的这一等式与 $\vec a\cdot\vec b=0$ 等价。
\end{exercise}
\begin{exercise}{11}{距离比轨迹与退化情形}
设 $A(2,0,-3)$、$B(4,-6,6)$。求满足 $|MA|=3|MB|$ 的动点 $M$ 的轨迹方程、球心与半径。再将距离比 $3$ 改为 $1$，求轨迹并说明其几何类型。
\end{exercise}
\begin{exercise}{12}{向量证明与比例推广}
在非退化三角形 $ABC$ 中，$M,N$ 分别为 $AB,AC$ 的中点。用向量证明 $MN\parallel BC$ 且 $|MN|=\frac12|BC|$。进一步，若 $\overrightarrow{AM}=t\overrightarrow{AB}$、$\overrightarrow{AN}=t\overrightarrow{AC}$，其中 $0<t<1$，求 $\overrightarrow{MN}$ 并证明相应的长度关系。
\end{exercise}
\begin{exercise}{13}{点的线性组合与原点选择}
设 $A,B$ 的位置向量为 $\vec a,\vec b$，以 $\alpha\vec a+\beta\vec b$ 作为点 $P$ 的位置向量。证明这一构造在任意平移坐标原点后仍表示同一几何点的充要条件是 $\alpha+\beta=1$。在此条件下，说明 $P$ 在线段 $AB$ 上的充要条件，其中 $A\ne B$。
\end{exercise}
\begin{exercise}{14}{线段参数与交点条件}
设 $A(0,0,0)$、$B(2,1,1)$、$C(1,0,a)$、$D(1,2,a-2)$。用线性组合分别表示直线 $AB,CD$ 上的点，求两直线相交的参数 $a$；相交时判断交点是否同时在线段 $AB,CD$ 上，并求其对两条线段的分比。
\end{exercise}
\subsection{内积、向量积与混合积}
\begin{exercise}{15}{三种乘积与体积}
设 $\vec a=(1,-2,-3)$，$\vec b=(2,2,-3)$，$\vec c=(2,1,-3)$。计算 $(2\vec a)\cdot\vec b$、$(3\vec a)\times(2\vec b)$ 及 $[(3\vec a)\times(2\vec b)]\cdot\vec c$，并求以 $\vec a,\vec b,\vec c$ 为棱的四面体体积。
\end{exercise}
\begin{exercise}{16}{混合积的线性性质}
设 $(\vec a\times\vec b)\cdot\vec c=2$。计算 $[(\vec a+\vec b)\times(\vec b+\vec c)]\cdot(\vec c+\vec a)$，并求使 $\vec a+\vec b,\vec b+\vec c,\vec c+t\vec a$ 共面的全部实数 $t$。
\end{exercise}
\begin{exercise}{17}{向量积展开与恒等式}
化简 $(2\vec a+\vec b)\times(\vec c-\vec a)+(\vec b+\vec c)\times(\vec a+\vec b)$。若 $\vec a,\vec b,\vec c$ 线性无关，判断该式是否可能为零向量，并证明。
\end{exercise}
\begin{exercise}{18}{内积消去与充要条件}
设 $\vec c\ne\vec0$ 且 $\vec a\cdot\vec c=\vec b\cdot\vec c$。证明 $(\vec a-\vec b)\cdot\vec c=0$，并举例说明能否由此推出 $\vec a=\vec b$。若另有 $\vec a-\vec b$ 与 $\vec c$ 平行，结论如何？
\end{exercise}
\begin{exercise}{19}{共线条件与方向判定}
设 $\vec a,\vec b$ 不共线。求使 $k\vec a+\vec b$ 与 $\vec a+k\vec b$ 共线的全部实数 $k$，并分别判断这两个向量同向还是反向。
\end{exercise}
\begin{exercise}{20}{三角形面积与高}
设 $A(2,-1,2)$、$B(1,1,3)$、$C(1,2,-1)$。用向量积求三角形 $ABC$ 的面积及从 $A$ 到直线 $BC$ 的距离，并说明三个点为何不共线。
\end{exercise}
\begin{exercise}{21}{面积与等面积轨迹}
设 $A(2,-1,3)$、$B(1,2,3)$、$C(0,1,4)$。求三角形 $ABC$ 的面积。保持 $A,B$ 固定，令 $C_t=C+t\overrightarrow{AB}$，证明三角形 $ABC_t$ 的面积与 $t$ 无关。
\end{exercise}
\begin{exercise}{22}{平行约束与最小模}
设 $\vec m=(2,-3,1)$。求与 $\vec m$ 共线且满足 $\vec r\cdot\vec m=14$ 的 $\vec r$。若去掉共线条件，证明所有满足该内积条件的向量中，刚求得的向量的模最小。
\end{exercise}
\begin{exercise}{23}{正交参数与投影系数}
设 $|\vec a|=2$，$|\vec b|=5$，夹角为 $\frac{2\pi}3$，$\vec p=3\vec a-\vec b$、$\vec q=\lambda\vec a+17\vec b$。求使 $\vec p\perp\vec q$ 的 $\lambda$，并求此时 $\vec q$ 在 $\vec a$ 方向上的标量投影。
\end{exercise}
\begin{exercise}{24}{内外积关系与信息充分性}
设 $|\vec a|=3$，$|\vec b|=26$，$|\vec a\times\vec b|=72$。求 $\vec a\cdot\vec b$ 的所有可能值，并说明还需知道夹角的什么信息才能唯一确定内积。
\end{exercise}
\begin{exercise}{25}{正交关系联立与夹角}
非零向量 $\vec a,\vec b$ 满足 $(\vec a+3\vec b)\perp(7\vec a-5\vec b)$ 及 $(\vec a-4\vec b)\perp(7\vec a-2\vec b)$。求 $|\vec a|:|\vec b|$ 与两向量夹角，并核验所得数据满足 $|\cos\theta|\le1$。
\end{exercise}
\begin{exercise}{26}{共面判定与有向体积}
设 $\vec a=(3,0,-1)$，$\vec b=(2,-4,3)$，$\vec c=(-1,-2,2)$。判断三向量是否共面，求它们所张成的平行六面体体积；再求使 $\vec a,\vec b,\vec c+t\vec a$ 共面的实数 $t$，或证明不存在。
\end{exercise}
\begin{exercise}{27}{拉格朗日恒等式与等号}
证明 $|\vec a\times\vec b|^2+(\vec a\cdot\vec b)^2=|\vec a|^2|\vec b|^2$，并由此证明 $|\vec a\cdot\vec b|\le |\vec a||\vec b|$，完整讨论等号成立条件，包括零向量情形。
\end{exercise}
\begin{exercise}{28}{余弦定理及逆命题}
设三角形 $ABC$ 的对边长为 $a,b,c$。用向量内积证明三条余弦定理 $a^2=b^2+c^2-2bc\cos A$、$b^2=a^2+c^2-2ac\cos B$、$c^2=a^2+b^2-2ab\cos C$，再据此证明 $A$ 为锐角的充要条件是 $a^2<b^2+c^2$。
\end{exercise}
\begin{exercise}{29}{三高共点的向量证明}
用向量证明非退化三角形的三条高所在直线共点。证明中说明两条高为何必有唯一交点，并证明该交点也在第三条高上。
\end{exercise}
\begin{exercise}{30}{正交分解与最小距离原理}
设 $\vec b\ne\vec0$。证明函数 $|\vec a-t\vec b|^2$ 在实数 $t$ 上有唯一最小值，求极小点并证明此时 $\vec a-t\vec b$ 与 $\vec b$ 垂直。由最小值非负推导内积不等式，并完整讨论等号条件。
\end{exercise}
\begin{exercise}{31}{混合积的参数讨论}
设 $\vec a=(1,1,0)$、$\vec b=(0,1,1)$、$\vec c_t=(1,t,1)$。求三向量共面的全部 $t$，共面时将 $\vec c_t$ 表为 $\vec a,\vec b$ 的线性组合；再求以这三向量为棱的平行六面体体积不超过 $2$ 的参数范围，并求对应四面体的最大体积。
\end{exercise}
\subsection{平面方程}
\begin{exercise}{32}{点法式与垂足}
求过 $A(1,2,3)$ 且垂直于直线 $\frac{x-2}{2}=\frac{y+1}{4}=\frac{z+2}{5}$ 的平面，并求该直线与平面的交点以及 $A$ 到直线的距离。
\end{exercise}
\begin{exercise}{33}{三点确定平面}
设 $A(2,2,1)$、$B(1,0,3)$、$C(0,3,-2)$。先证明三点不共线，再求所确定的平面方程，并求原点到该平面的距离。
\end{exercise}
\begin{exercise}{34}{平面平行于坐标轴}
求平行于 $y$ 轴且过 $A(1,0,1)$、$B(2,-1,3)$ 的平面方程，说明唯一性；再求该平面与 $xOz$ 坐标面的交线。
\end{exercise}
\begin{exercise}{35}{特殊平面的几何解释}
分别说明平面 $3y-1=0$、$4x+3y-12=0$、$x-3y=0$ 与各坐标轴、坐标面的相对位置；对每个平面给出一个法向量，并判断能否写成三个非零有限截距的截距式。
\end{exercise}
\begin{exercise}{36}{过两点的垂直平面}
求过 $A(2,-1,-1)$、$B(1,3,-4)$ 且垂直于 $x+2y-4z+8=0$ 的平面。证明所求平面唯一，并求这两个平面的交线方向向量。
\end{exercise}
\begin{exercise}{37}{双重垂直条件}
求过 $A(1,2,-1)$ 且同时垂直于 $2x+2y-z+1=0$、$3x-y-z-3=0$ 的平面，并说明其法向量与这两个已知平面的交线方向有何关系。
\end{exercise}
\begin{exercise}{38}{垂直平分面与等距轨迹}
求线段 $A(1,2,3)$、$B(2,-1,4)$ 的垂直平分面，并从距离相等条件证明该平面恰为所有满足 $|MA|=|MB|$ 的点的轨迹。
\end{exercise}
\begin{exercise}{39}{由两方向确定平面}
求过原点且方向平行于直线 $x=1,\ y=-1+t,\ z=2+t$ 与直线 $\frac{x+1}{1}=\frac{y+2}{2}=\frac{z-1}{1}$ 的平面。分别判断两已知直线是否位于所求平面内。
\end{exercise}
\begin{exercise}{40}{截距比例与多解}
求三个有向截距之比为 $1:3:2$ 且到原点距离为 $6$ 的全部平面，并求各平面在第一卦限是否围成有界四面体；能围成时求其体积。
\end{exercise}
\begin{exercise}{41}{过坐标轴的夹角约束}
求通过 $y$ 轴且与 $2x-\sqrt5y+z-7=0$ 成 $\frac\pi3$ 角的全部平面，说明求解时能否遗漏法向量某个分量为零的情形。
\end{exercise}
\begin{exercise}{42}{过两点的平面夹角}
求过 $M_1(0,0,1)$、$M_2(3,0,0)$ 且与 $xOy$ 面成 $60^\circ$ 角的全部平面，并说明两解之间的对称关系。
\end{exercise}
\begin{exercise}{43}{等距平行平面}
已知 $\Pi_1:2x+6y-3z+12=0$，$P(0,2,-1)$。求与 $\Pi_1$ 平行且不重合，并使 $P$ 到两个平面的距离相等的 $\Pi_2$；说明 $P$ 与两个平面的相对位置。
\end{exercise}
\begin{exercise}{44}{平面位置关系与定量计算}
逐组判定两平面的位置关系；相交时求夹角，平行且不重合时求距离：
\begin{enumerate}
\item $2x+y-z+1=0$ 与 $6x+3y-3z+3=0$；
\item $x-y-z+7=0$ 与 $-2x+4y-6z+11=0$；
\item $5x+3y-z+1=0$ 与 $x-6y-3z+2=0$；
\item $2x+y-z+1=0$ 与 $6x+3y-3z+5=0$。
\end{enumerate}
\end{exercise}
\begin{exercise}{45}{平面夹角与交线}
求平面 $x+y-11=0$ 与 $3x+8=0$ 的夹角及交线方程，并解释交线为何平行于 $z$ 轴。
\end{exercise}
\begin{exercise}{46}{坐标面夹角与法向量}
求 $2x+y-4z+5=0$ 与三个坐标面的夹角的余弦，证明这三个余弦的平方和为 $1$，并说明它们与法向量的方向余弦有何联系。
\end{exercise}
\begin{exercise}{47}{点面距离与正交投影}
求 $P(2,0,3)$ 到 $\Pi:4x-5y+3z-1=0$ 的距离、垂足坐标及 $P$ 关于 $\Pi$ 的对称点。
\end{exercise}
\begin{exercise}{48}{平行平面的双侧距离}
设 $A(0,-7,0)$、$B(2,-1,1)$、$C(2,2,2)$。求平行于三角形 $ABC$ 所在平面且与之相距 $2$ 的全部平面，并判定它们分居原平面的哪一侧。
\end{exercise}
\begin{exercise}{49}{切平面与线性近似}
求曲面 $z=x^2(1-\sin y)+y^2(1-\sin x)$ 在 $(1,0,1)$ 处的切平面，并用该切平面给出 $f(1+h,k)$ 的一阶近似，写明误差相对于 $\sqrt{h^2+k^2}$ 的阶。
\end{exercise}
\begin{exercise}{50}{有向截距与四面体体积}
平面 $\frac x2+\frac y4-\frac z3=1$ 与三个坐标面围成四面体。求各轴截距及四面体体积，并说明负截距为何不导致负体积。
\end{exercise}
\begin{exercise}{51}{平行平面距离与中间平面}
求 $3x+6y-2z-7=0$ 与 $3x+6y-2z+14=0$ 的距离，并求与两平面等距且位于两者之间的平面。
\end{exercise}
\begin{exercise}{52}{球面相切与切点}
球心为 $O(3,-5,-2)$ 的球面与 $2x-y-3z+11=0$ 相切。求球面方程与切点，并验证切点到球心的连线垂直于切平面。
\end{exercise}
\begin{exercise}{53}{过直线的平面约束}
设 $A(5,1,3)$、$B(1,6,2)$、$C(5,0,4)$、$D(4,0,6)$。
\begin{enumerate}
\item 求通过 $AB$ 且方向平行于 $CD$ 的平面，判断 $CD$ 是否在该平面内。
\item 求通过 $AB$ 且垂直于三角形 $ABC$ 所在平面的平面，并说明唯一性。
\end{enumerate}
\end{exercise}
\begin{exercise}{54}{平面夹角与角平分面}
对平面 $2x-y+z-7=0$ 与 $x+y+2z-11=0$，求夹角及两个角平分面方程。证明这两个角平分面互相垂直，并说明使用未经单位化的平面方程直接作和差为何需要条件。
\end{exercise}
\begin{exercise}{55}{平面束与唯一性}
设 $L$ 为平面 $x+y+z=0$ 与 $x-y=0$ 的交线。证明通过 $L$ 的全部平面可写成 $\alpha(x+y+z)+\beta(x-y)=0$，其中 $(\alpha,\beta)\ne(0,0)$。讨论另给一点 $P$ 时，通过 $L,P$ 的平面何时唯一，何时有无穷多个。
\end{exercise}
\begin{exercise}{56}{截距平面与体积极值}
一个平面在三个正半轴上截距分别为 $a,b,c>0$，且通过 $P(1,2,3)$。求它与三个坐标面围成的四面体体积的最小值及取得最小值的平面方程，并证明最小值确实能够取得。
\end{exercise}
\subsection{空间直线}
\begin{exercise}{57}{直线参数式与零分量}
求过 $M_1(5,2,-3)$、$M_2(-1,2,3)$ 的直线的参数方程与一般方程，并说明其对称式中应如何处理方向向量的零分量。
\end{exercise}
\begin{exercise}{58}{法线与点面距离}
求过 $P(2,-1,5)$ 且垂直于 $x-5y+3z+1=0$ 的直线，求其与平面的交点，并由此计算 $P$ 到平面的距离。
\end{exercise}
\begin{exercise}{59}{直线方程转换与点的判定}
将直线 $x-y+z-1=0,\ 2x+y+z-4=0$ 写为参数式和对称式，求它与 $xOy$ 面的交点，并验证该点满足原方程组。
\end{exercise}
\begin{exercise}{60}{两平面的公共平行方向}
求过 $P(-1,-3,4)$ 且方向同时平行于 $3x+2y+z+7=0$ 和 $x+2y-z-4=0$ 的直线，分别判定它与两个平面是平行还是包含关系。
\end{exercise}
\begin{exercise}{61}{公共垂直方向与空间位置}
求过 $P(-5,1,8)$ 且方向同时垂直于 $L_1:x-y+3z+2=0,\ 4x+y-2z+3=0$ 和 $L_2:(x,y,z)=(5-2t,4-t,3t+5)$ 的直线。判断所得直线是否分别与 $L_1,L_2$ 相交。
\end{exercise}
\begin{exercise}{62}{直线夹角与交点参数}
求过 $(2,-3,-1)$、与 $x$ 轴成 $\frac\pi3$ 角且与 $L_1:\frac{x-7}{1}=\frac{y-3}{2}=\frac{z-5}{2}$ 相交的全部直线。以交点在 $L_1$ 上的参数建立方程，并说明无增根、无漏解。
\end{exercise}
\begin{exercise}{63}{点到直线的垂足}
求过 $P(2,1,3)$ 且与 $\frac{x+1}{3}=\frac{y-1}{2}=\frac z{-1}$ 垂直相交的直线，并求垂足及 $P$ 到已知直线的距离。
\end{exercise}
\begin{exercise}{64}{垂线与最短距离}
过 $P(0,-1,1)$ 向直线 $L:y+1=0,\ x+2z-7=0$ 作垂线。求垂线方程和垂线段长度，并用一元二次函数的最小值验证该长度。
\end{exercise}
\begin{exercise}{65}{直线夹角与异面判定}
求直线 $\frac{x-1}{1}=\frac y{-4}=\frac{z+3}{1}$ 与 $\frac x2=\frac{y+2}{-2}=\frac z{-1}$ 的夹角，并判定两者相交还是异面。
\end{exercise}
\begin{exercise}{66}{一般式直线的方向}
求 $L_1:x-y+z-1=0,\ 2x+y+z-4=0$ 和 $L_2:x+y-4z+3=0,\ x-y-3z-2=0$ 的夹角，并求同时垂直于两条直线方向的单位向量。
\end{exercise}
\begin{exercise}{67}{线面夹角与平行判定}
求直线 $x+y+3z=0,\ x-y-z=0$ 与平面 $x-y-z+1=0$ 的夹角和距离。解释夹角为零是否足以推出直线位于平面内。
\end{exercise}
\begin{exercise}{68}{方向正交与实际相交}
设 $L_1:(x,y,z)=(-1,1,-3)+s(1,t,2)$，$L_2:\frac{x-2}{4}=\frac{y+1}{2}=\frac{z-2}{-1}$。求使两直线方向垂直的 $t$，再判断此时两直线是否垂直相交。
\end{exercise}
\begin{exercise}{69}{线面位置与交点}
判定 $L:\frac{x-2}{1}=\frac{y-1}{2}=\frac{z+1}{2}$ 与 $\Pi:x-5y+3z+1=0$ 的位置关系；相交时求交点，并求 $L$ 与 $\Pi$ 的夹角。
\end{exercise}
\begin{exercise}{70}{包含关系的充分证据}
判定 $L:\frac{x-3}{2}=\frac{y+1}{-1}=\frac{z-1}{2}$ 与 $\Pi:2x-2y-3z-5=0$ 的位置关系。分别核验方向条件及基点条件，并说明缺少任一条件会造成什么不确定性。
\end{exercise}
\begin{exercise}{71}{参数代入与交点唯一性}
求 $L:\frac{x+2}{2}=\frac{y-3}{-3}=\frac{z-5}{-1}$ 与 $\Pi:3x+y-2z-7=0$ 的交点。再将常数项 $-7$ 改为 $c$，讨论交点个数是否随 $c$ 改变。
\end{exercise}
\begin{exercise}{72}{一般式直线与法向量}
判定 $L:x+3y+2z+1=0,\ 2x-y-10z+3=0$ 与 $\Pi:4x-2y+z-2=0$ 的位置关系，并求其交点或证明无交点。要求用方向向量给出判定依据。
\end{exercise}
\begin{exercise}{73}{平面内的垂直直线族}
求平面 $\Pi:x+y+z+1=0$ 内方向垂直于 $L:y-z+1=0,\ x+2z=0$ 的全部直线，可用一个参数表示直线族；判定这些直线是否都与 $L$ 垂直相交，并说明理由。
\end{exercise}
\begin{exercise}{74}{两类线面位置关系}
逐项判定位置关系，并计算线面距离；相交时还须给出交点：
\begin{enumerate}
\item $L:\frac{x-2}{3}=\frac{y+2}{1}=\frac{z-3}{-4}$，$\Pi:x+y+z-3=0$；
\item $L:x+3y+2z+5=0,\ 2x-y-10z-3=0$，$\Pi:-x+\frac12y-\frac14z+1=0$。
\end{enumerate}
\end{exercise}
\begin{exercise}{75}{线面距离与等距性}
求 $L:\frac{x-1}{-2}=\frac{y+4}{-7}=\frac z3$ 与 $\Pi:4x-2y-2z-3=0$ 的距离，并证明 $L$ 上每个点到 $\Pi$ 的距离相同。
\end{exercise}
\begin{exercise}{76}{直线间距与位置关系}
判定 $L_1:\frac{x-2}{2}=\frac{y+3}{-3}=\frac z1$ 与 $L_2:x=-2y+1,\ z=y-3$ 的位置关系，并求两直线距离及一对实现最短距离的点。
\end{exercise}
\begin{exercise}{77}{点的平面投影与反射}
求 $P(-1,2,0)$ 在 $x+2y-z+1=0$ 上的投影点及点面距离，再求 $P$ 关于该平面的对称点。
\end{exercise}
\begin{exercise}{78}{点的直线投影与极小值}
求 $P(3,-1,2)$ 在 $L:x+y-z+1=0,\ 2x-y+z-4=0$ 上的投影点及点线距离，并证明该投影点是 $L$ 上距 $P$ 最近的唯一点。
\end{exercise}
\begin{exercise}{79}{平行与包含的参数区别}
设 $L:2x+y-1=0,\ 3x+z-2=0$，$\Pi_c:x+2y-z=c$。讨论所有实数 $c$ 下的线面位置关系、交点情况及距离，并特别给出 $c=1$ 时的结果。
\end{exercise}
\begin{exercise}{80}{异面直线公垂线}
设 $L_1:\frac{x-1}{1}=\frac y2=\frac{z+1}{1}$，$L_2:\frac{x+2}{-2}=\frac{y-1}{1}=\frac z1$。先验证两直线异面，再求公垂线方程、两个垂足以及两直线距离，并说明公垂线唯一的原因。
\end{exercise}
\begin{exercise}{81}{直线共面与距离公式的条件}
设 $L_1:P+t\vec u$、$L_2:Q+s\vec v$，其中 $\vec u,\vec v$ 非零且不平行。证明两直线相交的充要条件为 $(Q-P)\cdot(\vec u\times\vec v)=0$，并推导异面时的距离公式。再举例说明去掉不平行假设后，混合积为零为什么不能保证相交。
\end{exercise}
\begin{exercise}{82}{异面距离与参数退化}
设 $L_1:(x,y,z)=(t,0,0)$，$L_2:(x,y,z)=(0,1,a)+s(0,1,1)$。按 $a$ 分类两直线的位置关系，求它们的距离及一对实现此距离的点；异面时求公垂线方程，并核验其两个方向正交条件。
\end{exercise}
\subsection{曲面及其方程}
\begin{exercise}{83}{配方与球面位置}
求球面 $x^2+y^2+z^2-6x+8y+2z+10=0$ 的球心、半径，并判定三个坐标面分别与球面相交、相切还是相离。
\end{exercise}
\begin{exercise}{84}{直径端点与球面方程}
球面一条直径的端点为 $A(2,-3,5)$、$B(4,1,-3)$。求球面方程，并用 $(\overrightarrow{MA}\cdot\overrightarrow{MB})=0$ 说明该球面的几何刻画。
\end{exercise}
\begin{exercise}{85}{四点确定球面}
求过原点及 $A(1,3,0)$、$B(4,0,0)$、$C(0,0,-4)$ 的球面方程。证明这四点不共面，并说明所求球面唯一。
\end{exercise}
\begin{exercise}{86}{斜圆柱面与点线距离}
圆柱面以 $\frac x1=\frac{y-2}{-2}=\frac{z+1}{-2}$ 为轴，且过 $M(1,-2,1)$。求其半径和方程，可保留向量积的等价形式；证明沿轴向平移任意柱面上的点仍在柱面上。
\end{exercise}
\begin{exercise}{87}{二次曲面与截面辨析}
分别识别下列曲面的类型，指出中心或顶点、主要轴向，并写出一个能体现其类型的坐标面截线方程：
\begin{enumerate}
\item $x^2+y^2+z^2-2x+4y+2z=0$；
\item $4x^2+y^2-z^2=4$；
\item $4x^2-y^2-z^2=4$；
\item $x^2-y^2=9$；
\item $\frac{x^2}2+\frac{y^2}3-2z=0$；
\item $\frac{x^2}3-\frac{y^2}2-2z=0$。
\end{enumerate}
\end{exercise}
\begin{exercise}{88}{旋转曲面与参数表示}
将 $xOz$ 面内的圆 $x^2+z^2=9$ 绕 $z$ 轴旋转一周。求所得曲面方程，并给出覆盖整个曲面的双参数表示及参数范围。
\end{exercise}
\begin{exercise}{89}{旋转轴与母线}
说明 $x^2-y^2+z^2=1$ 可由哪个坐标面内的哪条曲线绕哪条坐标轴旋转得到，并给出一组参数方程；再用截面说明它的单叶性质。
\end{exercise}
\begin{exercise}{90}{柱面与缺失变量}
分别说明 $y^2=x$、$x^2+y^2=1$、$\frac{y^2}4+\frac{z^2}9=1$ 在空间中的几何类型、母线方向及一条准线，并给出每个柱面的参数表示。
\end{exercise}
\begin{exercise}{91}{标准化与曲面分类}
将下列方程化为标准形式，识别曲面类型及主要轴向，并判断是否有中心：
\begin{enumerate}
\item $x^2-2y^2+3z^2+6=0$；
\item $x^2+2y^2+3z^2-6=0$；
\item $x^2-y^2-z^2+2=0$；
\item $x^2+4y^2-8z=0$；
\item $z^2-6x^2-24y=0$。
\end{enumerate}
\end{exercise}
\begin{exercise}{92}{给定切点的球面多解}
求半径为 $4$ 且与平面 $x-2y+2z-3=0$ 相切于 $P(3,2,2)$ 的全部球面。先核验 $P$ 在已知平面上，再说明为何恰有两个球心。
\end{exercise}
\begin{exercise}{93}{旋转椭球与切锥体积}
将椭圆 $\frac{x^2}4+\frac{y^2}3=1$ 绕 $x$ 轴旋转得椭球面 $S_1$；将从 $(4,0)$ 向椭圆所作的切线绕 $x$ 轴旋转得圆锥面 $S_2$。
\begin{enumerate}
\item 求 $S_1,S_2$ 的方程及它们的相切圆所在平面。
\item 求由 $S_2$ 从相切圆至顶点的锥面部分及 $S_1$ 在相切圆右侧的曲面部分围成的有界区域的体积。列积分时明确积分区间，并说明为什么要扣除椭球右端的部分。
\end{enumerate}
\end{exercise}
\begin{exercise}{94}{斜圆柱面与坐标缺项}
证明曲面 $(x-z)^2+2y^2=2$ 是圆柱面，求其轴线与半径。说明沿 $(1,0,1)$ 平移为何保持曲面不变，并据此判断“圆柱面的直角坐标方程必缺少一个变量”这一说法是否成立。
\end{exercise}
\begin{exercise}{95}{含交叉项的曲面分类}
对实数 $\lambda$，讨论曲面 $x^2+y^2+z^2+2\lambda xz=1$ 的类型及有界性。可用正交坐标 $u=\frac{x+z}{\sqrt2}$、$v=\frac{x-z}{\sqrt2}$ 化为标准形式；须单独讨论 $\lambda=\pm1$，并给出无界曲面的主要轴向或母线方向。
\end{exercise}

\clearpage
\section{第9章 多元函数微分法及其应用}
\subsection{多元函数的基本概念}
\begin{exercise}{96}{定义域与边界}
求 $z=\ln(x-y)+\frac{\sqrt y}{\sqrt{4-x^2-y^2}}$ 的定义域，分别列出各项的约束，并求定义域的内部与边界，判断定义域是否为开集。
\end{exercise}
\begin{exercise}{97}{复合根式与对数约束}
求 $z=\arccos(2x)+\frac{\sqrt{y-\sqrt x}}{\ln(1-x^2-y^2)}$ 的定义域。说明对数真数为正与分母非零分别排除了哪些点，并判断定义域是否包含其全部边界点。
\end{exercise}
\begin{exercise}{98}{空间定义域与连通性}
求 $u=\sqrt{4-x^2-y^2-z^2}+\frac1{\sqrt{x^2+y^2+z^2}-1}$ 的定义域及边界，判断定义域的有界性、开闭性与连通性，并说明原点是否属于定义域。
\end{exercise}
\begin{exercise}{99}{函数复合与约分条件}
设 $f(u,v)=\frac{uv}{u^2+v}$。求 $g(x,y)=f(xy,\frac xy)$ 的表达式和自然定义域，并说明约分后的代数式能否直接用于原定义域之外的点。
\end{exercise}
\begin{exercise}{100}{变量代换与退化点}
已知对所有 $x\ne0$、$y\in\mathbb R$ 有 $f(x+y,\frac yx)=x^2-y^2$。求映射 $(x,y)\mapsto(x+y,\frac yx)$ 的像集，并在像集上求 $f(u,v)$；特别讨论 $v=-1$ 时的信息是否充分。
\end{exercise}
\begin{exercise}{101}{对数代换与可去延拓}
已知 $x>0$、$x\ne1$、$y\in\mathbb R$ 时
\[
f(x-y,\ln x)=\left(\frac yx-1\right)\frac{e^x}{e^y\ln(x^x)}.
\]
求 $f(u,v)$ 的表达式、由已知关系确定的定义域及 $f(1,1)$。对固定 $u$，讨论 $v\to0$ 时是否存在有限极限。
\end{exercise}
\begin{exercise}{102}{振荡因子的控制与推广}
用 $\varepsilon$--$\delta$ 定义证明 $\lim_{(x,y)\to(0,0)}(x^2+y^2)\sin\frac1{\sqrt{x^2+y^2}}=0$。若将正弦项换为任意有界函数 $g(x,y)$，原结论是否仍成立？证明。
\end{exercise}
\begin{exercise}{103}{路径判别与趋近尺度}
证明下列极限不存在，并分别比较沿直线和沿抛物线趋近所提供的信息：
\[
\lim_{(x,y)\to(0,0)}\frac{xy}{x^2+y^2},\qquad
\lim_{(x,y)\to(0,0)}\frac{x^2y}{x^4+y^2}.
\]
对第二个函数，另证明沿每条过原点的直线趋近时极限都为零。
\end{exercise}
\begin{exercise}{104}{等价无穷小与定义域延拓}
求下列极限，并说明哪些表达式在原分母为零的非孤立点处可以作连续延拓：
\begin{enumerate}
\item $\lim_{(x,y)\to(0,5)}\frac{\sin(xy)}x$；
\item $\lim_{(x,y)\to(0,0)}\frac{x^2+y^2}{|x|+|y|}$；
\item $\lim_{(x,y)\to(0,0)}\frac{xy}{\sqrt{2-e^{xy}}-1}$。
\end{enumerate}
\end{exercise}
\begin{exercise}{105}{径向极限与混合变量}
求 $\lim_{(x,y)\to(0,0)}\frac{\ln(1+x^2+y^2)}{(1-\cos\sqrt{x^2+y^2})(1+e^{xy})}$，并说明为什么可以分别处理径向因子与含 $xy$ 的因子，不能只检验坐标轴。
\end{exercise}
\begin{exercise}{106}{无穷远处的不同趋近方式}
求 $x,y$ 同时趋于 $+\infty$ 时 $\left(1+\frac1{xy}\right)^{x^2/(x+y)}$ 的极限，并给出控制其对数的估计。再固定 $y=b>0$，只令 $x\to+\infty$ 求极限，说明趋近条件对结果的影响。
\end{exercise}
\begin{exercise}{107}{指数衰减与区域约束}
求 $x,y$ 同时趋于 $+\infty$ 时 $(x^2+y^2)e^{-(x+y)}$ 的极限，并作估计证明。若只知道 $x+y\to+\infty$，分别讨论附加条件 $x,y\ge0$ 与允许 $x,y$ 为任意实数时，能否保证相同结论；不能时构造一条反例路径。
\end{exercise}
\begin{exercise}{108}{幂指型极限与对数估计}
求 $\lim_{(x,y)\to(0,0)}(x^2+y^2)^{x^2y^2}$。用 $x^2y^2\le\frac14(x^2+y^2)^2$ 控制其对数，并说明沿坐标轴趋近与一般路径趋近是否一致。
\end{exercise}
\begin{exercise}{109}{自然定义域上的连续性}
求 $f(x,y)=\frac{\sqrt{4x-y^2}}{\ln(1-x^2-y^2)}$ 的自然定义域，说明它在定义域内的连续性，并求 $\lim_{(x,y)\to(1/2,0)}f(x,y)$。判断根式等号边界上属于定义域的点是否连续。
\end{exercise}
\begin{exercise}{110}{连续延拓与可微性区别}
设 $f(x,y)=\frac{xy}{\sqrt{x^2+y^2}}$（非原点），$f(0,0)=a$。求使 $f$ 在原点连续的 $a$，并用定义判断此时两个偏导数及全微分是否存在。
\end{exercise}
\begin{exercise}{111}{振荡项与连续性}
讨论下列函数在原点的连续性。对连续者给出统一估计，对不连续者给出明确的趋近序列或路径：
\begin{enumerate}
\item $f(0,0)=0$，非原点处 $f(x,y)=(x^2+y^2)\ln(x^2+y^2)$；
\item $g(0,0)=0$，非原点处 $g(x,y)=\frac{x^2\sin\frac1{x^2+y^2}+y^2}{x^2+y^2}$。
\end{enumerate}
\end{exercise}
\begin{exercise}{112}{径向振荡与极限条件}
对实数 $\alpha$，设 $F_\alpha(0,0)=0$，非原点处
\[
F_\alpha(x,y)=(x^2+y^2)^{\alpha/2}\sin\frac1{\sqrt{x^2+y^2}}.
\]
完整讨论原点处的极限、连续性及局部有界性。没有极限时，构造两列趋近原点的点说明原因。
\end{exercise}
\begin{exercise}{113}{分子阶数与二重极限}
设 $p,q>0$，$G_{p,q}(0,0)=0$，非原点处 $G_{p,q}(x,y)=\frac{|x|^p|y|^q}{x^2+y^2}$。求使 $G_{p,q}$ 在原点连续的全部 $(p,q)$，并分别讨论原点邻域内有界的充要条件与二重极限存在的充要条件。
\end{exercise}
\subsection{偏导数}
\begin{exercise}{114}{偏导数定义与局部信息}
设 $f$ 在原点邻域有定义，且 $x\ne0$ 时 $\frac{f(x,0)-f(0,0)}x=2x+3$。用定义求 $f_x(0,0)$，并判断这些条件是否足以确定 $f_y(0,0)$，用两个具体函数说明。
\end{exercise}
\begin{exercise}{115}{根式函数的轴向差商}
设 $f(x,y)=e^{\sqrt{x^2+y^4}}$。用定义讨论原点两个偏导数是否存在，并进一步判断原点的连续性及可微性，给出相应理由。
\end{exercise}
\begin{exercise}{116}{偏导数与变步长差商}
设 $f_x(x_0,y_0)$ 存在。分别求
\[
\lim_{h\to0}\frac{f(x_0,y_0)-f(x_0-h,y_0)}h,\qquad
\lim_{h\to0}\frac{f(x_0-2h,y_0)-f(x_0,y_0)}h.
\]
再证明对任意实数 $a$，$\lim_{h\to0}\frac{f(x_0+ah,y_0)-f(x_0,y_0)}h=a f_x(x_0,y_0)$，并单独说明 $a=0$ 的情形。
\end{exercise}
\begin{exercise}{117}{不对称差商与条件辨析}
设 $f_x(x_0,y_0)$ 存在，求 $\lim_{h\to0}\frac{f(x_0+2h,y_0)-f(x_0-3h,y_0)}h$。反过来，仅由这个极限存在，能否推出 $f_x(x_0,y_0)$ 存在？给出证明或反例。
\end{exercise}
\begin{exercise}{118}{绝对值乘积的充分条件}
设 $f(x,y)=|x-y|g(x,y)$，$g$ 在原点邻域有定义。逐项判断下列条件是否足以保证 $f$ 在原点的两个偏导数存在，证明充分者并为不充分者给出反例：$g(0,0)=0$；$\lim g=0$ 且 $g(0,0)=0$；$\lim g$ 存在；$g$ 在原点连续且取值为零。再设 $\lim g=L$，求偏导数存在的充要条件。
\end{exercise}
\begin{exercise}{119}{轴向奇异与连续性}
设 $f(0,0)=0$，非原点处 $f(x,y)=\frac{\sqrt{|x|}}{x^2+y^2}\sin(x^2+y^2)$。讨论原点处两个偏导数是否存在，并证明函数在原点连续，说明连续性与偏导数存在性在本题中的区别。
\end{exercise}
\begin{exercise}{120}{偏导数与局部变化率}
设 $f(x,y)=x^3+3xy+4y-1$。求 $(1,2)$ 处的两个偏导数，并求使 $f(1+h,2+kh)-f(1,2)=o(h)$ 的常数 $k$。
\end{exercise}
\begin{exercise}{121}{特殊点求导与局部定义}
设 $f(x,y)=x+2y+(y-1)\arcsin\frac{x}{1+xy}$。说明原点坐标平移后的点 $(0,1)$ 有一个邻域包含于自然定义域内部，求 $f_x(0,1)$、$f_y(0,1)$ 及 $f_{xy}(0,1)$。
\end{exercise}
\begin{exercise}{122}{立方根的偏导数定义}
设 $f(x,y)=\sqrt[3]{x^5-y^3}$，立方根取实值。用定义求 $f_x(0,0)$、$f_y(0,0)$，并讨论在曲线 $x^5=y^3$ 上的非原点处，两个偏导数是否存在，说明直接套用链式公式的适用范围。
\end{exercise}
\begin{exercise}{123}{初等偏导数与特殊点}
分别求下列函数的两个偏导数，并完成相应判断：
\begin{enumerate}
\item $z=x+2y-\sqrt{x^2+y^2}$；另用定义讨论原点两个偏导数。
\item $z=e^x\sin(x^2+y)$；判断是否存在两个偏导数同时为零的点。
\end{enumerate}
\end{exercise}
\begin{exercise}{124}{幂函数与括号层次}
完成下列求导，并写明所用的实值定义区域：
\begin{enumerate}
\item 设 $u=(\frac xy)^z$，其中 $x/y>0$，求三个一阶偏导数。
\item 在 $x>y$ 上分别设 $U=\ln(x^2+y^2+z^2)+\arctan((x-y)^z)$、$V=\ln(x^2+y^2+z^2)+(\arctan(x-y))^z$。求 $U,V$ 关于 $x,y,z$ 的偏导数，并比较两者关于 $z$ 的求导过程。
\end{enumerate}
\end{exercise}
\begin{exercise}{125}{变底变指数的求导}
在 $xy^2+y>0$ 的区域上，求 $z=(xy^2+y)^x$ 的两个偏导数，并求 $(1,1)$ 处的全微分。写明对数求导中每个变量的作用。
\end{exercise}
\begin{exercise}{126}{齐次复合函数的求导}
在 $x,y>0$ 时设 $z=(\frac yx)^{x/y}$。求 $z_x(1,2)$、$z_y(1,2)$，并证明整个定义区域上恒有 $xz_x+yz_y=0$，解释该恒等式与同时缩放 $x,y$ 的关系。
\end{exercise}
\begin{exercise}{127}{分段函数与偏导连续性}
设 $f(0,0)=0$，非原点处 $f(x,y)=(x^2+y^2)\ln(x^2+y^2)$。求全平面上的两个偏导数，证明它们在原点连续，并判断 $f_{xx}(0,0)$ 是否存在。
\end{exercise}
\begin{exercise}{128}{由偏导数恢复函数}
设 $x>0$，$f_x(x,y)=x+\frac{y^2}x$，$f(1,y)=\sin y$。求 $f$ 及 $f_y$，并说明若定义域改为 $x\ne0$，仅给上述初值是否能唯一确定负半平面上的函数。
\end{exercise}
\begin{exercise}{129}{截线斜率与切向量}
求曲线 $z=\sqrt{1+x^2+y^2},\ x=1$ 在 $(1,1,\sqrt3)$ 处的切线关于 $y$ 轴的斜率，并给出切线参数方程与一个单位切向量。
\end{exercise}
\begin{exercise}{130}{截线倾角与切线方程}
求曲线 $z=\frac{x^2+y^2}4,\ y=4$ 在 $(2,4,5)$ 处沿 $x$ 增大方向的切线倾角，以及该点的切线方程。说明倾角应在哪个平面内理解。
\end{exercise}
\begin{exercise}{131}{坐标单调性与比较}
设 $f\in C^1(\mathbb R^2)$ 且处处 $f_x>0$、$f_y<0$。在 $x_1<x_2,y_1<y_2$；$x_1<x_2,y_1>y_2$；$x_1>x_2,y_1<y_2$；$x_1>x_2,y_1>y_2$ 四组条件中，哪些必推出 $f(x_1,y_1)<f(x_2,y_2)$？用一元中值定理证明，并说明其余条件为何不能保证该结论。
\end{exercise}
\begin{exercise}{132}{连续性与轴向不可导}
证明 $f(x,y)=\sqrt{x^2+y^2}$ 在原点连续但两个偏导数均不存在。进一步证明 $|f(P)-f(Q)|\le |P-Q|$，说明很强的连续性估计仍不能保证偏导数存在。
\end{exercise}
\begin{exercise}{133}{偏导存在与不连续反例}
设 $f(0,0)=0$，非原点处 $f(x,y)=\frac{2xy}{x^2+y^2}$。求原点两个偏导数并证明函数不连续；再求沿 $y=kx$ 的极限，说明两条坐标轴上的信息为何不足。
\end{exercise}
\subsubsection{高阶偏导数}
\begin{exercise}{134}{二阶偏导与调和性}
设 $z=\arctan\frac{x+y}{x-y}$，$x\ne y$。求 $z_{xx},z_{xy}$，并验证 $z_{xx}+z_{yy}=0$。说明混合偏导数可以交换的局部光滑性依据。
\end{exercise}
\begin{exercise}{135}{高阶求导与积分上限}
完成下列计算，并分别说明交换求导顺序的条件：
\begin{enumerate}
\item $f(x,y)=x\sin\frac yx$，求 $f_{xy}(2,\frac\pi2)$ 和 $f_{yx}(2,\frac\pi2)$。
\item $g(x,y)=\int_0^{xy}\frac{\sin t}{1+t^2}\,\dd t$，求 $g_{xx}(0,2)$，并判断是否可以先把 $x=0$ 代入后再连续求两次 $x$ 偏导。
\end{enumerate}
\end{exercise}
\begin{exercise}{136}{三阶偏导与对称性}
设 $z=\ln(e^x+e^y)$。求 $z_{xxy}$，化简后确定其零点集，并比较交换 $x,y$ 后该三阶偏导数的符号变化。
\end{exercise}
\begin{exercise}{137}{径向函数的微分恒等式}
在 $(x,y,z)\ne(0,0,0)$ 上设 $u=\sqrt{x^2+y^2+z^2}$。证明 $u_x^2+u_y^2+u_z^2=1$ 和 $u_{xx}+u_{yy}+u_{zz}=\frac2u$；再直接求导证明 $u^{-1}$ 的三个纯二阶偏导数之和为零。
\end{exercise}
\begin{exercise}{138}{混合偏导不交换的边界}
设 $f(0,0)=0$，非原点处 $f(x,y)=\frac{x^3y}{x^2+y^2}$。分别从定义求 $f_{xy}(0,0)$、$f_{yx}(0,0)$（约定 $f_{xy}=\partial_y(\partial_x f)$），证明两者不等，并讨论这两个混合偏导数在原点的连续性，指出交换定理缺少的条件。
\end{exercise}
\begin{exercise}{139}{轴向信息与混合偏导}
设 $F(0,0)=0$，非原点处 $F(x,y)=\frac{xy(x^2-y^2)}{x^2+y^2}$。用定义求原点的两个一阶偏导数和两个混合偏导数，讨论混合偏导交换的结论是否成立，并证明 $F$ 在原点连续。
\end{exercise}
\begin{exercise}{140}{参数族的一阶与二阶正则性}
对 $a>0$，设 $H_a(0,0)=0$，非原点处 $H_a(x,y)=(x^2+y^2)^a$。分别求使原点处两个一阶偏导数存在、两个纯二阶偏导数存在的 $a$ 的范围，并讨论两个混合偏导数在原点的值；需要时区分端点情形。
\end{exercise}
\subsection{全微分}
\begin{exercise}{141}{全增量与线性近似误差}
设 $f(x,y)=x^2y+xy^2$。求在 $(2,1)$ 处当 $\Delta x=0.1$、$\Delta y=-0.2$ 时的全增量和全微分，并计算两者之差。再对一般增量 $(h,k)$ 写出精确余项，证明其为 $o(\sqrt{h^2+k^2})$。
\end{exercise}
\begin{exercise}{142}{全微分与复合积分}
求下列函数的全微分，并完成相应说明：
\begin{enumerate}
\item $z=(x^2+y^2)e^{\arctan(y/x)}$，$x\ne0$；求在 $(1,0)$ 处的线性近似。
\item $z=\sin x+\sin(xy)+\int_0^{2x+y}e^{-t^2}\,\dd t$；说明变上限积分贡献的两个系数之间有何比例关系。
\end{enumerate}
\end{exercise}
\begin{exercise}{143}{全微分公式的有效区域}
设 $u=\arcsin\frac z{x+y}$，$x+y>0$。求全微分，并写出公式具有有限系数的开区域；说明 $z=\pm(x+y)$ 上不能直接沿用该公式的原因。
\end{exercise}
\begin{exercise}{144}{隐式微分与非退化条件}
设 $z=z(x,y)$ 由 $z-x-y+xe^z=0$ 确定。求 $\dd z$ 并写明非退化条件，再求经过 $(x,y,z)=(0,0,0)$ 的局部分支在原点的全微分。
\end{exercise}
\begin{exercise}{145}{含模长因子的可微性}
设 $f(x,y)=\frac{2x+3y}{1+xy\sqrt{x^2+y^2}}$。从定义证明它在原点可微并求全微分，给出余项关于 $r=\sqrt{x^2+y^2}$ 的一个阶数估计。
\end{exercise}
\begin{exercise}{146}{二阶渐近与一阶微分}
设 $f$ 在原点连续且 $\lim_{(x,y)\to(0,0)}\frac{f(x,y)-x-2y-1}{x^2+y^2}=1$。求 $f(0,0)$ 和 $\dd f|_{(0,0)}$，并写出含二次主项的渐近展开。若去掉原点连续的条件，原点处可微是否仍必然成立？说明理由。
\end{exercise}
\begin{exercise}{147}{全微分的相容性}
设 $u\in C^2(\mathbb R^2)$，$\dd u=3x^2y\,\dd x+f(x,y)\,\dd y$。求 $f_x$ 并说明能否唯一确定 $f$。若另给 $u(0,y)=y^2$，求 $u$ 与 $f$。
\end{exercise}
\begin{exercise}{148}{由全微分恢复函数}
已知 $\dd f=ye^y\,\dd x+[x(1+y)e^y+2y]\,\dd y$，$f(0,0)=2$。求 $f(x,y)$，并通过重新计算两个偏导数验证全微分的两个系数均得到满足。
\end{exercise}
\begin{exercise}{149}{可微定义中的余项尺度}
设 $\Delta f=f(x_0+h,y_0+k)-f(x_0,y_0)$，$r=\sqrt{h^2+k^2}$。判定下列条件中哪些足以保证原点平移后的点 $(x_0,y_0)$ 处可微；对不充分者给出反例：
\begin{enumerate}
\item $\Delta f\to0$；
\item 两个一阶偏导数在该点存在且连续；
\item 两个偏导数存在，且 $\Delta f-f_xh-f_yk\to0$；
\item 两个偏导数存在，且 $\frac{\Delta f-f_xh-f_yk}r\to0$。
\end{enumerate}
比较第三、四项条件中余项尺度的差别。
\end{exercise}
\begin{exercise}{150}{非线性路径与可微必要条件}
设 $f(0,0)=0$，非原点处 $f(x,y)=\frac{3x^2y}{x^4+y^2}$。求原点两个偏导数，并沿适当抛物线判断连续性，从而判定原点可微性。
\end{exercise}
\begin{exercise}{151}{连续但不可微的根式函数}
设 $f(x,y)=\sqrt{x^2+y^4}$。讨论原点处的连续性、两个偏导数和可微性；再说明若仅沿 $y$ 轴观察，为何会遗漏决定性信息。
\end{exercise}
\begin{exercise}{152}{振荡乘积的可微性}
设 $f(0,0)=0$，非原点处 $f(x,y)=xy\sin\frac1{\sqrt{x^2+y^2}}$。用可微定义证明它在原点可微并求全微分；再讨论两个偏导数在原点是否连续，给出明确路径或估计。
\end{exercise}
\begin{exercise}{153}{四类正则性质的逻辑关系}
对某点处的四种性质“连续”“两个偏导数连续”“可微”“两个偏导数存在”，说明以下四条推理链是否正确，并为每条错误链指出失效的箭头，给出相应反例：偏导数连续 $\Rightarrow$ 可微 $\Rightarrow$ 连续；可微 $\Rightarrow$ 偏导数连续 $\Rightarrow$ 连续；可微 $\Rightarrow$ 偏导数存在 $\Rightarrow$ 连续；可微 $\Rightarrow$ 连续 $\Rightarrow$ 偏导数存在。
\end{exercise}
\begin{exercise}{154}{连续与一阶余项检验}
设 $f(0,0)=0$，非原点处 $f(x,y)=\frac{x^2y}{x^2+y^2}$。证明原点连续且两个偏导数存在，再用可微定义检验是否存在全微分；写出使归一化余项不趋于零的具体路径。
\end{exercise}
\begin{exercise}{155}{四种性质的完整辨析}
设 $f(0,0)=0$，非原点处 $f(x,y)=\frac{xy}{\sqrt{x^2+y^2}}$。分别讨论原点的连续性、两个偏导数存在性、偏导数连续性及可微性。计算非原点偏导数，并用它们在坐标轴上的取值核验连续性判断。
\end{exercise}
\begin{exercise}{156}{可微但偏导数不连续}
设 $f(0,0)=0$，非原点处 $f(x,y)=(x^2+y^2)\sin\frac1{x^2+y^2}$。证明它在原点可微并求全微分，证明两个偏导数在原点均不连续；进一步判断它们在原点附近是否有界。
\end{exercise}
\begin{exercise}{157}{余项估计与充分条件的临界阶}
设 $f(0,0)=a$，且原点附近满足 $|f(x,y)-a-bx-cy|\le C(x^2+y^2)^{(1+\alpha)/2}$，其中 $C>0$、$\alpha>0$。证明 $f$ 在原点可微并求其全微分。若将 $\alpha>0$ 放宽为 $\alpha=0$，构造满足估计但不可微的函数。
\end{exercise}
\begin{exercise}{158}{乘积幂函数的可微阈值}
设 $a>0$，$F_a(x,y)=|xy|^a$。讨论原点处的连续性、偏导数及可微性，并求全微分存在的全部参数。再求使 $F_a$ 在原点的某个开邻域内具有连续一阶偏导数的参数范围，说明这与仅在原点可微的条件有何区别。
\end{exercise}
\subsection{多元复合函数的求导法则}
\begin{exercise}{159}{一元复合函数的依赖关系}
设 $z=u^2+v^2+\sin(uv)$，$u=e^x$、$v=x^2$。用链式法则求 $\frac{\dd z}{\dd x}$ 及其在 $x=0$ 处的值，再将 $u,v$ 直接代入求导，核验两种方法一致。
\end{exercise}
\begin{exercise}{160}{全导数与偏导数的区别}
设 $f(x,y)=\arcsin(xy)$，$z(x)=f(x,\ln x)$。在实值可导的开区间内求 $z'(x)$，并分别计算 $z'(1)$ 和 $f_x(1,0)$，说明中间变量的变化怎样进入全导数。
\end{exercise}
\begin{exercise}{161}{退化内映射与一阶导数}
设 $z=f(x\varphi(x),x^2+\varphi(x)^2)$，其中 $f$ 可微，$\varphi$ 可导。求 $z'(x)$，明确 $f$ 的偏导数取值点。若 $\varphi(0)=0$，证明 $z'(0)=0$，无需给出 $f$ 的具体表达式。
\end{exercise}
\begin{exercise}{162}{二阶链式法则}
设 $f\in C^2(\mathbb R^2)$，$z=f(x^2,\sin2x)$。求 $z''(x)$。若 $f_1(0,0)=a$、$f_2(0,0)=b$、$f_{22}(0,0)=c$，求 $z''(0)$，说明哪些数据实际参与结果。
\end{exercise}
\begin{exercise}{163}{内函数驻点与二阶信息}
设 $f\in C^2$ 且 $\dd f|_{(1,1)}=3\,\dd u+4\,\dd v$，$z=f(\cos x,1+x^2)$。求 $z'(0)$、$z''(0)$，并解释为什么不必知道 $f$ 的二阶偏导数的具体值。
\end{exercise}
\begin{exercise}{164}{角度比值与缩放关系}
设 $z=e^u\sin v$，$u=x-y$、$v=\frac yx$，$x\ne0$。求两个一阶偏导数，并证明 $xz_x+yz_y=(x-y)z$。
\end{exercise}
\begin{exercise}{165}{径向复合与原点延拓}
设 $u=\sqrt{x^2+y^2}$，$z=\ln(\arctan u)$。求原点外的两个偏导数及 $xz_x+yz_y$，并判断 $z$ 能否在原点作有限的连续延拓。
\end{exercise}
\begin{exercise}{166}{显含变量与中间变量}
设 $z=e^{x^2+y^2+u^2}$，$u=x^2\sin y$。求 $z_x,z_y$，写出显含 $x,y$ 与通过 $u$ 间接依赖的全部项。再求限制在 $y=0$ 上的一元函数的导数，核验与 $z_x(x,0)$ 一致。
\end{exercise}
\begin{exercise}{167}{链式法则与雅可比退化}
设 $z=f(x\arcsin y,\frac yx)$，$f\in C^1$，$x\ne0$、$|y|<1$。求 $z_x,z_y$，并求 $(1,0)$ 处两个偏导数；说明此点的 $z_x$ 是否与 $f$ 的选择有关。
\end{exercise}
\begin{exercise}{168}{二阶偏导与变量重新表示}
设 $u=x+y$、$v=x^2-y^2$，$z=f(u,v)$，其中 $f\in C^2$。求 $z_{xx},z_{yy}$，并将 $z_{xx}-z_{yy}$ 的结果仅用 $u,v$ 及 $f$ 的偏导数表示。
\end{exercise}
\begin{exercise}{169}{混合偏导的交叉项}
设 $z=f(e^{xy},x^2+y^2)$，$f\in C^2$。求 $z_{xy}$，并化简其在原点的值，说明为什么该点的结果只涉及 $f$ 的一个一阶偏导数。
\end{exercise}
\begin{exercise}{170}{由截面条件恢复外函数}
设 $z=e^{x-y}+f(x-2y)$，且对所有实数 $x$ 有 $z(x,0)=x^3$。先求 $f$，再求 $z_x,z_y,z_{xy}$，并验证恢复出的函数满足给定截面条件。
\end{exercise}
\begin{exercise}{171}{特殊点全微分与零变化方向}
设 $z=yf(x^2-y^2)$，$f$ 可导且 $f(0)=1$、$f'(0)=\frac12$。求 $\dd z|_{(1,1)}$，并求在该点使一阶变化为零的全部非零位移方向。
\end{exercise}
\begin{exercise}{172}{变上下限积分的混合求导}
设 $z=f(2x+y)+f(x-2y)+\int_{x-y}^{x^2+y^2}g(t)\,\dd t$，其中 $f$ 二阶可导，$g$ 一阶可导。求 $z_{xy}$ 及 $z_{xy}(0,0)$，说明上下限两项的符号如何确定。
\end{exercise}
\begin{exercise}{173}{齐次分量与微分恒等式}
设 $z=xy+xf(\frac yx)$，$x\ne0$，$f$ 二阶可导。证明 $xz_x+yz_y=xy+z$ 及 $x^2z_{xx}=y^2z_{yy}$。将 $z$ 分为一次和二次齐次分量，说明第一个恒等式为何不等于通常的一次齐次公式。
\end{exercise}
\begin{exercise}{174}{偏微分约束化为常微分方程}
在 $x,y>0$ 上设 $z=xf(u)$、$u=\frac xy$，$f\in C^1(0,+\infty)$，满足 $xz_x-yz_y=0$ 及 $f(1)=1$。求 $f(u)$，核验所得函数满足约束，并讨论它在 $u=0$ 能否作有限连续延拓。
\end{exercise}
\begin{exercise}{175}{多层依赖的二阶求导}
设 $z=f(2x-y)+g(x,xy)$，$f$ 二阶可导，$g\in C^2$。求 $z_{xx},z_{xy}$；再令 $g(u,v)=uv$，用直接展开的方法核验与 $g$ 有关的各项。
\end{exercise}
\begin{exercise}{176}{三元外函数与二阶链式法则}
设 $u=f(x,y,xe^y)$，$f\in C^2(\mathbb R^3)$。求 $u_{xx},u_{yy}$，并化简 $u_{yy}(0,y)$；说明外函数关于第三变量的一阶偏导项在何种二次求导中出现。
\end{exercise}
\begin{exercise}{177}{嵌套中间变量}
设 $u=f(x,y,z)$，$y=\varphi(x,t)$、$t=\psi(x,z)$，各函数可微，以 $x,z$ 为独立变量。求 $u_x,u_z$。再令 $\varphi(x,t)=x+t$、$\psi(x,z)=xz$，直接消去中间变量检验所得公式。
\end{exercise}
\begin{exercise}{178}{线性变换与二阶算子}
设 $f(x,y)=g(2x+y,3x-y)$，$g\in C^2(\mathbb R^2)$，且 $f_{xx}+f_{xy}-6f_{yy}=1$。求 $g_{uv}$。若另有 $g(u,0)=g(0,v)=0$，求 $g$ 并验证。
\end{exercise}
\begin{exercise}{179}{欧拉公式及二阶推广}
设 $f\in C^2(\mathbb R^2)$，对任意实数 $t$ 有 $f(tx,ty)=t^2f(x,y)$。证明 $xf_x+yf_y=2f$ 及 $x^2f_{xx}+2xyf_{xy}+y^2f_{yy}=2f$；写明两次对参数 $t$ 求导后在哪一点取值，并单独核验原点。
\end{exercise}
\begin{exercise}{180}{全微分形式不变性}
设 $z=\arctan(u+v)$，$u=2x-y^2$、$v=x^2y$。用全微分形式不变性求 $\dd z$，再求 $(1,0)$ 处的全微分，并通过直接代入后的偏导数核验。
\end{exercise}
\begin{exercise}{181}{复合全微分与切向位移}
设 $z=f(2x+3y,e^{xy})$，$f$ 可微。求全微分，并求在 $(0,0)$ 处使 $\dd z=0$ 对每个可微外函数 $f$ 都成立的全部非零位移方向。
\end{exercise}
\begin{exercise}{182}{链式法则所需的外函数条件}
设 $F(0,0)=0$，非原点处 $F(u,v)=\frac{uv}{\sqrt{u^2+v^2}}$，并令 $u(t)=v(t)=t$。证明 $F$ 在原点连续且两个偏导数存在，但 $F(u(t),v(t))$ 在 $t=0$ 不可导。说明只知道外函数偏导数存在时，链式法则的哪项充分条件没有满足。
\end{exercise}
\begin{exercise}{183}{极坐标复合与梯度模}
在 $x>0$ 上设 $r=\sqrt{x^2+y^2}$、$\theta=\arctan\frac yx$，$z=F(r,\theta)$，$F\in C^1$。证明 $z_x^2+z_y^2=F_r^2+\frac1{r^2}F_\theta^2$。若 $F=r^a\cos(b\theta)$、$a,b>0$，求使 $z_x^2+z_y^2$ 在整个右半平面上为常数的全部参数，并求该常数。
\end{exercise}
\subsection{隐函数的求导法则}
\begin{exercise}{184}{隐函数导数与退化点}
由 $e^x+\cos y=xy^2$ 确定的可微分支 $y=y(x)$，求 $y'$ 并写明分母非零条件。核验 $(0,\pi)$ 满足方程，并说明在该点能否直接使用这一求导公式及隐函数存在定理。
\end{exercise}
\begin{exercise}{185}{隐函数初值与切线}
设 $2^{xy}=x+y$ 在 $x=0$ 附近确定 $y=y(x)$。先求 $y(0)$ 并验证隐函数定理的条件，再求 $y'(0)$ 及曲线在该点的切线。
\end{exercise}
\begin{exercise}{186}{二阶隐式求导与常值分支}
设 $y-xe^y+x=0$。求经过 $x=0$ 的局部分支的 $y(0),y'(0),y''(0)$，并用隐函数定理与直接代入说明该局部分支是否恒为常数。
\end{exercise}
\begin{exercise}{187}{含未知函数的隐式求导}
设 $f\in C^1$，$y>0$，$xf(x+y)=e^x\ln y$ 在非退化点确定 $y=y(x)$。求 $y'$ 并写明非退化条件；再求经过 $x=0$ 的分支的初值和 $y'(0)$，结果可含 $f(1)$。
\end{exercise}
\begin{exercise}{188}{隐函数与外函数的二阶复合}
设 $z=f(\ln y-\sin x)$，$f$ 二阶可导且 $f'(0)=1$，$y=y(x)$ 由 $y-xe^{y-1}=1$ 在 $x=0$ 附近确定。求 $z'(0),z''(0)$，并解释为什么无需知道 $f''(0)$ 的值。
\end{exercise}
\begin{exercise}{189}{不同因变量的隐函数条件}
对方程 $xy-z\ln y+e^{xz}=1$，在 $(0,1,1)$ 处计算三个一阶偏导数，判定隐函数定理分别能否保证 $x=x(y,z)$、$y=y(x,z)$、$z=z(x,y)$。再令 $x=0,y=1$，说明本题为何不能唯一确定最后一种局部分支。
\end{exercise}
\begin{exercise}{190}{隐式复合中的结构性抵消}
设 $z+xe^{z-x-y}=x+y$ 在 $1+xe^{z-x-y}\ne0$ 的点附近确定 $z=z(x,y)$。求 $z_x,z_y$，并通过代换 $w=z-x-y$ 解释其中一个偏导数恒为常数的原因。
\end{exercise}
\begin{exercise}{191}{隐函数偏导及恒等式}
设 $xy=xf(z)+y\varphi(z)$，其中 $f,\varphi\in C^1$，$D=xf'(z)+y\varphi'(z)\ne0$。求 $z_x,z_y$，并证明 $xz_x+yz_y=\frac{xy}D$。
\end{exercise}
\begin{exercise}{192}{共同平移下的不变性}
在 $\tan z$ 有定义的局部分支上设 $(x-y)^3+z+\tan z=0$。求 $z_x+z_y$，说明隐函数定理的非退化条件自动成立，并证明在同一局部分支内同时平移 $x,y$ 相同的微小量时，$z$ 保持不变。
\end{exercise}
\begin{exercise}{193}{隐式全微分的非退化条件}
设 $e^z-xy^2+\sin(y+z)=0$ 在某点确定 $z=z(x,y)$。求 $\dd z$，明确给出保证该公式成立的偏导数条件，并写出 $z_x,z_y$；说明为何必须先检查该点满足原方程。
\end{exercise}
\begin{exercise}{194}{隐式微分与图形切平面}
设 $xyz+\sqrt{x^2+y^2+z^2}=\sqrt2$。证明在 $(1,0,1)$ 附近可确定 $z=z(x,y)$，求 $\dd z|_{(1,0)}$，并据此写出该分支图形在此处的切平面。
\end{exercise}
\begin{exercise}{195}{隐函数的二阶偏导}
设 $z^3-3xyz=1$。在 $z^2-xy\ne0$ 的局部分支上求 $z_{xx}$，可保留 $z,z_x$ 的中间表达式，但最终须消去 $z_x$。同时说明方程上的点是否自动满足这一非退化条件。
\end{exercise}
\begin{exercise}{196}{指定分支的二阶偏导和}
设 $z+e^x-y\ln(1+z^2)=0$。求 $(x,y)=(0,0)$ 对应的 $z$ 值，验证局部隐函数条件，并计算 $z_{xx}(0,0)+z_{yy}(0,0)$。
\end{exercise}
\begin{exercise}{197}{直接偏导与复合偏导}
设 $f(x,y,z)=yz^2e^x$，$z=z(x,y)$ 由 $x+y+z+xyz=0$ 确定，并令 $F(x,y)=f(x,y,z(x,y))$。求 $F_x(0,1)$，再分别求 $F_y(0,1)$ 与 $f_y(0,1,-1)$，说明两个记号的含义。
\end{exercise}
\begin{exercise}{198}{隐式中间变量的全微分}
设 $u=f(x,y,z)$，$f\in C^1$，其中 $z=z(x,y)$ 由 $xe^x-ye^y=ze^z$ 确定。求 $\dd u$ 并写明适用的非退化条件，明确区分外函数偏导数与复合函数的全微分系数。
\end{exercise}
\begin{exercise}{199}{低阶数据确定特殊二阶导数}
设 $f\in C^1$、$f(1,1)=0$，$z=z(x,y)$ 由 $(x+1)z-y^2=xf(x+z,x+y)$ 在 $(0,1,1)$ 附近确定。求 $\dd z|_{(0,1)}$、$z_{yy}(0,1)$，并说明求第二个量时为何不必假设 $f$ 有二阶偏导数。
\end{exercise}
\subsubsection{方程组确定的隐函数}
\begin{exercise}{200}{方程组的导数与投影退化}
设 $x+2y+3z=0$、$x^2+y^2+z^2=2$。在能以 $x$ 为参数的局部分支上求 $y',z'$；求曲线上使关于 $(y,z)$ 的雅可比行列式为零的全部点，并说明导数公式的适用范围。
\end{exercise}
\begin{exercise}{201}{隐式方程组与四个偏导}
设 $u^2+xv-y=0$、$yu+v^2-x=0$ 在某点确定 $u=u(x,y),v=v(x,y)$。写出雅可比非退化条件，求 $u_x,v_x,u_y,v_y$，并说明为什么不能不加检查就在 $(x,y,u,v)=(0,0,0,0)$ 使用所求分式。
\end{exercise}
\begin{exercise}{202}{方程组在特殊初值的简化}
设 $f\in C^1$、$f(-1)=1$，$y=y(x),z=z(x)$ 由 $z=xf(x-y)$、$x+y+ze^z=1$ 在 $x=0$ 附近确定。先求初值并核验雅可比条件，再求 $y'(0),z'(0)$，说明为何不需要 $f'(-1)$ 的具体值。
\end{exercise}
\begin{exercise}{203}{两个隐式中间函数的全导数}
设 $u=f(x,y,z)$，$f\in C^1$，$y=y(x)$、$z=z(x)$ 分别由 $e^{xy}+2y=0$ 与 $e^z-xz=2$ 确定。求各非退化分支上的 $\frac{\dd u}{\dd x}$；再求 $x=0$ 处的值，结果以 $f$ 在对应三元点的偏导数表示。
\end{exercise}
\begin{exercise}{204}{非退化条件的充分性}
分别考察方程 $(z-x-y)^3=0$ 与 $z^2-x^2=0$ 在原点附近能否唯一确定连续可微分支 $z=z(x,y)$。计算各方程对 $z$ 的偏导数，并用这两个例子说明隐函数定理中的偏导非零条件与分支存在、唯一之间的逻辑关系。
\end{exercise}
\begin{exercise}{205}{隐函数分支的临界参数}
设 $z+\lambda\sin z=x+y$，考虑经过 $(0,0,0)$ 的分支。求使隐函数定理适用的参数，并在这些参数下求原点的一阶、二阶偏导数。对 $\lambda=-1$，证明仍有唯一连续局部分支，再判断其原点可微性。
\end{exercise}
\subsection{多元函数微分法在几何上的应用}
\begin{exercise}{206}{参数曲线的切线与法平面}
设空间曲线为 $x=t+1$、$y=t^2-1$、$z=t^3$。求 $(2,0,1)$ 处的切线及法平面，写出对应参数与沿参数增大方向的单位切向量，并验证切线方向垂直于法平面。
\end{exercise}
\begin{exercise}{207}{参数法与交面法求切线}
求曲线 $y=x,\ z=x^2-2y$ 在原点的切线与法平面。分别以 $x$ 为参数和用两个曲面法向量的向量积求切向量，说明两种结果一致。
\end{exercise}
\begin{exercise}{208}{平行切线的数量与定位}
对曲线 $y=-x^2,\ z=x^3$，求所有与平面 $x+2y+z=4$ 平行的切线，给出切点及切线方程，并判断这些切线是否可能位于该平面内。
\end{exercise}
\begin{exercise}{209}{交曲线的正则性与法平面}
求曲线 $x^2+y^2+z^2=6,\ x+y+z=0$ 在 $(1,-2,1)$ 处的切线和法平面。先核验该点在曲线上，再证明两个曲面的法向量在此处不平行。
\end{exercise}
\begin{exercise}{210}{隐式曲面的切平面与法线}
求曲面 $x^2+\cos(xy)+yz+x=0$ 在 $(0,1,-1)$ 处的切平面，并求法线参数方程。指出该点附近可由隐函数定理选择哪些坐标作为因变量。
\end{exercise}
\begin{exercise}{211}{椭球切平面与轴截距}
求椭球面 $x^2+2y^2+3z^2=9$ 在 $(2,1,1)$ 处的切平面和法线，进一步求切平面的三个有向轴截距及其到原点的距离。
\end{exercise}
\begin{exercise}{212}{切平面存在与图形表示}
求曲面 $e^z-z+xy=3$ 在 $(2,1,0)$ 处的切平面和法线。说明虽然该点处 $F_z=0$，曲面仍有切平面的原因；再判断隐函数定理能否在此直接保证 $z=z(x,y)$。
\end{exercise}
\begin{exercise}{213}{非标准显式曲面的法向量}
求曲面 $y=x^2+3z^2-15$ 在 $(1,-2,2)$ 处的切平面和法线，并求法线与 $xOz$ 平面的夹角。说明法向量应由哪个隐式函数的梯度得到。
\end{exercise}
\begin{exercise}{214}{平行切平面的反求问题}
求抛物面 $z=4-x^2-y^2$ 上使切平面平行于 $2x+2y+z-1=0$ 的全部点、对应切平面及两个平面的距离，并说明切点是否唯一。
\end{exercise}
\begin{exercise}{215}{切平面包含整条直线}
设 $L:x+y+b=0,\ x+ay-z-3=0$ 位于曲面 $z=x^2+y^2$ 在 $(1,-2,5)$ 处的切平面内。求 $a,b$，并将直线参数式代入切平面，核验包含关系对全部参数成立。
\end{exercise}
\begin{exercise}{216}{零梯度与切平面判定}
平面 $z=0$ 也可写为 $z^3=0$。比较这两种定义函数在原点的梯度，并说明为什么梯度为零不能直接推出曲面没有切平面。再考察圆锥面 $z^2=x^2+y^2$ 在原点的各条母线方向，说明为何不存在能同时承载其一阶方向的单一平面。
\end{exercise}
\begin{exercise}{217}{椭球的指定法向切平面}
设 $\lambda>0$。在椭球面 $x^2+y^2+\lambda z^2=1$ 上求全部点，使其切平面平行于 $x+2y+z=0$；求对应切平面及它们之间的距离，并求切平面到原点的距离等于 $1$ 时的参数 $\lambda$。
\end{exercise}
\subsection{方向导数和梯度}
\begin{exercise}{218}{方向单位化与最大变化率}
求 $f(x,y)=3x^2y-y^2$ 在 $A(1,1)$ 处沿 $A$ 指向 $B(2,1+\sqrt3)$ 的方向导数，写出单位方向向量；再求该点方向导数的最大值并判断给定方向是否达到它。
\end{exercise}
\begin{exercise}{219}{沿曲线正向的方向导数}
设曲线 $y=x^2-5$ 的正向为 $x$ 增大的方向。求 $f(x,y)=x^2-xy+y^2$ 在 $(2,-1)$ 沿曲线切线正向的方向导数，并求反向的结果，说明归一化因子的来源。
\end{exercise}
\begin{exercise}{220}{方向角的相容性}
设 $f(x,y,z)=\ln(x+\sqrt{y^2+z^2})$。先验证 $\alpha=60^\circ,\beta=45^\circ,\gamma=60^\circ$ 可以作为一个空间方向的方向角，再求 $(1,0,1)$ 沿该方向的方向导数及该点方向导数的最大值。
\end{exercise}
\begin{exercise}{221}{单侧方向导数与可微性}
设 $f(x,y)=\sqrt{x^2+y^2}$。采用 $D_{\vec e}f(0)=\lim_{t\to0^+}\frac{f(t\vec e)-f(0)}t$ 的定义，对每个单位方向求方向导数，并讨论连续性、偏导数和可微性。若改用 $t\to0$ 的双侧极限，结论如何？
\end{exercise}
\begin{exercise}{222}{曲面外法向与方向导数}
设 $\vec n$ 为曲面 $2x^2+3y^2+z^2=6$ 在 $P(1,1,1)$ 处的外法向量，$u=\frac{\sqrt{6x^2+8y^2}}z$。求 $P$ 处的外单位法向量及沿该方向的方向导数，再给出沿内法向的结果。
\end{exercise}
\begin{exercise}{223}{梯度与等值线切向}
设 $f(x,y)=\arctan\frac xy$。求 $\nabla f(1,1)$ 及该点使方向导数为零的全部单位方向，说明它们与经过该点的等值线有什么关系。
\end{exercise}
\begin{exercise}{224}{最快增减与零变化方向}
设 $f(x,y,z)=x^3-xy^2-z$，$P(1,1,0)$。求该点增长最快的单位方向及变化率，再求下降最快的单位方向和全部零变化率单位方向，并说明后者形成的几何集合。
\end{exercise}
\begin{exercise}{225}{给定方向与梯度方向比较}
求 $f(x,y,z)=x^2+y^2+xz$ 在 $P(1,0,1)$ 处的梯度，以及沿 $\vec l=(3,-1,6)$ 的方向导数。判断该方向是否为最快增长方向，并给出判断依据。
\end{exercise}
\begin{exercise}{226}{沿另一函数梯度方向求导}
设 $u=\sqrt{x^2+y^2+z^2}$，$g=\frac12x^2+yz$。求在 $(1,1,-1)$ 处沿 $g$ 的梯度方向的 $u$ 的方向导数，判断该方向上 $u$ 是增加还是减少，并与 $u$ 自身的最快增加方向比较。
\end{exercise}
\begin{exercise}{227}{隐函数梯度与方向极值}
设 $x+y^2+z^3=3$ 在 $(x,y)=(1,1)$ 附近确定实值分支 $z=z(x,y)$。验证隐函数条件，求该点方向导数的最大值、最小值及分别取得它们的单位方向。
\end{exercise}
\begin{exercise}{228}{全方向导数存在仍需统一余项}
设 $r=\sqrt{x^2+y^2}$，$F(0,0)=0$，非原点处 $F(x,y)=r\frac{x^3y}{x^6+y^2}$。证明 $F$ 在原点连续且沿任意单位方向的方向导数均为零，再沿 $y=x^3$ 检验可微性。说明所有方向导数均由同一个线性式给出，为何仍不足以保证可微。
\end{exercise}
\begin{exercise}{229}{受切向约束的最速变化}
设温度场 $T(x,y,z)=x^2+2y^2+3z^2$，只能从 $P(1,1,1)$ 沿球面 $x^2+y^2+z^2=3$ 的切向单位方向瞬时移动。求允许方向中温度方向导数的最大值、最小值及取得它们的方向，并求全部零变化率允许方向。用梯度在切平面上的正交投影解释结果。
\end{exercise}
\subsection{多元函数的极值及其应用}
\begin{exercise}{230}{局部极小的必要条件}
设 $f$ 在开区域 $D$ 内可微，内点 $P$ 为局部极小点。证明 $f_x(P)=f_y(P)=0$，并判断能否推出 $f(Q)>f(P)$ 对所有 $Q\in D\setminus\{P\}$ 成立；不能时给出具体反例。
\end{exercise}
\begin{exercise}{231}{极值与坐标截线的导数}
设 $f$ 在内点 $(x_0,y_0)$ 的两个偏导数存在，且在该点取得局部极大值。仅利用两条坐标截线的一元极值必要条件证明两个偏导数为零，并举例说明逆命题不成立。
\end{exercise}
\begin{exercise}{232}{二重渐近确定极值}
设 $f$ 在原点连续，且 $\lim_{(x,y)\to(0,0)}\frac{f(x,y)}{\sin(x^2+y^2)}=1$。求 $f(0,0)$ 和两个偏导数，判断原点是否为严格极值点，并用统一估计证明其类型。
\end{exercise}
\begin{exercise}{233}{四阶主项的符号控制}
设 $f$ 在原点邻域连续，且 $\lim_{(x,y)\to(0,0)}\frac{f(x,y)-x^2y^2}{(x^2+y^2)^2}=-1$。求 $f(0,0)$，判断原点极值类型，并证明存在 $\delta>0$，使 $0<x^2+y^2<\delta^2$ 时 $f(x,y)\le-\frac12(x^2+y^2)^2$。
\end{exercise}
\begin{exercise}{234}{三次多项式的驻点分类}
求 $z=x^2y-xy^2-xy$ 的全部驻点，逐一判定极值或鞍点并求极值；再判断函数在全平面是否有全局最大值或最小值。
\end{exercise}
\begin{exercise}{235}{指数因子的驻点与全局性质}
求 $f(x,y)=e^{-x}(x-y^3+3y)$ 的全部驻点及局部极值，分别给出二阶判别依据；再利用固定 $x$ 改变 $y$ 判断其全局有界性。
\end{exercise}
\begin{exercise}{236}{局部极值与全局界}
求 $f(x,y)=xe^{-(x^2+y^2)/2}$ 的全部驻点与极值，并证明所得极大值、极小值分别为全局最大值、最小值，写出全部取等点。
\end{exercise}
\begin{exercise}{237}{截面数据恢复与正定性}
设 $f(x,y)=2y^2+g(x-y)$，且对所有实数 $x$ 有 $f(x,0)=x^2$。先恢复 $g$ 与 $f$，再求极值，并分别用二阶判别及平方和表示证明其全局性质。
\end{exercise}
\begin{exercise}{238}{非孤立驻点与退化判别}
求 $f(x,y)=\frac12x^2+xy+\frac12y^2-2x-2y+5$ 的全部驻点和极值。说明为何二阶判别式为零仍可由配方得到结论，并区分严格极小与非严格极小。
\end{exercise}
\begin{exercise}{239}{隐函数的存在与极值}
设正值函数 $z=z(x,y)$ 由 $(x^2+y^2)z+\ln z+2(x+y+1)=0$ 确定。先证明每个 $(x,y)$ 对应唯一正根且隐函数处处可微，再求全部极值，判断所求极大值是否为全局最大值。
\end{exercise}
\begin{exercise}{240}{无盖水箱的全局最省材料}
设 $V>0$，用铁板制作容积为 $V$ 的无盖长方体水箱。求最省材料的长、宽、高及所需铁板面积，并证明所得驻点给出全局最小值，不能只列驻点方程。
\end{exercise}
\begin{exercise}{241}{闭区域内点与分段边界}
求 $f(x,y)=x^2+2y^2-x^2y^2$ 在 $D=\{(x,y):x^2+y^2\le4,\ y\ge0\}$ 上的最大值、最小值及全部取值点。要求分别列出内部驻点、直径边界与半圆边界的候选点并比较。
\end{exercise}
\begin{exercise}{242}{条件极值与偏导关系}
设 $f,\varphi\in C^1$，$P$ 为约束 $\varphi=0$ 下的局部极值点，且 $\varphi_y(P)\ne0$。证明 $f_x\varphi_y-f_y\varphi_x=0$。据此判断 $f_x(P)\ne0$ 是否必有 $f_y(P)\ne0$；再分别举例说明 $f_x(P)=0$ 时 $f_y(P)$ 可为零也可非零。
\end{exercise}
\begin{exercise}{243}{线性约束与距离最小化}
求 $f(x,y)=x^2+y^2-3$ 在 $x-y+1=0$ 下的全部条件极值，判断是否有最大值，并用原点到直线的距离解释最小值及其唯一性。
\end{exercise}
\begin{exercise}{244}{球内接长方体与全局最大体积}
设 $R>0$。求内接于半径为 $R$ 的球面的长方体的最大体积及长、宽、高。分别用拉格朗日乘数法与基本不等式证明结论，并说明退化边界不能取得最大值。
\end{exercise}
\begin{exercise}{245}{由梯度恢复函数后的区域最值}
设 $f(0,0)=1$，$f_x=6x$、$f_y=2y-1$。恢复 $f$，再求它在 $x^2+\frac{y^2}2\le1$ 上的最大值、最小值及全部取值点；需比较区域内部与椭圆边界。
\end{exercise}
\begin{exercise}{246}{相同二阶数据与不同极值类型}
构造三个二元多项式，使它们在原点的函数值、一阶偏导数和全部二阶偏导数完全相同，但原点分别为严格极小点、严格极大点和鞍点。逐一验证，并说明二阶判别式为零时不能仅凭这些数据决定极值类型。
\end{exercise}
\begin{exercise}{247}{四次型的参数分类与约束最值}
设 $F_a(x,y)=x^4+y^4+2ax^2y^2$，$a\in\mathbb R$。按参数完整分类原点的严格极小、非严格极小或鞍点性质；再求 $F_a$ 在单位圆 $x^2+y^2=1$ 上的最大值、最小值及全部取值点，单独讨论临界参数。
\end{exercise}
\subsection{二元函数的泰勒公式}
\begin{exercise}{248}{多项式泰勒展开与零余项}
设 $f(x,y)=2x^2-xy-y^2+3x+5y-1$。求它在 $(1,1)$ 处任意 $n\ge2$ 阶带拉格朗日余项的泰勒公式，并解释余项为什么恒为零；用展开式恢复原多项式作核验。
\end{exercise}
\begin{exercise}{249}{非原点的二阶泰勒公式}
求 $f(x,y)=x^3y^2$ 在 $(1,-1)$ 处带皮亚诺余项的二阶泰勒公式。令 $h=x-1,k=y+1$，明确写出常数、线性、混合二次项，并说明余项相对于 $h^2+k^2$ 的阶。
\end{exercise}
\begin{exercise}{250}{不同截断阶数的余项}
对 $f(x,y)=x^4+xy+(1+y)^2$，分别写出带皮亚诺余项的一阶、二阶麦克劳林公式，并求两种截断的精确余项，说明同一个 $x^4$ 项在两种余项条件下是否都可忽略。
\end{exercise}
\begin{exercise}{251}{三阶展开与二重极限}
求 $f(x,y)=\ln(1+x+y)$ 的三阶麦克劳林公式，余项取 $o((x^2+y^2)^{3/2})$；据此求 $\lim_{(x,y)\to(0,0)}\frac{\ln(1+x+y)+xy-x-y}{x^2+y^2}$，说明混合二次项如何抵消。
\end{exercise}
\begin{exercise}{252}{多元中值估计与符号区域}
设 $f$ 在全平面可微，$f_x>1$、$f_y<-1$、$f(0,0)=0$。分别讨论四个开象限中是否必有 $f(x,y)>x-y$，对必然成立或必然反向的情形作证明，对不能确定的情形给出两个线性函数作为例子。
\end{exercise}
\begin{exercise}{253}{二阶泰勒多项式的唯一性}
设两个次数不超过二的多项式 $P,Q$ 满足 $P(x,y)-Q(x,y)=o(x^2+y^2)$。证明 $P\equiv Q$。说明只沿两条坐标轴核验余项为什么不足以确定混合二次项，并给出明确的多项式反例。
\end{exercise}
\begin{exercise}{254}{退化二次项之后的极值判定}
设 $F_a(x,y)=e^{x+y}-1-x-y-a xy$。求原点处的三阶泰勒公式，并对全部实数 $a$ 判定原点是严格极小点、非严格极小点还是鞍点。二阶主项退化的参数须单独讨论，必要时沿直线检验更高阶项或使用指数函数不等式。
\end{exercise}
\subsubsection{综合提高题}
\begin{exercise}{255}{偏导符号与四点排序}
设 $f\in C^1(\mathbb R^2)$，$f_x>0$、$f_y<0$，记 $a=f(x,y)$、$b=f(x,y+1)$、$c=f(x+1,y)$、$d=f(x+1,y+1)$。确定四个数中必然的最大者与最小者并证明；再说明 $a,d$ 的次序能否确定，给出不同次序及相等情形的例子。
\end{exercise}
\begin{exercise}{256}{增长阶与可微性的四个命题}
设 $f$ 在原点连续。逐项判定下列命题，证明正确者并为错误者给出反例：
\begin{enumerate}
\item $\lim\frac{f(x,y)}{|x|+|y|}$ 为有限数，则 $f$ 在原点可微；
\item $\lim\frac{f(x,y)}{x^2+y^2}$ 为有限数，则 $f$ 在原点可微；
\item $f$ 在原点可微，则第一项中的极限存在；
\item $f$ 在原点可微，则第二项中的极限存在。
\end{enumerate}
所有极限均指 $(x,y)\to(0,0)$。再讨论第一项的极限值为零时可得到什么加强结论。
\end{exercise}
\begin{exercise}{257}{可微性与偏导数的定义区域}
设 $f(x,y)=(x^2+y^2)\sin\frac1{xy}$（$xy\ne0$），在 $xy=0$ 时取零。判断原点可微性并求全微分；在 $xy\ne0$ 处求 $f_x$，判断它能否连续延拓到原点；另用定义讨论坐标轴非原点处相应偏导数是否存在。
\end{exercise}
\begin{exercise}{258}{由二阶渐近判断极值与微分}
设 $\lim_{(x,y)\to(0,0)}\frac{f(x,y)-f(0,0)-x^3-2y^3}{\cos\sqrt{x^2+y^2}-1}=-2$。求 $f-f(0,0)$ 的二次主项，讨论原点连续性、两个偏导数、可微性及极值类型，并说明能否为极大点。
\end{exercise}
\begin{exercise}{259}{仅有两个偏导数时的结论边界}
已知 $f_x(0,0)=1$、$f_y(0,0)=2$。分别判断能否推出原点连续、存在全微分 $\dd x+2\,\dd y$、沿任意单位方向的方向导数为 $\cos\alpha+2\cos\beta$、沿 $x$ 轴负向的单侧方向导数为 $-1$。证明必然成立的结论，并构造反例说明其余结论不一定成立。
\end{exercise}
\begin{exercise}{260}{复合数据能否确定外函数偏导}
设 $f\in C^1((0,+\infty)^3)$，对所有 $x>y>0$ 有
\[
f(x-y,x+y,xy)=(x^2-y^2)^{xy}+(xy)^{2x}.
\]
判断这些条件能否唯一确定 $f_1(1,3,2)$，若不能，构造两个满足条件但该偏导数不同的函数。另令 $h(x,y)=f(x-y,x+y,xy)$，求可由已知数据确定的 $h_x(2,1)$，说明外函数偏导与复合函数导数的信息差别。
\end{exercise}
\begin{exercise}{261}{二阶偏导与截面初值}
设 $f_{yy}(x,y)=2$，$f(x,0)=1$、$f_y(x,0)=x$。求 $f$ 并证明在该数据下唯一，再判定原点是否为极值点。
\end{exercise}
\begin{exercise}{262}{沿曲线求导恢复二阶数据}
设 $f\in C^2$，$f_{xx}=f_{yy}$，且 $f(x,3x)=x$、$f_x(x,3x)=x^3$。求 $f_{xx}(x,3x)$ 和 $f_{xy}(x,3x)$。要求分别写出对两条给定关系求导得到的方程，再联立消元。
\end{exercise}
\begin{exercise}{263}{三元外函数的全微分}
设 $z=f(x,x+y,\frac xy)$，$f$ 可微，$y\ne0$。求 $\dd z$，并化简 $(x,y)=(0,1)$ 处的表达式，明确各偏导数在哪个三元点取值。
\end{exercise}
\begin{exercise}{264}{嵌套乘积的混合偏导}
设 $u=f(x,xy,xyz)$，其中 $f\in C^2$，$x,y,z$ 为独立变量。求 $u_{xy}$，化简 $u_{xy}(0,y,z)$，并说明不能先把 $x=0$ 代入使函数成为常数再求该混合偏导的原因。
\end{exercise}
\begin{exercise}{265}{零全微分的充要条件}
已知 $\lim_{(x,y)\to(0,0)}\frac{f(x,y)-f(0,0)}{\sqrt{x^2+y^2}}=0$。从定义求 $\dd f|_{(0,0)}$。再证明这一极限条件等价于原点可微且两个一阶偏导数均为零。
\end{exercise}
\begin{exercise}{266}{累次极限与联合趋近}
设 $x,y>0$，$f(x,y)=\frac y{1+xy}-\frac{1-y\sin(\pi x/y)}{\arctan x}$。先求 $g(x)=\lim_{y\to+\infty}f(x,y)$，再求 $\lim_{x\to0^+}g(x)$。判断联合趋近 $x\to0^+,y\to+\infty$ 时是否有相同极限，可比较 $y=x^{-2}$ 与 $y=x^{-3}$ 两条路径。
\end{exercise}
\begin{exercise}{267}{有理化与连续参数}
设 $a>0$，$f(0,0)=0$，非原点处 $f(x,y)=\frac{\sqrt{a+x^2y^2}-1}{x^2+y^2}$。讨论原点连续性的全部参数，并在连续的参数下进一步求两个偏导数及全微分；证明须覆盖任意趋近路径。
\end{exercise}
\begin{exercise}{268}{复合约束中的隐函数导数}
设 $u=f(x,y,z)$，$\varphi(x^2,y,z)=0$、$y=\sin x$，$f,\varphi\in C^1$ 且 $\varphi_3\ne0$。求 $\frac{\dd u}{\dd x}$。若 $\varphi(0,0,z_0)=0$，再写出经过该点的局部分支在 $x=0$ 的全导数，明确两组偏导数的取值点。
\end{exercise}
\begin{exercise}{269}{径向调和方程与初值}
设 $f\in C^2(0,+\infty)$，$z=f(\sqrt{x^2+y^2})$ 在原点外满足 $z_{xx}+z_{yy}=0$。推导 $f''(r)+\frac1r f'(r)=0$；若 $f(1)=0$、$f'(1)=1$，求 $f$，并讨论其原点处的有限连续延拓是否可能。
\end{exercise}
\begin{exercise}{270}{二阶方程的可逆变量变换}
设 $z\in C^2$ 满足 $z_{xx}+z_{xy}-z_{yy}=0$。求 $a,b$，使可逆变换 $u=x+ay,v=x+by$ 将方程化为 $z_{uv}=0$。计算变换后的非零比例系数，并说明为何不能选择 $a=b$。
\end{exercise}
\begin{exercise}{271}{隐式齐次关系与退化海森矩阵}
设 $z\ne0$，由 $\frac xz=\varphi(\frac yz)$ 确定 $C^2$ 局部分支，且 $x-y\varphi'(y/z)\ne0$。证明 $xz_x+yz_y=z$ 及 $z_{xx}z_{yy}=z_{xy}^2$；在第二个证明中说明为什么向量 $(x,y)$ 必非零，以及它与二阶偏导矩阵的关系。
\end{exercise}
\begin{exercise}{272}{空间交曲线的二阶隐式求导}
设 $a>0$，$z=x^2+y^2$、$x^2+2y^2+3z^2=4a^2$ 在关于 $(y,z)$ 的雅可比行列式非零的点附近确定 $y=y(x),z=z(x)$。求 $y'',z''$，最终表达式仅含 $x,y,z$；写出雅可比非退化条件，并利用 $z\ge0$ 将条件简化。
\end{exercise}
\begin{exercise}{273}{极值点数据与二阶复合导数}
设 $z=f(xy,yg(x))$，$f\in C^2$，$(1,1)$ 是 $f$ 的极值点，$g$ 由 $(x-1)^2-\int_x^{g(x)}e^{-t^2}\,\dd t=0$ 确定。求 $g(1),g'(1)$，并将 $z_{xy}(1,1)$ 用 $f$ 在 $(1,1)$ 的二阶偏导数表示。若该点为严格极小点，能否断言所得值严格为正？说明理由。
\end{exercise}
\begin{exercise}{274}{曲线数据确定切平面}
设 $f$ 可微，$f(t^2,t+1)=t^2(t-1)$，且 $f_1(1,2)=1$。求 $f(1,2)$、$f_2(1,2)$，再求曲面 $z=f(x,y)$ 在对应点的切平面与法线，说明曲线数据提供了哪个导数约束。
\end{exercise}
\begin{exercise}{275}{包含给定直线的椭球切平面}
在椭球面 $x^2+2y^2+3z^2=21$ 上求全部点 $M$，使过 $M$ 的切平面包含直线 $\frac{x-6}2=\frac{y-3}1=\frac{2z-1}{-2}$。求对应切平面，并把直线的基点条件与方向条件分别写出。
\end{exercise}
\begin{exercise}{276}{全微分相容性与最速增长}
设 $(3x^2+axy+2y^2)\,\dd x+(bxy+x^2)\,\dd y$ 是某个 $C^2$ 函数 $z$ 的全微分，且 $z(0,0)=1$。求 $a,b,z$，再求 $(1,0)$ 处方向导数的最大值及达到它的单位方向；用两个偏导数核验恢复结果。
\end{exercise}
\begin{exercise}{277}{二值隐函数的分支极值}
对方程 $x^2-6xy+10y^2-2yz-z^2+18=0$，先显式解出两个实值分支，证明它们在全平面正则。分别求两分支的全部极值点、极值，并判断相应极值是否为各分支的全局最值。
\end{exercise}
\begin{exercise}{278}{倒数约束与全局极值}
设 $x,y,z,a>0$，$\frac1x+\frac1y+\frac1z=\frac1a$。求 $f=xyz$ 的全部条件极值及取值点，证明其全局性质，并判断最大值是否存在，可将变量换为三个倒数。
\end{exercise}
\begin{exercise}{279}{双约束的可行参数区间}
在 $z=x^2+y^2$、$x+y+z=4$ 的约束下求 $u=x^2+y^2+z^2$ 的最大值、最小值及全部取值点。先推导可行的 $z$ 区间，再将目标化为 $z$ 的一元函数，验证区间端点确可实现。
\end{exercise}
\begin{exercise}{280}{约束曲面上的方向变化率极值}
在曲面 $2x^2+2y^2+z^2=1$ 上，求使 $f(x,y,z)=x^2+y^2+z^2$ 沿固定方向 $\vec l=(1,-1,0)$ 的方向导数最大的全部点及最大值；同时求最小值与取值点，并说明方向导数中的单位化。
\end{exercise}
\begin{exercise}{281}{椭球切面截体的最小体积}
设 $a,b,c>0$。在椭球面 $\frac{x^2}{a^2}+\frac{y^2}{b^2}+\frac{z^2}{c^2}=1$ 的 $x,y,z>0$ 部分求点 $P$，使其切平面与三个坐标面围成的四面体体积最小。求点、切平面和最小体积，并解释趋向坐标边界不能产生更小体积。
\end{exercise}
\begin{exercise}{282}{切平面投影域上的最值}
设 $f(x,y)=x^3+y^3-(x+y)^2+3$。求曲面 $z=f(x,y)$ 在 $(1,1,1)$ 处的切平面 $T$，并写出 $T$ 与三个坐标面围成的有界区域在 $xOy$ 面上的投影 $D$。再求 $f$ 在 $D$ 上的最大值、最小值及全部取值点，分别检查三条边与内部。
\end{exercise}
\begin{exercise}{283}{由边际成本恢复总成本与约束优化}
某企业生产甲、乙两种产品的固定成本为 $10000$ 万元，产量为 $x,y$ 件，边际成本分别为 $C_x=20+\frac12x$、$C_y=6+y$，单位为万元/件。求总成本 $C(x,y)$。在 $x+y=50$、$x,y\ge0$ 时求最省成本的产量及最小成本；先按实数产量求解，再判断非负整数产量限制是否改变最优结果。
\end{exercise}

\clearpage
\section{第10章\quad 重积分}
\subsection{二重积分的概念与性质}

\begin{exercise}{284}{积分方程与参数讨论}
设 $D=\{(x,y):0\leq y\leq x\leq1\}$，$\lambda\in\mathbb R$。讨论连续函数方程
\[
 f(x,y)=xy+\lambda\iint_D f(u,v)\,\dd u\dd v
\]
的解是否存在、是否唯一；存在时求出 $f(x,y)$，并给出 $\lambda=1$ 时的结果。
\end{exercise}

\begin{exercise}{285}{二重和的积分极限}
将下列极限分别表示为二重积分并计算，说明两者积分区域的差别：
\[
 \lim_{n\to\infty}\sum_{i=1}^n\sum_{j=1}^n
 \frac{n}{(n+i)(n^2+j^2)},\qquad
 \lim_{n\to\infty}\sum_{i=1}^n\sum_{j=1}^i
 \frac{n}{(n+i)(n^2+j^2)}.
\]
第二个极限只须化为一个定积分，无须求出初等形式。
\end{exercise}

\begin{exercise}{286}{黎曼和与误差估计}
计算
\[
 L=\lim_{n\to\infty}\frac{\pi}{2n^4}
 \sum_{i=1}^n\sum_{j=1}^n i^2\sin\frac{j\pi}{2n}.
\]
利用 $x^2$ 与 $\sin y$ 在相应区间上的单调性，说明该和式与 $L$ 之差为 $O(1/n)$。
\end{exercise}

\begin{exercise}{287}{积分比较与严格性}
设 $D$ 为顶点 $(1,0),(1,1),(0,1)$ 的三角形闭区域。比较
$\iint_D\ln(x+y)\,\dd\sigma$ 与 $\iint_D[\ln(x+y)]^2\,\dd\sigma$ 的大小，并证明所得不等式为严格不等式，说明边界上取等号是否影响结论。
\end{exercise}

\begin{exercise}{288}{幂函数比较与严格性}
设 $D=\{(x,y):(x-1)^2+(y-1)^2\leq2\}$，$u=(x+y)/4$。确定 $u$ 在 $D$ 上的取值范围，严格比较
\[
 I_1=\iint_Du\,\dd\sigma,\quad I_2=\iint_D\sqrt u\,\dd\sigma,\quad I_3=\iint_D\sqrt[3]u\,\dd\sigma.
\]
进一步证明 $J(p)=\iint_Du^p\,\dd\sigma$ 在 $p>0$ 时严格递减。
\end{exercise}

\begin{exercise}{289}{积分估值与分离变量}
设 $D=[0,\pi]\times[0,\pi]$。先用被积函数的最大、最小值估计
$I=\iint_D\sin^2x\sin^2y\,\dd\sigma$，再利用变量分离求其准确值；求 $\sin^2x\sin^2y$ 在 $D$ 上的平均值，并解释平均值与最大值不同的原因。
\end{exercise}

\begin{exercise}{290}{最优函数界与积分估值}
设 $D=\{(x,y):x^2+y^2\leq4\}$。求使 $m\leq x^2+4y^2+9\leq M$ 在 $D$ 上成立的最大 $m$ 和最小 $M$，据此估计二重积分；再计算其准确值，检验估计结果。
\end{exercise}

\begin{exercise}{291}{分片常值函数的积分}
设 $0<a<b<3$，$D=[0,3]\times[0,3]$，
\[
 f_{a,b}(x,y)=\begin{cases}1,&0\leq y\leq a,\\2,&a<y\leq b,\\3,&b<y\leq3.\end{cases}
\]
求 $\iint_D f_{a,b}\,\dd\sigma$，给出 $a=1,b=2$ 的结果；若只改变 $y=a$ 与 $y=b$ 两条线段上的有界函数值，积分是否改变？说明理由。
\end{exercise}

\begin{exercise}{292}{平均值与中值点}
设 $D=\{(x,y):0\leq y\leq x\leq1\}$。求 $f(x,y)=xy$ 在 $D$ 上的平均值，并求出 $D$ 内所有使 $f(\xi,\eta)$ 等于该平均值的点，说明其中是否存在区域内部的点。
\end{exercise}

\begin{exercise}{293}{缩小圆域上的平均值极限}
设 $D_r=\{(x,y):x^2+y^2\leq r^2\}$，$r>0$，
$A(r)=\frac1{\pi r^2}\iint_{D_r}e^{x^2-y^2}\cos(x+y)\,\dd x\dd y$。
求 $\lim_{r\to0^+}A(r)$；进一步利用二阶展开与圆域对称性，求 $\lim_{r\to0^+}[A(r)-1]/r^2$，并说明余项积分后仍可忽略的依据。
\end{exercise}

\begin{exercise}{294}{局部连续性与积分极限}
设 $f$ 在原点的某邻域内连续，$f(0,0)=1$；$g$ 在 $0$ 的某邻域内连续可微，且 $g(0)=0,g'(0)=3$。令 $D_r=\{x^2+y^2\leq r^2\}$。
求 $\lim_{r\to0^+}\frac{\iint_{D_r}f(x,y)\,\dd x\dd y}{g(r^2)}$，并用连续性的定义证明分子为 $\pi r^2+o(r^2)$，不能仅引用形式上的无穷小替换。
\end{exercise}

\begin{exercise}{295}{对称性使用的存在性条件}
有人利用奇偶性声称
$\iint_{[-1,1]\times[-2,2]}(x^2\tan y+\sin x^3)\,\dd\sigma=0$。
指出被积函数在区域内的奇点，并判断该式能否作为通常的有界函数二重积分成立。
再设 $0<h<\pi/2$，计算同一被积函数在 $[-1,1]\times[-h,h]$ 上的二重积分，写明对称性成立所需的条件。
\end{exercise}

\begin{exercise}{296}{变量互换与参数积分}
设 $D=\{(x,y):x^2+y^2\leq4,x\leq0,y\leq0\}$，$f$ 在 $[-2,0]$ 上连续且严格为正，$a,b\in\mathbb R$。
不用求 $f$ 的原函数，计算
$\iint_D\frac{af(x)+bf(y)}{f(x)+f(y)}\,\dd\sigma$。
给出积分关于 $a,b$ 的依赖关系，并说明 $a=b$ 时的结果与一般公式是否一致。
\end{exercise}

\begin{exercise}{297}{三角形区域上的对称分解}
设 $D=\{(x,y):-1\leq x\leq y\leq1\}$，$D_1=D\cap\{x\geq0,y\geq0\}$。
分别计算 $\iint_Dx^3y\,\dd\sigma$ 与 $\iint_D\cos x\sin y\,\dd\sigma$，从而求两者之和；证明后一个积分等于 $2\iint_{D_1}\cos x\sin y\,\dd\sigma$，指出所用的区域变换。
\end{exercise}

\begin{exercise}{298}{对称性与圆域积分}
设 $D=\{(x,y):x^2+y^2\leq1\}$。不先作二次积分，利用对称性求
\[
 \iint_D(x^2+2y)\,\dd x\dd y,
\]
并求 $\iint_D(ax^2+by^2+cxy+dy)\,\dd x\dd y$，其中 $a,b,c,d$ 为实常数。
\end{exercise}

\begin{exercise}{299}{保号性与中值定理的条件}
设 $D$ 为单位闭圆盘，$f$ 在 $D$ 上连续且非负。证明 $\iint_Df\,\dd\sigma=0$ 当且仅当 $f\equiv0$；若将连续条件去掉，仅要求有界可积，考察“原点取 $1$、其余点取 $0$”的函数，说明结论是否保持。
再令 $E$ 为两个互不相交的单位闭圆盘之并，在两圆盘上分别令 $g=0$ 与 $g=2$。判断能否找到 $P\in E$ 使 $g(P)$ 等于 $g$ 的积分平均值，并指出中值结论还需要的区域条件。
\end{exercise}

\begin{exercise}{300}{双矩积分方程的全部参数分支}
设 $D=[0,1]^2$，$\lambda\in\mathbb R$，连续函数 $f$ 满足
\[
 f(x,y)=x+y+\lambda\iint_D(xu+yv)f(u,v)\,\dd u\dd v.
\]
记 $A=\iint_Duf(u,v)\,\dd u\dd v$，$B=\iint_Dvf(u,v)\,\dd u\dd v$。建立关于 $A,B$ 的线性方程，完整讨论无解、唯一解及无穷多解的参数情形，写出存在时的全部 $f$，并代回验证。
\end{exercise}

\subsection{二重积分的计算方法}
\subsubsection{积分区域与交换积分次序}

\begin{exercise}{301}{交换积分次序与计算}
设 $f$ 在 $D=\{(x,y):0\leq x\leq4,\ x\leq y\leq2\sqrt{x}\}$ 上连续。
把 $\int_0^4\dd x\int_x^{2\sqrt{x}}f(x,y)\,\dd y$ 改写成先对 $x$ 积分的累次积分，并用所得表达式计算 $f(x,y)=y$ 时的积分。
\end{exercise}

\begin{exercise}{302}{半圆域与积分次序}
对连续函数 $f$，将
\[
 \int_0^1\dd y\int_{-\sqrt{1-y^2}}^{\sqrt{1-y^2}}f(x,y)\,\dd x
\]
交换积分次序。再取 $f(x,y)=x^2+y^2$，写出极坐标表达式并计算其值。
\end{exercise}

\begin{exercise}{303}{分块区域的合并}
设 $f$ 在下式涉及的闭区域上连续，将
\[
 \int_0^1\dd x\int_{-\sqrt{x}}^{\sqrt{x}}f(x,y)\,\dd y
 +\int_1^4\dd x\int_{x-2}^{\sqrt{x}}f(x,y)\,\dd y
\]
改写为一个先对 $x$ 积分的累次积分，并由此求该积分区域的面积。
\end{exercise}

\begin{exercise}{304}{分段边界与换序}
设 $f$ 连续，把
\[
 \int_0^1\dd y\int_0^{2y}f(x,y)\,\dd x
 +\int_1^3\dd y\int_0^{3-y}f(x,y)\,\dd x
\]
交换积分次序；取 $f(x,y)=y$，计算积分并求该区域的形心纵坐标。
\end{exercise}

\begin{exercise}{305}{三角函数边界的换序}
设 $f$ 在 $D=\{(x,y):0\leq x\leq\pi,-\sin(x/2)\leq y\leq\sin x\}$ 上连续。
将 $\iint_Df\,\dd x\dd y$ 改为先对 $x$ 积分的表达式，写明必须分段的纵坐标，并求 $D$ 的面积。
\end{exercise}

\begin{exercise}{306}{抛物线区域的分段换序}
设 $f$ 在 $D=\{(x,y):1\leq y\leq2,y\leq x\leq y^2\}$ 上连续。
将 $\int_1^2\dd y\int_y^{y^2}f(x,y)\,\dd x$ 交换积分次序，写明分段点；再取 $f(x,y)=x-y$ 计算，并从区域位置验证结果为正。
\end{exercise}

\begin{exercise}{307}{分块区域的重新组合}
设 $f$ 连续，记
\[
 I=\int_1^2\dd x\int_x^{x^2}f(x,y)\,\dd y
 +\int_2^8\dd x\int_x^8f(x,y)\,\dd y.
\]
分别用不等式描述两项对应的区域，将 $I$ 全部改为先对 $x$ 积分，并尽量合并具有相同积分限的部分；计算 $f\equiv1$ 时的值。
\end{exercise}

\begin{exercise}{308}{变限二重积分求导}
设 $f$ 在 $[1,+\infty)$ 上连续，$F(t)=\int_1^t\dd y\int_y^t f(x)\,\dd x$，$t\geq1$。
先交换积分次序，再求 $F'(t)$ 与 $F'(3)$；若 $f$ 处处为正，判断 $F$ 在 $(1,+\infty)$ 上的单调性并说明理由。
\end{exercise}

\subsubsection{直角坐标计算与分区}

\begin{exercise}{309}{矩形域与参数变化}
设 $a,b>0$，$D_{a,b}=[0,a]\times[0,b]$。求 $I(a,b)=\iint_{D_{a,b}}e^{x-y}\,\dd x\dd y$，给出 $I(1,1)$，并判断固定 $a$ 时 $I(a,b)$ 随 $b$ 的变化及其极限。
\end{exercise}

\begin{exercise}{310}{区域分割与积分份额}
设 $D=[0,1]^2$，以直线 $x+y=1$ 分割为 $D_-=\{(x,y)\in D:x+y\leq1\}$ 与 $D_+=D\setminus D_-$。计算 $x+y$ 在 $D,D_-,D_+$ 上的积分，并判断两个等面积部分的积分是否相等。
\end{exercise}

\begin{exercise}{311}{含参积分与零参数}
设 $\alpha\in\mathbb R$，求 $I(\alpha)=\int_0^1\int_0^1xe^{\alpha xy}\,\dd y\dd x$。分别处理 $\alpha=0$ 与 $\alpha\ne0$，求 $I(1)$，并检验所得表达式在 $\alpha=0$ 处是否连续。
\end{exercise}

\begin{exercise}{312}{三角形域与尺度参数}
设 $a>1$，$D_a$ 为直线 $x=a,y=1,y=x$ 围成的闭区域。求 $\iint_{D_a}xy\,\dd x\dd y$，给出 $a=2$ 的结果，并求 $xy$ 在 $D_a$ 上的平均值及其在 $a\to1^+$ 时的极限。
\end{exercise}

\begin{exercise}{313}{曲边域与加权平均}
设 $D$ 为 $y=\sqrt{x}$ 与 $y=x^2$ 围成的闭区域，面密度为 $\rho(x,y)=\sqrt y$。计算其质量与关于 $y$ 轴的一阶矩 $\iint_Dx\sqrt y\,\dd x\dd y$，进而求质心横坐标。须写出两曲线交点及积分限。
\end{exercise}

\begin{exercise}{314}{线性密度与形心比较}
设 $a>0$，$D_a=\{(x,y):x\geq0,y\geq0,x+y\leq a\}$，面密度 $\rho=3x+2y$。求质量 $M(a)$ 与质心横坐标，给出 $M(2)$，并比较质心横坐标与均匀薄片的形心横坐标。
\end{exercise}

\begin{exercise}{315}{三角形积分的参数表达}
设 $a>0$，$D_a=\{(x,y):0\leq y\leq x\leq a\}$。求 $I(a)=\iint_{D_a}x\cos(x+y)\,\dd x\dd y$，计算 $I(\pi)$；利用积分定义判断 $I(a)$ 在足够小的正 $a$ 下的符号，检验公式。
\end{exercise}

\begin{exercise}{316}{双曲线边界与参数}
设 $a>1$，$D_a$ 为 $x=a,xy=1,y=x$ 围成的第一象限闭区域。求 $\iint_{D_a}x^2/y^2\,\dd x\dd y$，给出 $a=2$ 的结果；说明该积分在每个固定 $a>1$ 时不存在奇点。
\end{exercise}

\begin{exercise}{317}{非均匀薄片的一阶矩}
设 $D=\{(x,y):0\leq x\leq1,x\leq y\leq5x\}$，密度为 $x+6y$。计算质量与关于 $y$ 轴的一阶矩，求质心横坐标；判断能否仅用几何形心代替质心，并用计算结果说明。
\end{exercise}

\begin{exercise}{318}{梯形区域与边界参数}
设 $0<h<\pi-1$，$D_h=\{(x,y):0\leq x\leq1,0\leq y\leq h+x\}$。求 $I(h)=\iint_{D_h}(1+x)\sin y\,\dd x\dd y$，给出 $h=1$ 的结果，并从积分区域与被积函数判断 $I(h)$ 的单调性。
\end{exercise}

\begin{exercise}{319}{缩放区域与积分符号}
设 $a>0$，$D_a=\{(x,y):0\leq y\leq a,y/2\leq x\leq y\}$。求
$I(a)=\iint_{D_a}(x^2+y^2-x)\,\dd x\dd y$，计算 $I(2)$，并完整讨论 $I(a)$ 的正负与零点。
\end{exercise}

\begin{exercise}{320}{平行四边形上的变量代换}
设 $D=\{(x,y):1\leq y\leq3,y-1\leq x\leq y\}$。以 $u=y-x,v=y$ 把积分化为矩形域积分，计算 $\iint_D(x^2+y^2)\,\dd x\dd y$ 与 $\iint_D(y^2-x^2)\,\dd x\dd y$，并说明后者为何严格为正。
\end{exercise}

\begin{exercise}{321}{三角形薄片的一阶矩}
设 $D$ 为顶点 $O(0,0),A(1,2),B(2,1)$ 的三角形闭区域。求 $\iint_Dx\,\dd x\dd y$，再利用 $x,y$ 互换的对称性求均匀薄片的形心，并用三角形顶点坐标的平均值检验结果。
\end{exercise}

\begin{exercise}{322}{积分次序与非初等原函数}
设 $D=\{(x,y):0\leq y\leq2,1\leq x\leq y+1\}$。将 $\iint_D\cos(y^2)\,\dd x\dd y$ 写成两种次序的累次积分，选取能直接计算的次序求值，并说明为何无须求 $\cos(y^2)$ 的初等原函数。
\end{exercise}

\begin{exercise}{323}{可去奇点与曲边区域}
令 $f(x)=\sin x/x$（$x\ne0$），$f(0)=1$。设 $D$ 由 $y=x,y=x^2$ 围成。计算 $\iint_Df(x)\,\dd x\dd y$，并证明在整个 $D$ 上被积函数连续；若仅改变原点处的取值，积分是否改变？
\end{exercise}

\begin{exercise}{324}{三角形域与符号估计}
设 $D=\{(x,y):0\leq x\leq y\leq1\}$。求 $I=\iint_D(e^{-y^2}-1)\,\dd x\dd y$；不依赖准确结果，直接利用被积函数证明 $-1/2<I<0$，再检验所求值。
\end{exercise}

\begin{exercise}{325}{密度积分与高斯型因子}
设 $D=\{(x,y):0\leq x\leq y\leq1\}$。求 $\iint_Dx^2e^{-y^2}\,\dd x\dd y$；若薄片密度为 $e^{-y^2}$，再求其关于 $y$ 轴的转动惯量与质量之比。
\end{exercise}

\begin{exercise}{326}{曲线间宽度与精确积分}
设 $0<a<1$，$D_a=\{(x,y):a\leq x\leq1,x^2\leq y\leq x\}$。求 $I(a)=\iint_{D_a}e^{x^2}\,\dd x\dd y$ 的一元积分表达式；在 $a=1/2$ 时将其化为含 $\int_{1/2}^1e^{x^2}\,\dd x$ 的表达式，并证明 $I'(a)<0$。允许保留所指定的非初等定积分。
\end{exercise}

\begin{exercise}{327}{换序与线性分解}
将 $I=\int_0^1\dd x\int_x^1(y+xye^{y^2})\,\dd y$ 交换积分次序并计算。另求 $y$ 项与 $xye^{y^2}$ 项对积分的贡献，比较两项大小。
\end{exercise}

\begin{exercise}{328}{参数上限与积分求导}
设 $a>0$，$I(a)=\int_0^a\dd x\int_x^a\frac{x}{\sqrt{1+y^3}}\,\dd y$。交换次序求 $I(a)$，给出 $I(1)$；由积分表达式与所得公式分别求 $I'(a)$，检验两者一致。
\end{exercise}

\begin{exercise}{329}{角点反常积分与换序}
设 $D=\{(x,y):1<x\leq2,0\leq y\leq\ln x\}$，定义
\[
 I=\lim_{\varepsilon\to0^+}\int_{1+\varepsilon}^2\dd x\int_0^{\ln x}\frac{e^{xy}}{x^x-1}\,\dd y.
\]
证明该极限存在并计算；写出形式上先对 $x$ 积分的表达式，指出角点 $(1,0)$ 使本题不能直接按连续函数的有界区域积分处理的原因。
\end{exercise}

\begin{exercise}{330}{换序与分部积分}
设 $f\in C^3[0,2]$，$f(0)=f'(0)=f''(0)=-1$，$f(2)=-1/2$。计算
\[
 I=\int_0^2\dd x\int_0^x\sqrt{(2-x)(2-y)}f'''(y)\,\dd y.
\]
先将 $I$ 化为仅含端点值 $f(0),f'(0),f''(0),f(2)$ 的表达式；说明为什么未知的 $f'(2),f''(2)$ 不影响结果。
\end{exercise}

\begin{exercise}{331}{平面支集与卷积截面}
设 $f(x)=g(x)=a$（$0\leq x\leq1$），在其余点取零，其中 $a>0$。求
$H(y)=\int_{-\infty}^{+\infty}f(x)g(y-x)\,\dd x$ 的分段表达式，进而求 $\iint_{\mathbb R^2}f(x)g(y-x)\,\dd x\dd y$。须从两个因子的非零条件确定支集。
\end{exercise}

\begin{exercise}{332}{函数支集与圆外区域}
设 $f(x,y)=x^2y$（$1\leq x\leq2,0\leq y\leq x$），其余点取零。求 $f$ 在 $D=\{x^2+y^2\geq2x\}$ 上的积分，并求在补集上的积分；用两者之和检验对整个支集的积分。
\end{exercise}

\begin{exercise}{333}{绝对值分区与最小值}
设 $0\leq c\leq1$，$D=[-1,1]\times[0,1]$。求 $I(c)=\iint_D|y-cx^2|\,\dd x\dd y$，计算 $I(1)$，并求使 $I(c)$ 最小的 $c$。写明去掉绝对值所需的分界曲线。
\end{exercise}

\begin{exercise}{334}{最大值函数与分界双曲线}
设 $0<c\leq4$，$D=[0,2]^2$。求 $I(c)=\iint_D\max\{xy,c\}\,\dd x\dd y$，给出 $I(1)$；从几何分区说明 $c=4$ 时公式应为何值，并检验结果。
\end{exercise}

\begin{exercise}{335}{对角线分割与连续参数}
设 $\alpha\in\mathbb R$，$D=[0,1]^2$。沿 $y=x$ 分区求
$I(\alpha)=\iint_De^{\alpha\max\{x^2,y^2\}}\,\dd x\dd y$，给出 $I(1)$，并用极限检验 $\alpha=0$ 时的值。
\end{exercise}

\begin{exercise}{336}{根式边界与缩放规律}
设 $R>0$，$D_R=\{(x,y):\sqrt{|x|}+\sqrt{|y|}\leq R\}$。求 $D_R$ 的面积及
$\iint_{D_R}(\sqrt{|x|}+\sqrt{|y|})\,\dd x\dd y$，给出 $R=1$ 的结果；说明两者分别按 $R$ 的几次幂缩放，并验证函数平均值小于 $R$。
\end{exercise}

\subsubsection{极坐标与圆形边界}

\begin{exercise}{337}{扇环域与角函数平均}
设 $0<\alpha<\beta<\pi/2$，$D=\{(r\cos\theta,r\sin\theta):1\leq r\leq3,\alpha\leq\theta\leq\beta\}$。
求 $\iint_D\arctan(y/x)\,\dd x\dd y$ 及该函数在 $D$ 上的平均值；给出边界直线为 $x=\sqrt3y$ 与 $y=\sqrt3x$ 时的结果。
\end{exercise}

\begin{exercise}{338}{圆域根式积分}
设 $D=\{(x,y):x^2+y^2\leq1,x\geq0,y\geq0\}$。计算
$I=\iint_D\sqrt{\frac{1-x^2-y^2}{1+x^2+y^2}}\,\dd x\dd y$；用被积函数的范围独立证明 $0<I<\pi/4$，检验准确值。写出极坐标中的面积因子。
\end{exercise}

\begin{exercise}{339}{半球体积与半径确定}
设 $a,V>0$，$D_a=\{x^2+y^2\leq a^2\}$。求
$\iint_{D_a}\sqrt{a^2-x^2-y^2}\,\dd x\dd y$ 关于 $a$ 的表达式；证明令其等于 $V$ 时半径 $a$ 唯一确定，求 $V=\pi$ 时的 $a$，并给出体积解释。
\end{exercise}

\begin{exercise}{340}{四分之一圆域上的平均距离}
设 $a>0$。计算 $\int_0^a\dd x\int_0^{\sqrt{a^2-x^2}}\sqrt{x^2+y^2}\,\dd y$，进而求四分之一圆盘上点到原点距离的平均值，说明所得平均值与 $a$ 的比例为何不随 $a$ 改变。
\end{exercise}

\begin{exercise}{341}{平移圆域与积分零点}
设 $h\in\mathbb R$，$D_h=\{(x,y):x^2+(y-h)^2\leq1\}$。求 $I(h)=\iint_{D_h}(x^2-y)\,\dd x\dd y$，给出 $h=0$ 的结果，并求所有使 $I(h)=0$ 的 $h$。
\end{exercise}

\begin{exercise}{342}{环域与径向对数积分}
设 $R>1$，$D_R=\{(x,y):1\leq x^2+y^2\leq R^2\}$。求 $I(R)=\iint_{D_R}x^2\ln(x^2+y^2)\,\dd x\dd y$，给出 $R=e$ 的结果；利用对称性说明把 $x^2$ 换为 $y^2$ 后积分不变，并求换为 $x^2+y^2$ 后的值。
\end{exercise}

\begin{exercise}{343}{互换对称性与孤立不连续点}
设 $D=\{x^2+y^2\leq1,x\geq0,y\geq0\}$。将
$x\cos(x^2+y^2)/(x+y)$ 在原点定义为零，求其在 $D$ 上的积分；判断此延拓在原点是否连续，说明该点是否影响积分。可将分子 $x$ 与 $y$ 互换后相加简化计算。
\end{exercise}

\begin{exercise}{344}{原点奇性与径向截断}
设 $D=\{(x,y):0\leq x\leq1,0\leq y\leq x\}$。用截断区域 $D\cap\{x^2+y^2\geq\varepsilon^2\}$ 定义并计算
$\iint_D(x^2+y^2)^{-1/2}\,\dd x\dd y$，证明极限有限，并说明极坐标面积因子如何抵消原点奇性。
\end{exercise}

\begin{exercise}{345}{扇形边界下的径向矩}
设 $a>0$，$D_a=\{(x,y):0\leq x\leq a,0\leq y\leq x\}$。计算 $\iint_{D_a}\sqrt{x^2+y^2}\,\dd x\dd y$，求 $a=1$ 的结果；求距离的平均值，并判断其是否小于区域内最大距离。
\end{exercise}

\begin{exercise}{346}{参数扇域与退化积分}
设 $a\geq1$，定义 $I(a)=\int_0^1\dd x\int_x^{ax}\frac1{\sqrt{x^2+y^2}}\,\dd y$，原点处按反常积分极限理解。用适当换元求 $I(a)$，证明其收敛并讨论单调性；求 $a=1$ 的值，解释积分区域退化时结果如何理解。
\end{exercise}

\begin{exercise}{347}{两圆之间的极坐标区域}
设 $D$ 为第一象限内由 $y=0,y=x$ 及圆弧 $x^2+y^2=4x,x^2+y^2=9x$ 围成的闭区域。求其极坐标范围，计算 $\iint_Dy^2/x^2\,\dd x\dd y$；另求 $D$ 的面积，从而求 $(y/x)^2$ 的平均值。
\end{exercise}

\begin{exercise}{348}{平移旋转与半圆积分}
设 $D=\{(x,y):(x-1)^2+(y-1)^2\leq2,y\geq x\}$。作适当的平移与旋转，把 $x-y$ 化为一个坐标，求 $\iint_D(x-y)\,\dd x\dd y$，并求 $x-y$ 在 $D$ 上的平均值。
\end{exercise}

\begin{exercise}{349}{偏心圆盘的矩}
设 $a>0$，$D_a=\{(x,y):x^2+y^2\leq ax\}$。求 $\iint_{D_a}(x+y^2)\,\dd x\dd y$，给出 $a=1$ 的结果；分别用平移圆心后的对称性确定两个分项的缩放次数。
\end{exercise}

\begin{exercise}{350}{圆域积分方程的可解性}
设 $D=\{(x,y):x^2+y^2\leq y,x\geq0\}$，$\lambda\in\mathbb R$。讨论连续函数方程
\[
 f(x,y)=\sqrt{1-x^2-y^2}-\lambda\iint_Df(u,v)\,\dd u\dd v
\]
的存在性与唯一性，存在时求其解；给出 $\lambda=8/\pi$ 的结果。必须算出区域面积及根式项的积分。
\end{exercise}

\begin{exercise}{351}{最小值函数与径向分区}
设 $D=\{x^2+y^2\leq3/16\}$。求两函数 $\sqrt{3/16-x^2-y^2}$ 与 $2(x^2+y^2)$ 在 $D$ 内的相等点集，计算其最小值的二重积分；再利用“最小值加最大值等于两者之和”求最大值的积分。
\end{exercise}

\begin{exercise}{352}{绝对值与挖去圆盘}
设 $D=\{(x,y):x^2+y^2\geq x,|x|\leq1,|y|\leq1\}$。把 $D$ 解释为正方形挖去一个圆盘后的区域，求 $\iint_D|xy|\,\dd x\dd y$；同时求 $\iint_Dxy\,\dd x\dd y$，说明两个积分能否使用相同的奇偶性结论。
\end{exercise}

\subsubsection{一般换元与无界积分}

\begin{exercise}{353}{线性变换与雅可比}
设 $a,b>0$，$D_{a,b}=\{0\leq x+y\leq a,0\leq x-y\leq b\}$。求
$\iint_{D_{a,b}}(x+y)\sin(x-y)\,\dd x\dd y$，给出 $a=b=\pi$ 的结果。写出正、逆变换及雅可比绝对值，并解释为何不能遗漏该因子。
\end{exercise}

\begin{exercise}{354}{椭圆变换与幂参数}
设 $a,b>0$，$p\geq0$，$D=\{x^2/a^2+y^2/b^2\leq1\}$。求
$\iint_D(x^2/a^2+y^2/b^2)^p\,\dd x\dd y$，分别给出 $p=0,1$ 的结果，并从中求原二次函数在椭圆内的平均值。
\end{exercise}

\begin{exercise}{355}{和与比值变量的变换}
设 $D=\{x\geq0,y\geq0,x+y\leq1\}$。在原点给 $e^{x/(x+y)}$ 任意有限值。用 $u=x+y,v=x/(x+y)$ 计算积分，并求以该函数为面密度时的总质量及 $x+y$ 的加权平均值；说明变换在原点失效是否影响结果。
\end{exercise}

\begin{exercise}{356}{幂乘积变换与区域面积}
设 $D$ 为第一象限内由 $xy=4,xy=8,xy^3=5,xy^3=15$ 围成的闭区域。取 $u=xy,v=xy^3$，证明逆变换在相应矩形内唯一，求雅可比绝对值并计算 $D$ 的面积；再求 $\iint_Dxy\,\dd x\dd y$。
\end{exercise}

\begin{exercise}{357}{菱形域的积分恒等式}
设 $f$ 在 $[-1,1]$ 上连续，$D=\{|x|+|y|\leq1\}$。证明
$\iint_Df(x+y)\,\dd x\dd y=\int_{-1}^1f(u)\,\dd u$。
进一步设 $g$ 也连续，推导 $\iint_Df(x+y)g(x-y)\,\dd x\dd y$ 的乘积公式，并检验 $f=g\equiv1$ 时的面积。
\end{exercise}

\begin{exercise}{358}{无界曲边区域的收敛性}
设 $D=\{(x,y):x\geq0,4x^2\leq y\leq9x^2\}$。先以 $0\leq y\leq R$ 截断，计算 $\iint_Dxe^{-y^2}\,\dd x\dd y$ 并证明其收敛；再判断将 $e^{-y^2}$ 换为 $1$ 后积分是否收敛。
\end{exercise}

\begin{exercise}{359}{无界扇形上的高斯积分}
设 $\alpha>0$。证明并计算
$I(\alpha)=\int_0^{+\infty}\dd x\int_x^{\sqrt3x}e^{-\alpha(x^2+y^2)}\,\dd y$，给出 $I(1)$；利用平面缩放独立说明 $I(\alpha)$ 与 $\alpha$ 的关系。
\end{exercise}

\begin{exercise}{360}{第一象限反常积分的参数界限}
设 $p\in\mathbb R$，以半径 $R$ 的第一象限圆盘递增截断定义
$I(p)=\iint_{x,y\geq0}(1+x^2+y^2)^{-p}\,\dd x\dd y$。
求积分收敛的全部 $p$ 及收敛时的值；特别证明 $I(2)=\pi/4$，并分别说明临界值与临界值以下的发散情形。
\end{exercise}

\begin{exercise}{361}{换元区域与雅可比的双重核查}
设 $D=\{0\leq y\leq x\leq1\}$，取 $u=x+y,v=x-y$。有人将 $D$ 的像写成矩形 $[0,2]\times[0,1]$，并把雅可比因子取为 $1$。求真实像区域、逆变换与正确的面积因子，计算 $\iint_D(x+y)\,\dd x\dd y$；再用原坐标积分检验，分别说明两项错误各自改变了什么。
\end{exercise}

\begin{exercise}{362}{圆域奇点的两个临界指数}
设 $p\in\mathbb R$，$D=\{x^2+y^2\leq1\}$。以 $D\cap\{x^2+y^2\geq\varepsilon^2\}$ 截断，分别讨论
\[
 \iint_D(x^2+y^2)^{-p/2}\,\dd\sigma,\qquad
 \iint_D\frac{x^2}{(x^2+y^2)^{p/2}}\,\dd\sigma
\]
的收敛参数范围及收敛时的值；写出临界参数处的截断表达式，解释分子 $x^2$ 如何改变奇性阶数。
\end{exercise}

\subsection{三重积分}

\begin{exercise}{363}{空间区域与幂函数积分}
设 $\Omega=\{(x,y,z):0\leq y\leq x\leq1,\ 0\leq z\leq xy\}$。
写出先对 $z$ 积分与先对 $y$ 积分的两种累次积分，计算
$\iiint_\Omega xy^2z^3\,\dd x\dd y\dd z$，并说明曲面 $z=xy$ 与 $z=0$ 如何确定区域的下边界 $y=0$。
\end{exercise}

\begin{exercise}{364}{三重积分的次序选择}
设 $\Omega=\{(x,y,z):0\leq x\leq\pi/2,\ 0\leq y\leq\sqrt{x},\ 0\leq z\leq\pi/2-x\}$。
将 $\sin x/x$ 在 $x=0$ 处作连续延拓，计算
\[
 \iiint_\Omega\frac{y\sin x}{x}\,\dd x\dd y\dd z,
\]
并解释先对 $y,z$ 积分为何能消去分母 $x$ 带来的计算困难。
\end{exercise}

\begin{exercise}{365}{椭球截面与参数积分}
设 $0<h\leq3$，$\Omega_h=\{(x,y,z):x^2+y^2/4+z^2/9\leq1,\ 0\leq z\leq h\}$。
用水平截面法求 $J(h)=\iiint_{\Omega_h}z^2\,\dd V$，求 $J(1)$，并说明 $J'(h)$ 的几何意义。
\end{exercise}

\begin{exercise}{366}{球面与抛物面的交域}
设 $\Omega$ 由 $z=\sqrt{2-x^2-y^2}$ 与 $z=x^2+y^2$ 围成。求两曲面的交线及区域的柱面坐标范围，计算体积与 $\iiint_\Omega z\,\dd V$，进而求均匀物体的形心高度。
\end{exercise}

\begin{exercise}{367}{抛物面内水平截层}
设 $0<a<b$，$\Omega_{a,b}=\{(x,y,z):a\leq z\leq b,x^2+y^2\leq z\}$。求 $\iiint_{\Omega_{a,b}}(x^2+y^2)^2\,\dd V$，给出 $a=1,b=2$ 的结果；求该函数在区域上的平均值。
\end{exercise}

\begin{exercise}{368}{旋转曲面与参数截面}
将平面 $x=0$ 内的抛物线 $y^2=2z$ 绕 $z$ 轴旋转，所得曲面与平面 $z=H$（$H>0$）围成 $\Omega_H$。求曲面方程与 $\iiint_{\Omega_H}(x^2+y^2+z^2)\,\dd V$，并给出 $H=1$ 的结果。
\end{exercise}

\begin{exercise}{369}{球域幂函数与坐标面}
设 $R>0$，$\Omega_R=\{x^2+y^2+z^2\leq R^2,x,y,z\geq0\}$。计算 $\iiint_{\Omega_R}xyz\,\dd V$，给出 $R=1$ 的结果；若改在整个球域上积分，结果如何？说明区域改变对奇偶性的影响。
\end{exercise}

\begin{exercise}{370}{偏心球与圆锥的交域}
设 $\Omega=\{x^2+y^2+z^2\leq2z,z\geq\sqrt{x^2+y^2}\}$。写出球面坐标范围，计算 $\iiint_\Omega z^3\,\dd V$；再用水平截面说明截面半径在何处改变表达式，写出另一种计算式。
\end{exercise}

\begin{exercise}{371}{半球与卦限的对称性辨析}
设 $R>0$，$\Omega_1=\{x^2+y^2+z^2\leq R^2,z\geq0\}$，$\Omega_2=\Omega_1\cap\{x\geq0,y\geq0\}$。对 $f=x,y,z,xyz$，分别判断
\[
 \iiint_{\Omega_1}f\,\dd V=4\iiint_{\Omega_2}f\,\dd V
\]
是否成立，给出理由；对成立的非零情形，计算两边的共同值。
\end{exercise}

\begin{exercise}{372}{球域的二次型积分}
设 $a,b,c>0$，$\Omega=\{x^2+y^2+z^2\leq1\}$。利用坐标互换求
$\iiint_\Omega(x^2/a^2+y^2/b^2+z^2/c^2)\,\dd V$；进一步求加入混合项 $\lambda xy+\mu yz+\nu zx$ 后的积分，说明这些系数为何不影响结果。
\end{exercise}

\begin{exercise}{373}{球内二次矩与转动惯量}
设 $R>0$，$\Omega=\{x^2+y^2+z^2\leq R^2\}$。求 $\iiint_\Omega y^2\,\dd V$，并由对称性求密度为 $1$ 的实心球关于任一直径的转动惯量；检验结果关于 $R$ 的缩放次数。
\end{exercise}

\begin{exercise}{374}{圆锥截球与平均高度}
设 $\Omega$ 由圆锥面 $z=\sqrt{x^2+y^2}$ 与上半球面 $z=\sqrt{1-x^2-y^2}$ 围成。先说明 $x$ 项的积分是否为零，再求 $\iiint_\Omega(x+z)\,\dd V$ 与该区域的平均高度。
\end{exercise}

\begin{exercise}{375}{径向积分比值的单调性}
设 $f$ 在 $[0,+\infty)$ 上连续且严格为正，$t>0$。令
\[
 F(t)=\frac{\iiint_{x^2+y^2+z^2\leq t^2}f(x^2+y^2+z^2)\,\dd V}
 {\iint_{x^2+y^2\leq t^2}f(x^2+y^2)\,\dd\sigma}.
\]
把 $F(t)$ 化为一元积分之比，证明其在 $(0,+\infty)$ 上严格递增，并求 $\lim_{t\to0^+}F(t)/t$。证明中须明确用到 $f>0$ 的位置。
\end{exercise}

\begin{exercise}{376}{小圆柱上的积分极限}
设 $f$ 在 $0$ 的右邻域连续可微，$f(0)=0$，$H>0$，$\Omega_t=\{x^2+y^2\leq t^2,0\leq z\leq H\}$。求
$\lim_{t\to0^+}t^{-4}\iiint_{\Omega_t}f(x^2+y^2)\,\dd V$，给出 $H=1$ 的结果；用 $f(u)=f'(0)u+o(u)$ 的余项估计证明极限。
\end{exercise}

\begin{exercise}{377}{球冠的坐标范围与截面验证}
设 $0<h<1$，$\Omega_h=\{x^2+y^2+z^2\leq1,z\geq h\}$。分别用水平截面与球面坐标写出体积积分，明确球面坐标中径向变量的下界及极角范围。计算体积与 $\iiint_{\Omega_h}z\,\dd V$，由两种表达式检验结果，并判断 $h\to1^-$ 时平均高度的极限。
\end{exercise}

\begin{exercise}{378}{非均匀抛物体的参数质心}
设 $\Omega=\{x^2+y^2\leq z\leq1\}$，体密度为 $\rho_\alpha=1+\alpha z$。
先求密度在整个区域上非负的全部 $\alpha$。在该范围内求质量与全部质心坐标，证明质心高度随 $\alpha$ 严格增加，求其最小值与 $\alpha\to+\infty$ 时的极限，并确认质量不会为零。
\end{exercise}

\subsection{重积分的应用}

\begin{exercise}{379}{坐标面截三角形的面积}
设 $a,b,c>0$。求平面 $x/a+y/b+z/c=1$ 被三个坐标面截出的三角形面积。分别写出向 $xOy$ 面投影与向 $yOz$ 面投影的面积积分，并验证两种结果相同。
\end{exercise}

\begin{exercise}{380}{柱内球面的面积}
设 $a>0$。求球面 $x^2+y^2+z^2=a^2$ 满足 $x^2+y^2\leq ax$ 部分的面积，明确是否包含上、下半球的部分；再求该面积占整个球面面积的比例，并说明比例为何与 $a$ 无关。
\end{exercise}

\begin{exercise}{381}{球带面积与球面总面积}
设 $R>0$，$-R\leq a<b\leq R$。用曲面面积积分求球面 $x^2+y^2+z^2=R^2$ 上 $a\leq z\leq b$ 的球带面积，进而推导整球表面积；若积分表达式在极点或赤道出现无穷因子，说明相应极限的处理方法。
\end{exercise}

\begin{exercise}{382}{曲边薄片的质量与质心}
设 $D=\{(x,y):0\leq x\leq1,x^2\leq y\leq x\}$，面密度为 $\mu(x,y)=x^2y$。求质量及质心坐标；再求同一区域均匀薄片的形心，比较两个横坐标，并从密度随 $x$ 的变化解释结果。
\end{exercise}

\begin{exercise}{383}{偏心球体的非均匀质心}
设球体 $\Omega=\{x^2+y^2+z^2\leq z\}$ 的体密度为 $\rho=x^2+y^2+z^2$。利用对称性确定质心的两个坐标，计算质量及质心竖坐标，并比较它与均匀球体的球心高度。
\end{exercise}

\begin{exercise}{384}{抛物体截面的形心}
设 $H>0$，$\Omega_H=\{(x,y,z):x^2+y^2\leq z\leq H\}$。求其体积与全部形心坐标，给出 $H=1$ 的结果；解释形心高度为何不能直接取上下端高度的算术平均值。
\end{exercise}

\begin{exercise}{385}{斜圆锥的截面与形心}
设 $\Omega=\{(x,y,z):0\leq z\leq1,x^2+(y-z)^2\leq(1-z)^2\}$ 为均匀锥体。求水平截面的面积及圆心，利用截面法求全部形心坐标，并证明形心位于底面圆心与顶点的连线上。
\end{exercise}

\begin{exercise}{386}{椭圆薄片的惯量}
设 $a,b>0$，$D=\{x^2/a^2+y^2/b^2\leq1\}$，面密度为 $1$。求薄片关于 $y$ 轴、$x$ 轴的转动惯量与关于原点的极惯量；比较前两者并给出相等的充要条件。
\end{exercise}

\begin{exercise}{387}{平行直线的惯量最小值}
设 $D=\{(x,y):-1\leq x\leq1,x^2\leq y\leq1\}$，面密度为常数 $\mu>0$。求关于直线 $y=c$ 的转动惯量 $I(c)$，给出 $c=-1$ 的结果，并求使 $I(c)$ 最小的 $c$，说明其与质心的关系。
\end{exercise}

\begin{exercise}{388}{抛物面下立体的轴惯量}
设 $a,\rho>0$，$\Omega=\{(x,y,z):|x|\leq a,|y|\leq a,0\leq z\leq x^2+y^2\}$，体密度为 $\rho$。求质量及关于 $z$ 轴的转动惯量，给出惯量与质量之比，并检验其量纲。
\end{exercise}

\begin{exercise}{389}{圆柱体的轴向引力}
设 $R,h>0$，均匀圆柱体 $\Omega=\{x^2+y^2\leq R^2,0\leq z\leq h\}$ 的体密度为 $1$，引力常数为 $G$。求其对 $M_0(0,0,a)$（$a>h$）处单位质量质点的引力向量；先用对称性确定方向，再计算非零分量，并检验 $a\to+\infty$ 时引力大小与 $G\pi R^2h/a^2$ 的比值。
\end{exercise}

\begin{exercise}{390}{曲面面积的投影因子与边界极限}
设 $S:z=\sqrt{1-x^2-y^2}$ 为单位上半球面。解释为何投影圆盘面积不能直接作为 $S$ 的面积。
先在 $x^2+y^2\leq(1-\varepsilon)^2$（$0<\varepsilon<1$）上计算对应球面面积，再令 $\varepsilon\to0^+$；指出面积因子在赤道附近虽无界，面积仍有限的具体积分依据。
\end{exercise}

\begin{exercise}{391}{椭圆薄片的非对称密度与中心惯量}
设 $a,b>0$，$D=\{x^2/a^2+y^2/b^2\leq1\}$，面密度为 $1+\lambda x/a$。
求使密度非负的全部 $\lambda$，在该范围内求质量、质心、关于 $y$ 轴的惯量，以及关于过质心且平行于 $y$ 轴的直线的惯量。证明最后一个惯量的最大值与最小值分别在何处取得，并解释质量为何不随 $\lambda$ 改变。
\end{exercise}

\clearpage
\section{第11章\quad 曲线积分}
\subsection{对弧长的曲线积分}

\begin{exercise}{392}{弧长积分与取向}
设分段光滑曲线弧 $L$ 的长度为 $a>0$，$k$ 为常数。从弧长积分的定义求 $\int_Lk\,\dd s$，并判断将曲线反向后结果是否改变；若已知 $\int_L(k+1)\,\dd s=3a$，求 $k$。
\end{exercise}

\begin{exercise}{393}{弧长积分的保序性}
设 $L$ 为长度非零的分段光滑曲线弧，$f$ 在 $L$ 上连续且非负。证明 $\int_Lf\,\dd s\geq0$；进一步证明，若 $f$ 在 $L$ 的某个内点处为正，则积分严格为正，并据此给出积分等于零的充要条件。
\end{exercise}

\begin{exercise}{394}{积分比较与严格不等式}
设 $f,g$ 在长度非零的分段光滑曲线弧 $L$ 上连续，且 $f\geq g$。证明
$\int_Lf\,\dd s\geq\int_Lg\,\dd s$。
判断仅有“某一点 $f>g$”能否推出严格不等式，并说明连续性在判断中的作用。
\end{exercise}

\begin{exercise}{395}{积分绝对值不等式的等号}
设实值连续函数 $f$ 定义在长度非零的空间光滑曲线弧 $L$ 上。证明
$|\int_Lf\,\dd s|\leq\int_L|f|\,\dd s$，并给出等号成立的充要条件；用单位圆上的 $f(x,y)=x$ 说明不等式可以严格成立。
\end{exercise}

\begin{exercise}{396}{圆弧上径向函数的积分}
设 $R>0$，$L_R$ 为下半圆 $y=-\sqrt{R^2-x^2}$。求 $\int_{L_R}\sqrt{x^2+y^2}\,\dd s$ 及径向距离的平均值，给出 $R=2$ 的结果，并解释两者随 $R$ 的缩放次数。
\end{exercise}

\begin{exercise}{397}{线段上的仿射函数平均值}
设 $L$ 连接 $(-1,0)$ 与 $(0,2)$。计算 $\int_L(2x-y+5)\,\dd s$；证明任意仿射函数在线段上的弧长平均值等于其在中点处的值，并用该结论检验结果。
\end{exercise}

\begin{exercise}{398}{水平圆周与参数积分}
设 $R>0,h\in\mathbb R$，$L$ 为 $x^2+y^2=R^2,z=h$ 的圆周。求 $\oint_L\arctan\frac{z^2}{x^2+y^2}\,\dd s$，给出 $R=h=2$ 的结果；固定 $R$ 时讨论积分随 $|h|$ 的单调性与极限。
\end{exercise}

\begin{exercise}{399}{大圆上的轮换对称性}
设 $\Gamma$ 为球面 $x^2+y^2+z^2=9$ 与平面 $x+y+z=0$ 的交圆。求 $\oint_\Gamma z^2\,\dd s$，再求 $\oint_\Gamma xy\,\dd s$；说明坐标轮换与平面方程分别提供了什么关系。
\end{exercise}

\begin{exercise}{400}{椭圆约束下的二次函数}
设椭圆 $L:x^2/4+y^2/3=1$ 的周长为 $a$。求所有使 $\lambda x^2+\mu y^2$ 在 $L$ 上为常数的 $(\lambda,\mu)$，据此求 $\oint_L(2xy+3x^2+4y^2)\,\dd s$。不要求求出周长 $a$。
\end{exercise}

\begin{exercise}{401}{折线对称性与积分比例}
设 $L_1:y=|x|$（$-1\leq x\leq1$），$L_2:y=x$（$0\leq x\leq1$）。分别计算
$I_1=\int_{L_1}(xy^2-|x|+2y)\,\dd s$ 与 $I_2=\int_{L_2}x\,\dd s$，求其比例；将 $2y$ 换成 $cy$ 后，求使新 $I_1=0$ 的 $c$。
\end{exercise}

\begin{exercise}{402}{线段上函数积分的排序}
设 $L$ 连接 $(1,0)$ 与 $(0,1)$，
$I=\int_Lx^2\,\dd s,J=\int_L(x-y)\,\dd s,K=\int_L(x+y)/2\,\dd s$。
求三者并严格排序；若 $I_p=\int_Lx^p\,\dd s$（$p>0$），求使 $J<I_p<K$ 成立的全部 $p$。
\end{exercise}

\begin{exercise}{403}{互换对称性与分母条件}
设 $0<R<(\sqrt3/2)\ln\pi$，$L:x^2+xy+y^2=R^2$ 的周长为 $a$。证明 $\sin(e^x)+\sin(e^y)>0$ 在 $L$ 上成立，求
$\oint_L\frac{2\sin(e^x)-3\sin(e^y)}{\sin(e^x)+\sin(e^y)}\,\dd s$。
说明去掉对 $R$ 的限制后，能否仅凭对称性就忽略积分存在性。
\end{exercise}

\begin{exercise}{404}{球面截圆上的交叉矩}
设 $L$ 为 $x^2+y^2+z^2=1$ 与 $x+y+z=0$ 的交圆。求 $\oint_L(xy+z)\,\dd s$，并求任意常数 $a,b,c$ 下 $\oint_L(ax y+byz+czx)\,\dd s$。解释线性项与交叉二次项所用对称性为何不同。
\end{exercise}

\begin{exercise}{405}{指数曲线与弧长因子}
设 $b>0$，$L_b:y=e^x$（$0\leq x\leq b$）。求 $\int_{L_b}\sqrt{1+y^2}\,\dd s$，给出 $b=1$ 的结果；判断其是否等于 $L_b$ 的长度并给出严格比较。
\end{exercise}

\begin{exercise}{406}{抛物线参数与弧长平均}
设 $b>0$，$L_b:y^2=2x$（$0\leq y\leq b$）。计算 $\int_{L_b}y\,\dd s$，给出 $b=\sqrt3$ 的结果；写出曲线长度的一元积分式，证明 $y$ 的弧长平均值位于 $(0,b)$。
\end{exercise}

\begin{exercise}{407}{双曲线弧上的因子约简}
设 $0<a<1$，$L_a$ 为 $xy=1$ 上 $(a,1/a)$ 至 $(1,1)$ 的弧段。计算 $\int_{L_a}\frac{x+y}{\sqrt{x^2+y^2}}\,\dd s$，给出 $a=1/2$ 的结果；用 $x$ 与用 $y$ 作参数分别写出积分，验证换向不改变结果。
\end{exercise}

\begin{exercise}{408}{平移圆弧与角度参数}
设 $a>0,0<b\leq\pi$，$L:x=a(1+\cos t),y=a\sin t$（$0\leq t\leq b$）。求 $\int_Lxy\,\dd s$，给出 $b=\pi/2$ 的结果；固定 $b$，说明积分关于 $a$ 的缩放次数。
\end{exercise}

\begin{exercise}{409}{极坐标闭曲线与几何中心}
设 $L$ 的极坐标方程为 $r=\cos\theta+\sin\theta$，$-\pi/4\leq\theta\leq3\pi/4$。求其直角坐标方程、圆心和半径，计算 $\oint_L(x+y)\,\dd s$；再求 $x+y$ 的弧长平均值并用圆心坐标解释。
\end{exercise}

\begin{exercise}{410}{空间螺旋弧与收敛极限}
设 $T>0$，$\Gamma_T:x=e^t\cos t,y=e^t\sin t,z=e^t$（$0\leq t\leq T$）。求
$\int_{\Gamma_T}(x^2+y^2+z^2)^{-1}\,\dd s$，给出 $T=2$ 的结果；讨论 $T\to+\infty$ 时曲线长度与该积分是否各有有限极限。
\end{exercise}

\begin{exercise}{411}{对数曲线的参数选择}
设 $b>1$，$L_b:y=\ln x$（$1\leq x\leq b$）。求 $\int_{L_b}\frac{e^y}{1+x^2}\,\dd s$，给出 $b=\sqrt3$ 的结果；证明该积分随 $b$ 严格增加，并求 $b\to+\infty$ 时的极限。
\end{exercise}

\begin{exercise}{412}{椭圆弧与根式抵消}
设 $L$ 为椭圆 $x^2+4y^2=2x+3$ 在 $x$ 轴上方的部分。选取参数使 $\sqrt{1+3y^2}\,\dd s$ 化简，计算其积分；若改取完整椭圆，结果如何？说明对称性依据。
\end{exercise}

\begin{exercise}{413}{双纽线的极坐标分段}
设 $a>0$，$L:(x^2+y^2)^2=a^2(x^2-y^2)$ 为完整双纽线。计算 $\oint_L(|x|-|y|)\,\dd s$，写明两叶曲线的角度范围，并从曲线方程证明被积函数处处非负。
\end{exercise}

\begin{exercise}{414}{空间线段的弧长积分}
设 $L$ 连接 $A(-1,0,2)$ 与 $B(3,-2,1)$。用同一个参数分别计算 $\int_Lxyz\,\dd s$ 与 $\int_L|xyz|\,\dd s$，明确第二个积分的分段点，并判断二者是否互为绝对值。
\end{exercise}

\begin{exercise}{415}{空间截圆上的二次矩}
设 $|c|<2\sqrt3$，$L_c$ 为 $x^2+y^2+z^2=6$ 与 $x+z=c$ 的交圆。求圆心、半径及 $\oint_{L_c}(xz+y^2)\,\dd s$，给出 $c=2$ 的结果；可取沿平面内两条正交直径的坐标简化二次项。
\end{exercise}

\begin{exercise}{416}{封闭边界与线密度质量}
设 $L$ 为直线 $y=2x$ 与抛物线 $y^2=x$ 围成区域的完整边界。求两交点，分段计算 $\oint_Ly\,\dd s$；把 $y$ 视为线密度时，求两段细丝的质量比例，并说明端点密度为零是否妨碍总质量为正。
\end{exercise}

\begin{exercise}{417}{折线上的弧长积分}
设 $L$ 为顶点 $O(0,0),A(1,0),B(0,1)$ 的三角形边界。分边计算
$\oint_L\frac{x}{1+y}\,\dd s$，并求密度为 $\rho(x,y)=1/(1+y)$ 的三角形细丝的质心横坐标。
\end{exercise}

\begin{exercise}{418}{弧长平均值与平方积分}
设 $L$ 为长度 $\ell>0$ 的光滑曲线弧，$f$ 为实值连续函数。证明
$|\int_Lf\,\dd s|^2\leq\ell\int_Lf^2\,\dd s$，给出等号成立的充要条件。
比较“$\int_Lf\,\dd s=0$”与“$\int_Lf^2\,\dd s=0$”分别能推出什么；用单位圆上的一个具体连续函数说明两条件不能互相替代。
\end{exercise}

\begin{exercise}{419}{圆形细丝的非均匀密度}
设 $L:x^2+y^2=1$，线密度为 $\rho_\lambda(x,y)=1+\lambda x$。
求使密度处处非负的全部 $\lambda$，在该范围内求质量、质心及关于 $y$ 轴的转动惯量。进一步求过质心且平行于 $y$ 轴的直线的惯量，讨论其随 $|\lambda|$ 的变化，并说明曲线反向是否影响各物理量。
\end{exercise}

\subsection{对坐标的曲线积分}

\begin{exercise}{420}{坐标微分与弧长微分}
设 $L$ 为从 $(c,a)$ 到 $(c,b)$ 的有向线段，其中 $a<b$，$P,Q$ 连续。求
$\int_LP(x,y)\,\dd x$ 和 $\int_LQ(x,y)\,\dd y$ 的一元积分表达式；再取 $P\equiv1$，比较 $\int_LP\,\dd x$ 与 $\int_LP\,\dd s$。
\end{exercise}

\begin{exercise}{421}{常向量场与端点}
设 $L$ 为由 $A(x_1,y_1)$ 到 $B(x_2,y_2)$ 的分段光滑曲线，$k,m$ 为常数。求
$\int_Lk\,\dd x$、$\int_Lk\,\dd y$ 及 $\int_L(k\,\dd x+m\,\dd y)$；据此说明常向量场所做的功与路径无关，并求反向积分。
\end{exercise}

\begin{exercise}{422}{三角根式与绝对值}
设 $0<b\leq\pi$，$L_b:y=\sin x$ 从 $(0,0)$ 到 $(b,\sin b)$。求 $I(b)=\int_{L_b}\sqrt{1-y^2}\,\dd x$ 的分段表达式，给出 $I(\pi)$，并指出直接以 $\cos x$ 代替根式会在哪个区间产生错误。
\end{exercise}

\begin{exercise}{423}{抛物线上的奇函数消去}
设 $a>0$，$L_a:y=x^2$ 从 $(-a,a^2)$ 到 $(a,a^2)$。计算 $\int_{L_a}(x+e^y)\,\dd y$，给出 $a=1$ 的结果；分别说明 $x$ 项与 $e^y$ 项是否为零，不能仅根据端点纵坐标相同判断整个积分。
\end{exercise}

\begin{exercise}{424}{线段参数与方向符号}
设 $a,b>0$，$L$ 为从 $(a,0)$ 到 $(0,b)$ 的线段。求 $\int_Lxy\,\dd y$ 与 $\int_Lxy\,\dd x$，给出 $a=3,b=1$ 时前者的值；解释两积分符号不同的原因。
\end{exercise}

\begin{exercise}{425}{双曲线上的对数积分}
设 $b>1$，$L$ 为 $xy=1$ 上从 $(1,1)$ 到 $(b,1/b)$ 的弧。求 $\int_Lx\ln y\,\dd y$，给出 $b=2$ 的结果；用起终点及 $\ln y$ 的符号独立判断结果的正负。
\end{exercise}

\begin{exercise}{426}{曲线约束与恰当微分}
设 $L:y=\cos x$ 从 $(0,1)$ 到 $(\pi,-1)$。计算 $\int_Le^y\sin x\,\dd x+(1-y)\,\dd y$；判断该微分形式在整个平面内是否与路径无关，解释本题沿给定曲线仍可快速求值的原因。
\end{exercise}

\begin{exercise}{427}{回折抛物线的参数选择}
设 $L:y^2=x$ 从 $(1,-1)$ 到 $(1,1)$。取 $y$ 为参数计算 $\int_L\frac{xy\,\dd x-x^2\,\dd y}{1+x}$；说明直接以 $x$ 为全程单调参数会遗漏什么，并写出以 $x$ 分段时的两段方向。
\end{exercise}

\begin{exercise}{428}{摆线一拱与端点项}
设 $a>0$，$L:x=a(t-\sin t),y=a(1-\cos t)$，$t:0\to2\pi$。计算 $\int_L(2a-y)\,\dd x+\dd y$，并求 $\int_Ly\,\dd x$；用两者之和与水平位移检验结果。
\end{exercise}

\begin{exercise}{429}{螺旋线上的多项式三角积分}
设 $T>0$，$L_T:x=2\cos t,y=2\sin t,z=t$，$t:0\to T$。求 $\int_{L_T}y\,\dd x+z\,\dd y-x^2\,\dd z$，给出 $T=\pi$ 的结果；另求将路径改为同端点直线段后的积分，判断是否与路径无关。
\end{exercise}

\begin{exercise}{430}{椭圆弧上的端点法}
设 $L$ 为椭圆 $x^2+4y^2=4$ 在第一象限内按逆时针方向的弧。找出 $y\,\dd x+x\,\dd y$ 的原函数，计算积分；若终点改为该弧上任意点 $(u,v)$，写出积分值并求其最大值。
\end{exercise}

\begin{exercise}{431}{角形式与绕行次数}
设 $a>0$，$L$ 为顺时针圆周 $x^2+y^2=a^2$。参数化计算 $\oint_L\frac{y\,\dd x-x\,\dd y}{x^2+y^2}$；若沿同一圆周顺时针走 $m$ 圈（$m$ 为正整数），结果如何？说明为何不能在整个圆盘上直接使用无奇点的格林公式。
\end{exercise}

\begin{exercise}{432}{圆周约束与分母退化}
设 $r>0$，$L_r:(x-1)^2+(y-2)^2=r^2$ 取逆时针方向。讨论 $\oint_{L_r}\frac{\dd x+\dd y}{x^2+y^2-2x-4y}$ 何时有定义，有定义时计算，给出 $r=3$ 的结果；特别说明 $r=\sqrt5$ 的情形。
\end{exercise}

\begin{exercise}{433}{空间径向场与路径替换}
设 $L$ 从 $A(1,2,3)$ 到 $B(-1,4,0)$。证明 $\int_Lx\,\dd x+y\,\dd y+z\,\dd z$ 与路径无关并计算；若终点可在球面 $x^2+y^2+z^2=17$ 上任意变化，积分是否改变？
\end{exercise}

\begin{exercise}{434}{斜截圆的观察方向}
设 $\Gamma:x^2+y^2=1,x+z=1$，从 $x$ 轴正方向观察为顺时针。写出符合方向的参数方程，计算
$\oint_\Gamma yz\,\dd x-\frac{x}{x^2+y^2}\,\dd y-\frac{y}{x+z}\,\dd z$。
分别核查两个分母在曲线上是否为零，并给出反向后的值。
\end{exercise}

\begin{exercise}{435}{折线拐点与分段积分}
设 $L:y=1-|x|$，从 $(-1,0)$ 到 $(1,0)$。计算 $\int_Ly^2\,\dd x+x^2\,\dd y$；把顶点 $(0,1)$ 改为 $(0,h)$（$h>0$）后求积分关于 $h$ 的表达式，说明两侧 $\dd y$ 的符号。
\end{exercise}

\begin{exercise}{436}{三角形边界上的指数场}
设 $h>0$，$L_h$ 为依次经过 $O(0,0),A(1,0),B(1,h),O$ 的闭折线。分边求 $\oint_{L_h}e^y\,\dd x+e^x\,\dd y$，给出 $h=2$ 的结果；说明不能只根据闭合性将积分判为零。
\end{exercise}

\begin{exercise}{437}{复合圆弧与全微分}
设 $L$ 从 $O(0,0)$ 沿第一象限的圆弧 $x^2+y^2=2x$ 到 $A(2,0)$，再沿 $x^2+y^2=4$ 的第一象限圆弧到 $B(0,2)$。求 $\int_Ly\,\dd x+(x-1)\,\dd y$，分别用分段法和原函数法计算；说明若中间点及路径改变但端点不变，结果是否保持。
\end{exercise}

\begin{exercise}{438}{角点的单侧极限与曲线积分}
设 $L$ 从原点沿 $y=\sqrt x$ 到 $(1,1)$，再沿 $y=x$ 回原点。分段计算 $\oint_L\arctan(x/y)\,\dd x-\dd y$，各段端点按单侧极限处理；判断被积函数能否在原点作连续延拓，说明仅改原点处的有限取值是否影响积分。
\end{exercise}

\begin{exercise}{439}{空间折线上的坐标微分}
设 $L$ 依次经过 $A(0,0,0),B(0,0,2),C(1,0,2),D(1,3,2)$。求 $\int_Lz^2\,\dd x+x^2\,\dd y+y^2\,\dd z$；再按 $A\to(1,0,0)\to(1,3,0)\to D$ 的顺序移动，求新积分并比较，指出哪一项体现了路径影响。
\end{exercise}

\begin{exercise}{440}{空间交圆的平面化}
设 $L:x^2+y^2+z^2=6,y+z=2$，从 $z$ 轴正方向观察为逆时针。利用 $y+z=2$ 将
$\oint_Lz\,\dd x+x\,\dd y-\frac{y+z}{2}\,\dd z$ 化为平面投影上的积分，求值并说明投影方向；再用一组正交参数检验符号。
\end{exercise}

\begin{exercise}{441}{折点路径与全微分}
设 $L$ 为圆锥面 $z=\sqrt{x^2+y^2}$ 与平面 $y=x$ 的交线上从 $A(1,1,\sqrt2)$ 经原点到 $B(-1,-1,\sqrt2)$ 的曲线弧。
说明 $L$ 在原点处的光滑性，并分别用分段参数法和全微分法计算
$\int_Lyz\,\dd x+zx\,\dd y+xy\,\dd z$。
\end{exercise}

\begin{exercise}{442}{有向单位切向量}
设 $L$ 是圆周 $x^2+y^2=2y$ 上从原点到 $(0,2)$ 的右半圆弧，$P,Q$ 连续。
求与 $L$ 同向的单位切向量，将 $\int_LP\,\dd x+Q\,\dd y$ 化为对弧长的积分；再取 $P=x,Q=y$ 计算，并检验反向后的符号。
\end{exercise}

\begin{exercise}{443}{两类曲线积分与方向}
设 $L$ 为 $x=t,y=t^2,z=t^3$ 上 $t$ 从 $1$ 变到 $0$ 的曲线弧，弧长为 $a$。
不用求出 $a$ 的具体数值，计算
\[
 \int_L\frac{\dd x+2x\,\dd y+3y\,\dd z}{\sqrt{1+4x^2+9y^2}}.
\]
再将三个坐标微分的系数解释为沿正参数方向的单位切向量分量，说明取向如何决定结果。
\end{exercise}

\begin{exercise}{444}{抛物线弧上的积分转换}
设 $f$ 连续，$L$ 为 $y=x^2$ 上从 $(1,1)$ 到 $(0,0)$ 的曲线弧。
\begin{enumerate}[label=(\arabic*)]
\item 把 $I=\int_L f(x,y)\,\dd x+2xf(x,y)\,\dd y$ 转化为对弧长的积分，写明符号的来源。
\item 当 $f(x,y)=1/(x^2+3y+1)$ 时计算 $I$。
\item 当 $f$ 在 $L$ 上非负时，判断 $I$ 的符号；给出等号成立的条件。
\end{enumerate}
\end{exercise}

\begin{exercise}{445}{非单调参数与重复经过的弧段}
设 $\gamma(u)=(u,u^2)$（$0\leq u\leq1$）。曲线 $C$ 由 $u=4t(1-t)$、$t:0\to1$ 描述，因此先沿抛物线从原点走到 $(1,1)$，再原路返回。
分别计算 $\int_Cx\,\dd s$ 与 $\int_C(-y\,\dd x+x\,\dd y)$，比较它们与只走去程的相应积分。解释第一类积分中的参数速度为何取绝对值，以及为何不能把 $u(0)=u(1)$ 直接代入第一类积分的积分限判为零。
\end{exercise}

\begin{exercise}{446}{路径族与保守场参数}
设 $\alpha\in\mathbb R,p\geq1$，$L_p:y=x^p$ 从 $(0,0)$ 到 $(1,1)$。求
$W_\alpha(p)=\int_{L_p}\alpha y\,\dd x+x\,\dd y$。
求使 $W_\alpha(p)$ 对所有 $p\geq1$ 均相同的 $\alpha$，证明相应向量场在整个平面内保守；对其余 $\alpha$ 讨论 $W_\alpha$ 关于 $p$ 的单调性与 $p\to+\infty$ 时的极限，并解释路径趋近折线时的结果。
\end{exercise}

\subsection{格林公式及其应用}
\subsubsection{封闭曲线、补线与奇点}

\begin{exercise}{447}{闭曲线积分与全微分}
设 $L$ 为任意分段光滑有向闭曲线。求 $\oint_L(y-\cos x)\,\dd x+(x+\sin y)\,\dd y$，并找出原函数；说明即使闭曲线自交，原函数法是否仍然适用。
\end{exercise}

\begin{exercise}{448}{环流数据与面积确定}
设 $D$ 为有界区域，边界分段光滑且取正向，已知 $\oint_{\partial D}(1+3y)\,\dd x+(5x-1)\,\dd y=4$。求 $D$ 的面积；把 $3,5$ 分别换为参数 $a,b$ 后，讨论何时可以由同一积分数据唯一确定面积，何时数据本身不可能成立。
\end{exercise}

\begin{exercise}{449}{负向三角形边界}
设 $a,b>0$，$L$ 为顶点 $(0,0),(a,0),(a,b)$ 的三角形边界，取顺时针方向。求 $\oint_L(2x+y-4)\,\dd x+(5x-3y+6)\,\dd y$，给出 $a=3,b=2$ 的结果；指出其中哪些项对任意闭曲线的积分均为零。
\end{exercise}

\begin{exercise}{450}{双曲线区域的负向环流}
设 $D$ 由 $xy=1,y=x,x=2$ 围成，$L$ 为其负向边界。计算 $\oint_L(2x+y^2)\,\dd x+(3y-x^2)\,\dd y$，并由格林公式中的被积函数及取向先判断符号；另写出正向边界的结果。
\end{exercise}

\begin{exercise}{451}{圆周环流与四次缩放}
设 $R>0$，$L$ 为正向圆周 $x^2+y^2=R^2$。求 $\oint_Lx^2y\,\dd x-xy^2\,\dd y$，并求积分与圆盘面积之比；分别从变量缩放与格林公式解释结果关于 $R$ 的次数。
\end{exercise}

\begin{exercise}{452}{扇环边界与反正切场}
设 $0<r<R$，$D=\{(\rho\cos\theta,\rho\sin\theta):r\leq\rho\leq R,\pi/4\leq\theta\leq\pi/3\}$。计算其正向边界上的
$\oint_{\partial D}\frac1x\arctan\frac yx\,\dd x+\frac2y\arctan\frac xy\,\dd y$，给出 $r=1,R=2$ 的结果；检查坐标轴上的奇点为何不妨碍使用格林公式。
\end{exercise}

\begin{exercise}{453}{抛物线弓形的面积与环流}
设 $L$ 为 $y^2=2x,y=x-4$ 围成区域的顺时针边界。求交点、区域面积和形心纵坐标，据此计算 $\oint_L(e^x+2y)\,\dd x-(2e^y-xy)\,\dd y$，说明指数项对闭合积分的贡献。
\end{exercise}

\begin{exercise}{454}{菱形域与奇函数消去}
设 $a>0$，$L_a$ 为菱形 $|x|+|y|\leq a$ 的逆时针边界。计算 $\oint_{L_a}(x^2y-xy^3)\,\dd x+(y^2-2xy)\,\dd y$，给出 $a=1$ 的结果；写出格林公式转换后各项被保留或消去的依据。
\end{exercise}

\begin{exercise}{455}{补线与对称抵消}
设 $h>0$，$L_h$ 依次连接 $(-1,0),(0,h),(1,0)$。用补线法求 $\int_{L_h}xy\,\dd x+x^2\,\dd y$，给出 $h=1$ 的结果；再用两段参数验证，并解释其与 $h$ 无关的原因。
\end{exercise}

\begin{exercise}{456}{半圆弧补线与端点项}
设 $L:y=\sqrt{4x-x^2}$ 从 $O(0,0)$ 到 $A(4,0)$。计算 $\int_L(x^2-3y)\,\dd x+(2y^3+x)\,\dd y$，分别标明原函数项与面积项的贡献；若改走下半圆且端点顺序不变，两结果相差多少？
\end{exercise}

\begin{exercise}{457}{曲线约束下的分母简化}
设 $c>0$，$L:y^2=2x$ 从 $(2,2)$ 到 $(2,-2)$。求
$\int_L\frac{(y^2+e^x)\,\dd x+(x^2+y^4)\,\dd y}{y^2-2x+c}$，给出 $c=3$ 的结果；判断 $c=0$ 时能否沿同一路径定义该积分。
\end{exercise}

\begin{exercise}{458}{挖去奇点后的格林公式}
设 $L$ 为 $y^2=x+3,x=2$ 围成区域的逆时针边界。确认原点相对区域的位置，使用挖去小圆的方法计算 $\oint_L\frac{x\,\dd y-y\,\dd x}{x^2+y^2}$；说明仅计算区域内的偏导数差为零为何不足以判定积分为零。
\end{exercise}

\begin{exercise}{459}{移动圆周与奇点穿越}
设 $R>0$，$L_R:(x-1)^2+y^2=R^2$ 取顺时针方向。对 $R<1,R=1,R>1$ 分别讨论
$\oint_{L_R}\frac{y\,\dd x-x\,\dd y}{4x^2+y^2}$ 的存在性并求有定义时的值；可作线性变换把分母化为平方和，须说明方向是否改变。
\end{exercise}

\subsubsection{路径无关、原函数与全微分方程}

\begin{exercise}{460}{参数判定与定义域}
在单位圆盘 $D=\{x^2+y^2<1\}$ 内，求使 $\frac{ax\,\dd x+2y\,\dd y}{x^2+y^2-1}$ 与路径无关的全部 $a$，并求归一化原函数 $u(0,0)=0$。说明单位圆周为何不能包含在该原函数的定义域内。
\end{exercise}

\begin{exercise}{461}{独立函数项的系数匹配}
求使 $(ax^2\cos y-y^3\sin x)\,\dd x+(by^2\cos x+x^3\sin y)\,\dd y$ 在平面内与路径无关的 $a,b$，并找出一个原函数；对所得原函数求偏导数验证，不能只列出系数条件。
\end{exercise}

\begin{exercise}{462}{任意闭曲线积分为零的条件}
设 $F\in C^1(\mathbb R^2)$。证明对任意分段光滑闭曲线 $\Gamma$ 有 $\oint_\Gamma F(x,y)(y\,\dd x+x\,\dd y)=0$ 的充要条件为 $xF_x=yF_y$；分别判断 $F=x+y$ 与 $F=e^{xy}$ 是否满足，并对满足者求一个原函数。
\end{exercise}

\begin{exercise}{463}{分母直线与分支原函数}
求使 $\frac{(ay-x)\,\dd x-y\,\dd y}{(x+y)^2}$ 为全微分的 $a$，并在半平面 $x+y>0$ 内求一个原函数。说明该原函数的表达式如何扩展到 $x+y<0$，以及为何两半平面的积分常数可独立选取。
\end{exercise}

\begin{exercise}{464}{复合原函数与端点公式}
设 $f\in C^1(\mathbb R)$，$L$ 从 $A(a_1,a_2)$ 到 $B(b_1,b_2)$。证明 $\int_Lf(xy)(y\,\dd x+x\,\dd y)$ 与路径无关，并写出仅含一元定积分的端点公式；取 $f(u)=1/(1+u^2)$，$A=(1,0),B=(2,1)$ 计算。
\end{exercise}

\begin{exercise}{465}{由保守场反求未知函数}
设 $f\in C^1(\mathbb R)$，$f(0)=1$，且 $[xy^2-f(x)]\,\dd x+yf(x)\,\dd y$ 与路径无关。求 $f$ 与一个原函数，并计算从 $(0,0)$ 到 $(1,1)$ 的积分；说明初值在唯一确定 $f$ 时的作用。
\end{exercise}

\begin{exercise}{466}{齐次函数与区域条件}
设 $f\in C^1(D)$，$D=\{y>0\}$，且 $f(tx,ty)=t^{-2}f(x,y)$（$t>0$）。证明任意 $D$ 内简单闭曲线 $L$ 均满足 $\oint_Lf(y\,\dd x-x\,\dd y)=0$。再取 $f=1/(x^2+y^2)$，说明把定义域换为去掉原点的整个平面后，同一结论是否仍成立。
\end{exercise}

\begin{exercise}{467}{识别乘积型原函数}
设 $L:x^2/4+y^2/9=1$ 取正向。求 $\oint_L(3x^2\cos x-x^3\sin x)y\,\dd x+(x^3\cos x+2y^2)\,\dd y$，并求从 $(2,0)$ 到 $(0,3)$ 的任意曲线路径积分；必须写出原函数。
\end{exercise}

\begin{exercise}{468}{含非初等项的端点消去}
设 $L:y=\sqrt{2x-x^2}$ 沿上半圆从 $(0,0)$ 到 $(2,0)$，即顺时针方向。证明 $\omega=(e^x+2x\sin y)\,\dd x+(x^2+e^{-y^2})\cos y\,\dd y$ 与路径无关，允许用一元定积分表示原函数，并求 $\int_L\omega$。解释为何计算结果无需求出 $e^{-y^2}\cos y$ 的初等原函数。
\end{exercise}

\begin{exercise}{469}{多项式原函数与端点数据}
设 $L$ 从 $A(1,2)$ 到 $B(-1,1)$。证明 $(x^3-2xy-y^2)\,\dd x-(x+y)^2\,\dd y$ 是全微分并计算其积分；若路径增加一段闭合绕行，说明积分是否变化。
\end{exercise}

\begin{exercise}{470}{径向指数场与路径替换}
设 $L:y=x^3/2$ 从 $(0,0)$ 到 $(2,4)$。计算 $\int_Le^{x^2+y^2}(x\,\dd x+y\,\dd y)$，并求同端点任意分段光滑路径的积分；推广写出 $e^u$ 换为连续函数 $g(u)$ 后的端点公式。
\end{exercise}

\begin{exercise}{471}{线性场的保守性}
证明 $\int_{(-1,-2)}^{(3,4)}(x+2y)\,\dd x+(2x-y)\,\dd y$ 与路径无关并求值。再把第二项中的 $2x$ 改为 $kx$，求仍保持路径无关的 $k$，并说明必要性。
\end{exercise}

\begin{exercise}{472}{混合项配对与原函数}
求 $\int_{(2,0)}^{(4,3)}(xy^2+x-y)\,\dd x+(x^2y-x+y)\,\dd y$。将原函数整理成两个平方项之和，并据此求从 $(0,0)$ 到任意 $(x,y)$ 的积分何时为零。
\end{exercise}

\begin{exercise}{473}{星形线面积与伸缩}
设 $a,b>0$，$L:x=a\cos^3t,y=b\sin^3t$，$0\leq t\leq2\pi$。用曲线积分求围成区域的面积，给出 $a=b$ 时的结果；说明四个尖点是否妨碍按分段光滑边界使用格林公式。
\end{exercise}

\begin{exercise}{474}{双纽线两叶的有向面积}
设 $a>0$，双纽线 $L:(x^2+y^2)^2=a^2(x^2-y^2)$ 的两叶均取正向。用格林公式求总面积；若只把左叶反向，求 $\frac12\int_L(x\,\dd y-y\,\dd x)$，解释它此时是否仍表示总面积。
\end{exercise}

\begin{exercise}{475}{全微分与归一化原函数}
验证 $(2x+\sin y)\,\dd x+x\cos y\,\dd y$ 为全微分，求满足 $u(0,0)=0$ 的原函数，并求其在 $y=0$ 上的最小值；通过偏导数复核原函数。
\end{exercise}

\begin{exercise}{476}{半平面原函数与全局障碍}
求 $\frac{y\,\dd x-x\,\dd y}{x^2+y^2}$ 在右半平面 $x>0$ 内满足 $u(1,0)=0$ 的原函数；判断其能否成为去掉原点的整个平面上的单值连续可微原函数，并用闭圆周积分说明理由。
\end{exercise}

\begin{exercise}{477}{参数匹配与初值归一化}
设 $\dd u=(x+a)e^{x+y}\,\dd x+(xe^{x+y}+bx-2y)\,\dd y$。
\begin{enumerate}[label=(\arabic*)]
\item 求 $a,b$，并说明为何必须在整个平面上匹配偏导数。
\item 已知 $u(0,0)=1$，求 $u(x,y)$。
\item 判断等值线 $u=1$ 在原点附近能否写成 $x=x(y)$，给出依据。
\end{enumerate}
\end{exercise}

\begin{exercise}{478}{全微分方程逐项判定}
分别判断下列方程是否为全微分方程，写出偏导数依据，并求其中全微分方程的隐式通解：
\begin{enumerate}[label=(\arabic*)]
\item $(x^2+y)\,\dd x-x\,\dd y=0$；
\item $(e^y-x)\,\dd x+(y+e^x)\,\dd y=0$；
\item $\sin x\cos y\,\dd x+\cos x\sin y\,\dd y=0$；
\item $xy\,\dd x+(x^2+y^2)\,\dd y=0$。
\end{enumerate}
\end{exercise}

\begin{exercise}{479}{全微分方程与初始等值线}
分别求下列方程的隐式通解，以及经过指定点的等值线方程，并检验所求原函数：
\begin{enumerate}[label=(\arabic*)]
\item $(xy^2-2x+y)\,\dd x-(2y-x-x^2y)\,\dd y=0$，指定点 $(0,1)$；
\item $(x^2+e^{-3y})\,\dd x+(y-3xe^{-3y})\,\dd y=0$，指定点 $(1,0)$。
\end{enumerate}
\end{exercise}

\begin{exercise}{480}{环域上的局部条件与全局原函数}
设 $A=\{1<x^2+y^2<4\}$，
$\omega=\frac{-ay\,\dd x+ax\,\dd y}{x^2+y^2}+b(x\,\dd x+y\,\dd y)$，$a,b\in\mathbb R$。
证明偏导数交叉条件在 $A$ 上始终成立，求任意以原点为圆心的正向圆周（半径在 $(1,2)$ 内）上的积分，据此求 $\omega$ 在 $A$ 内存在单值原函数的充要条件。另求 $1\leq x^2+y^2\leq4$ 的完整正向边界积分，明确内外边界方向。
\end{exercise}

\begin{exercise}{481}{只检验圆周能否判断路径无关}
在去掉原点的平面上，设
\[
 \omega_{a,\lambda}=\frac{-ay\,\dd x+ax\,\dd y}{x^2+y^2}
 +\lambda(y^3\,\dd x+x^3\,\dd y).
\]
求所有以原点为圆心的正向圆周上的积分，确定使这些积分全部为零的参数；再计算正向椭圆 $x^2/4+y^2=1$ 上的积分，求使该形式在整个定义域内与路径无关的全部参数。解释为何圆周检验与最终条件不同。
\end{exercise}

\subsection{曲线积分的应用}

\begin{exercise}{482}{变限积分曲线的弧长}
设 $0<b<\pi/2$，$L_b:y=\int_0^x\tan t\,\dd t$（$0\leq x\leq b$）。求弧长 $s(b)$，给出 $b=\pi/4$ 的结果，并讨论 $b\to\pi/2^-$ 时弧长是否有有限极限。
\end{exercise}

\begin{exercise}{483}{空间螺旋弧的长度与位移}
设 $T>0$，$\Gamma:x=e^t\cos t,y=e^t\sin t,z=e^t$（$0\leq t\leq T$）。求弧长并给出 $T=1$ 的结果；计算两端点距离，证明它严格小于弧长，并说明等号为何不可能成立。
\end{exercise}

\begin{exercise}{484}{零密度点与抛物线细丝}
设 $a>0$，$L_a:y=x^2$（$-a\leq x\leq a$），线密度 $\rho=\sqrt y$。求质量与质心横坐标，给出 $a=1$ 的质量；解释密度在原点为零为何不影响质量严格为正。
\end{exercise}

\begin{exercise}{485}{空间细丝的质量与质心}
设 $L:x=t,y=t^2/2,z=t^3/3$（$0\leq t\leq1$），线密度为 $\sqrt{1+2y+3xz}$。求质量及质心横坐标；指出密度与弧长因子的配合如何把根式积分化为多项式积分。
\end{exercise}

\begin{exercise}{486}{细棒密度与质心位置}
细棒位于 $x$ 轴的 $[0,1]$，密度为 $\rho_c(x)=1+c x(1-x)$，$c\geq0$。求质量与质心，给出 $c=1$ 的结果；用关于中点的对称性解释质心为何不随 $c$ 改变。
\end{exercise}

\begin{exercise}{487}{非均匀螺旋弹簧的重心}
设 $L:x=2\cos t,y=2\sin t,z=3t$（$0\leq t\leq2\pi$），线密度为 $x^2+y^2+z^2$。求质量及全部重心坐标，判断重心高度与几何中点高度 $3\pi$ 的大小关系，并用密度随高度的变化说明结果。
\end{exercise}

\begin{exercise}{488}{半圆铁丝的两个轴惯量}
设 $L:x^2+y^2=1,x\geq0$，线密度 $|y|$。求质量及关于 $x,y$ 两轴的转动惯量，检验两惯量之和与质量的关系；说明半圆区域本身的面积积分为何不能代替本题弧长积分。
\end{exercise}

\begin{exercise}{489}{空间大圆的坐标轴惯量}
设 $a>0$，$\Gamma:x^2+y^2+z^2=a^2,x+y+z=0$，线密度为 $\sqrt{x^2+y^2+z^2}$。求质量与关于三个坐标轴的转动惯量，再求关于过原点且垂直于圆平面的轴的惯量，比较两类轴的结果。
\end{exercise}

\begin{exercise}{490}{圆锥螺线上的变力功}
设 $T>0$，质点从原点沿 $\Gamma_T:x=t\cos t,y=t\sin t,z=t$（$0\leq t\leq T$）在力 $\bm F=(y,-x,z)$ 作用下移动。求功 $W(T)$，给出 $T=\pi/2$ 的结果；求使 $W(T)=0$ 的正 $T$ 并讨论功的符号。
\end{exercise}

\begin{exercise}{491}{中心回复力与椭圆路径}
设 $a,b,k>0$，$\bm F(x,y)=-k(x,y)$。质点沿椭圆 $x^2/a^2+y^2/b^2=1$，从 $(-a,0)$ 顺时针移动至 $(0,b)$。求功并讨论其正负；若沿同椭圆反方向到达相同终点，功是否改变？给出原函数依据。
\end{exercise}

\begin{exercise}{492}{径向斥力与平面距离}
设力方向背离原点，大小为 $k/|z|$，$k>0$。质点沿 $x=at,y=bt,z=ct$（$c\ne0$）从 $t=1$ 移动到 $t=2$。求功，分别核查 $c>0,c<0$ 的表达式；再把终点改为 $t=\lambda>1$，求功关于 $\lambda$ 的关系。
\end{exercise}

\begin{exercise}{493}{细丝惯量的平行轴公式及密度条件}
设分段光滑细丝 $L$ 上的连续线密度 $\rho\geq0$，总质量 $M>0$，质心纵坐标为 $\bar y$。
证明关于直线 $y=c$ 的转动惯量满足
$I(c)=I(\bar y)+M(c-\bar y)^2$，由此证明过质心的水平轴唯一使惯量最小。
若允许有符号密度，取单位圆上 $\rho\equiv-1$，计算 $I(c)$，说明最小化结论为何失效。
\end{exercise}

\begin{exercise}{494}{半圆细丝的质心与旋转力功}
设 $R>0$，上半圆细丝 $L:x^2+y^2=R^2,y\geq0$ 的密度为 $1+\lambda y/R$。
求密度非负的全部 $\lambda$、质量、质心及关于原点的极惯量。再设力场为
$\bm F=(-x-\alpha y,-y+\alpha x)$，$\alpha\in\mathbb R$，求质点沿 $L$ 从 $(R,0)$ 逆时针走到 $(-R,0)$ 的功，讨论功的符号，并解释密度参数是否进入功的表达式。
\end{exercise}

\subsubsection{综合提高题}

\begin{exercise}{495}{边界互换对称性}
设 $L$ 为 $y=x^3$ 与 $y=\sqrt[3]x$ 在第一象限围成区域的完整边界。求 $\oint_L(x-y)^{15}\,\dd s$，并证明将 $15$ 换成任意正奇数结论不变；说明此处不需要指定边界方向的原因。
\end{exercise}

\begin{exercise}{496}{象限分解与两种积分}
设 $L:x^2+y^2=1$，$L_1$ 为第一象限弧，$I=\oint_L(x+y+|x|+|y|)\,\dd s$，$J=\int_{L_1}y\,\dd s$。
\begin{enumerate}[label=(\arabic*)]
\item 用对称性确定 $I,J$ 的关系。
\item 计算两者。
\item 把单位圆换为半径 $R>0$ 的圆，求相应两积分并检验关系是否保留。
\end{enumerate}
\end{exercise}

\begin{exercise}{497}{悬链线上的对数绝对值}
设 $a>0$，$L_a:y=(e^x+e^{-x})/2$（$-a\leq x\leq a$）。求 $\int_{L_a}|\ln(2y-e^x)|\,\dd s$，给出 $a=1$ 的结果；写明对数真数为何为正及去掉绝对值的分段点。
\end{exercise}

\begin{exercise}{498}{空间三角形边界与非负弧长}
设 $L$ 为 $O(0,0,0),A(2,2,0),B(2,2,2)$ 构成的三角形完整边界。分边求 $\oint_Lz\sin(xy)\,\dd s$，指出哪条边贡献为零；将边界方向全部反转后，积分是否改变？须根据弧长微分说明。
\end{exercise}

\begin{exercise}{499}{扇形完整边界的径向积分}
设 $R>0,0<\alpha<\pi$，$L$ 为扇形 $0\leq r\leq R,0\leq\theta\leq\alpha$ 的完整边界。求 $\oint_Le^{\sqrt{x^2+y^2}}\,\dd s$，给出 $R=4,\alpha=\pi/4$ 的结果；分别列出两条半径和圆弧的贡献，讨论漏去半径将产生的误差。
\end{exercise}

\begin{exercise}{500}{等值曲线上的符号判断}
设 $f\in C^1(\mathbb R^2)$，$\Gamma$ 为等值线 $f=1$ 上从第二象限点 $M$ 到第四象限点 $N$ 的非零长度光滑弧。分别判断 $\int_\Gamma f\,\dd x$、$\int_\Gamma f\,\dd y$、$\int_\Gamma f\,\dd s$、$\int_\Gamma(f_x\,\dd x+f_y\,\dd y)$ 的符号并证明；说明这些结论是否需要坐标沿路径单调。
\end{exercise}

\begin{exercise}{501}{曲边闭域的直接积分与补线}
设 $L$ 为 $y=x^2,y^2=x$ 围成区域的顺时针边界。计算 $\oint_L\frac1{\sqrt{1+y}}\,\dd x+x^2\,\dd y$，写出两段方向；另用格林公式写出同值二重积分，核查方向符号。
\end{exercise}

\begin{exercise}{502}{曲线族中的功最小化}
设 $a>0$，$L_a:y=a\sin x$ 从 $(0,0)$ 到 $(\pi,0)$。求 $I(a)=\int_{L_a}(1+y^3)\,\dd x+(2x+y)\,\dd y$，找出使 $I(a)$ 最小的曲线并求最小值；证明该点在全部 $a>0$ 中给出全局最小值。
\end{exercise}

\begin{exercise}{503}{曲线约束与空间闭合积分}
设 $\Gamma$ 为 $x^2+y^2+z^2=4,x^2+y^2=2x,z\geq0$ 的完整上半交线，包含闭合点 $(2,0,0)$，从 $x$ 轴正方向看为逆时针。求
$\oint_\Gamma\frac{y^2\,\dd x+z^2\,\dd y+x^2\,\dd z}{2x+z^2}$。
先用两曲面方程化简分母，参数化时在闭合点处分段或取极限；说明该点是否造成积分发散。
\end{exercise}

\begin{exercise}{504}{第二类积分的估值与极限}
\begin{enumerate}[label=(\arabic*)]
\item 设有向分段光滑曲线 $L$ 长 $K$，$P,Q$ 连续，$M=\max_L\sqrt{P^2+Q^2}$。证明 $|\int_LP\,\dd x+Q\,\dd y|\leq KM$，并说明常向量场沿直线段时何时取等号。
\item 设 $L_R:x^2+y^2=R^2$ 取逆时针方向。用上述不等式估计 $I_R=\oint_{L_R}\frac{y\,\dd x-x\,\dd y}{(x^2+xy+y^2)^2}$，求 $R\to+\infty$ 时的极限，并用缩放说明估计中的 $R$ 次数是否准确。
\end{enumerate}
\end{exercise}

\begin{exercise}{505}{绝对值边界与线性变换}
设 $a>0$，$L_a:|x|+|x+y|=a$ 取逆时针方向。求
$\oint_{L_a}\frac{(x-y)\,\dd x+(x+y)\,\dd y}{|x|+|x+y|}$，给出 $a=1$ 的结果；用 $u=x,v=x+y$ 求包围区域面积，并说明变换是否保持方向。
\end{exercise}

\begin{exercise}{506}{隐函数导数与格林公式}
设 $F\in C^1$，$F(z+x,z-y)=0$ 在单位圆盘邻域内唯一确定 $z=z(x,y)\in C^1$，且沿解有 $F_u+F_v\ne0$。推导 $z_x-z_y$ 的值，计算正向单位圆周 $L$ 上
$\oint_L(z-3y)\,\dd x+(z+x^2)\,\dd y$；说明非退化条件在隐式求导中的作用。
\end{exercise}

\begin{exercise}{507}{正函数环流的下界与等号}
设 $a>0$，$L:(x-a)^2+(y-a)^2=a^2$ 取正向，$f$ 在 $[0,2a]$ 上连续且严格为正。证明
$\oint_Lxf(y)\,\dd y-\frac y{f(x)}\,\dd x\geq2\pi a^2$，并求等号成立的充要条件。证明须说明为何只假设 $f$ 连续即可使用所需的偏导数。
\end{exercise}

\begin{exercise}{508}{余弦弧补线与奇偶性}
设 $L:y=\cos x$ 从 $(-\pi/2,0)$ 到 $(\pi/2,0)$。求
$\int_L(e^x-y^3+xy^2)\,\dd x-(x^2y+2x+\sin y)\,\dd y$。
将结果分为端点项和补线后的面积积分项，并说明奇函数项在哪一步消去。
\end{exercise}

\begin{exercise}{509}{最优区域与带奇点环流}
在平面有界单连通闭区域中，令 $D_1$ 为使 $I(D)=\iint_D(4-x^2-y^2)\,\dd x\dd y$ 最大的闭圆盘。
\begin{enumerate}[label=(\arabic*)]
\item 求 $D_1$ 与 $I(D_1)$，证明任何额外包含负值区域或遗漏正值区域都不能改善积分。
\item 计算正向边界上的
$\oint_{\partial D_1}\frac{(xe^{x^2+4y^2}+y)\,\dd x+(4ye^{x^2+4y^2}-x)\,\dd y}{x^2+4y^2}$；区分可写为全微分的部分与角形式部分。
\end{enumerate}
\end{exercise}

\begin{exercise}{510}{路径无关诱导的微分方程}
设 $f\in C^1(\mathbb R)$，$f(0)=0$，且 $[f(x)-e^x]\sin y\,\dd x-f(x)\cos y\,\dd y$ 与路径无关。求 $f$，再求一个原函数并求从 $(0,0)$ 到 $(1,\pi/2)$ 的积分。
\end{exercise}

\begin{exercise}{511}{同圆不同弧的端点一致性}
设 $L$ 为圆周 $x^2+y^2=2y$ 上从 $O(0,0)$ 到 $A(1,1)$ 的任一段有向弧。求 $\int_L(1-2xy-y^2)\,\dd x-(x+y)^2\,\dd y$，证明选择较长或较短的圆弧结果相同，并写出原函数。
\end{exercise}

\begin{exercise}{512}{右半平面中的反求函数}
设 $f$ 在 $x>0$ 可导，$f(1)=0$，$yf(x)\,\dd x+[2xf(x)-x^2]\,\dd y$ 在右半平面内与路径无关。
\begin{enumerate}[label=(\arabic*)]
\item 求 $f(x)$，写出由偏导数条件得到的一阶微分方程。
\item 求从 $(1,0)$ 到 $(2,3)$ 的积分，并验证所求原函数。
\end{enumerate}
\end{exercise}

\begin{exercise}{513}{端点恒等式反求场分量}
设 $Q\in C^1(\mathbb R^2)$，$2xy\,\dd x+Q(x,y)\,\dd y$ 与路径无关，且从原点到 $(t,1)$ 与到 $(1,t)$ 的积分对所有实数 $t$ 相等。求 $Q$ 及满足 $u(0,0)=0$ 的原函数，说明“对所有 $t$”如何确定任意函数项。
\end{exercise}

\begin{exercise}{514}{偏导信息、端点积分与最小值}
设 $f_x=(2x+1)e^{2x-y}$，$f(0,y)=y+1$，$L_t$ 从 $(0,0)$ 到 $(1,t)$。
\begin{enumerate}[label=(\arabic*)]
\item 求 $f(x,y)$，核查边界条件。
\item 求 $I(t)=\int_{L_t}f_x\,\dd x+f_y\,\dd y$。
\item 在全部 $t\in\mathbb R$ 中求 $I(t)$ 的最小值，并证明最小点唯一。
\end{enumerate}
\end{exercise}

\begin{exercise}{515}{三角形有向面积公式}
设三角形顶点按逆时针顺序为 $M_i(x_i,y_i)$（$i=1,2,3$）。用曲线积分证明
$S=\frac12[(x_1y_2-x_2y_1)+(x_2y_3-x_3y_2)+(x_3y_1-x_1y_3)]$。
再说明顶点顺序反转及三个点共线时该表达式的含义，给出不预设方向时的面积公式。
\end{exercise}

\begin{exercise}{516}{全微分参数与等值线方程}
设 $\dd u=(xy^2+a\sin x\cos x\sin3y)\,\dd x-(bx^2y+\cos2x\cos3y)\,\dd y$。
\begin{enumerate}[label=(\arabic*)]
\item 求 $a,b$。
\item 求全部原函数 $u$，并给出 $u(0,0)=0$ 的归一化结果。
\item 求将该微分形式置零所得方程的隐式通解，并求经过 $(0,0)$ 的等值线。
\end{enumerate}
\end{exercise}

\begin{exercise}{517}{不同曲线的弧长比较}
设 $s_1$ 为 $y=\sin x$（$0\leq x\leq2\pi$）的弧长，$s_2$ 为椭圆 $x^2+2y^2=2$ 的周长。写出两者的定积分表达式，通过变量代换确定其关系；不要求计算非初等积分的具体值，并说明比较为何不能仅靠图形外观。
\end{exercise}

\begin{exercise}{518}{微分方程曲线的均匀重心}
设 $y(x)$ 满足 $2xy'-4y=2\ln x-1$、$y(1)=1/4$。求 $y(x)$，化简弧长因子，求均匀曲线弧 $1\leq x\leq e$ 的质量（单位线密度）及重心的两个坐标；用 $1<\bar x<e$ 检验结果。
\end{exercise}

\begin{exercise}{519}{旋转径向力与角度功}
质点在 $(x,y)\ne(0,0)$ 受到力，大小为到原点距离的倒数的两倍，方向为 $\overrightarrow{OM}$ 顺时针旋转 $90^\circ$ 的方向。写出力的坐标分量，求从 $A(0,\pi/2)$ 沿 $y=(\pi/2)\cos x$ 到 $B(\pi/2,0)$ 的功；再判断换成第一象限内同端点的任意路径，功是否改变，并说明定义域条件。
\end{exercise}

\clearpage
\section{第12章\quad 曲面积分}
\subsection{第一类曲面积分}
\begin{exercise}{520}{上半球的对称分解}
设 $a>0$，$S$ 为球面 $x^2+y^2+z^2=a^2$ 的上半球面，$S_1$ 为其在第一卦限的部分。分别判断 $\iint_Sx\,\dd S,\iint_Sy\,\dd S,\iint_Sz\,\dd S,\iint_Sxyz\,\dd S$ 中哪些等于相应 $S_1$ 上积分的四倍；计算这四个全上半球积分，并说明每一次使用对称性的依据。
\end{exercise}
\begin{exercise}{521}{切平面与奇函数消去}
先求抛物面 $z=x^2+y^2$ 在 $(0,1,1)$ 处的切平面，令 $\Sigma$ 为该平面被柱面 $x^2+(y-1)^2=1$ 截出的部分。求
\[I(\lambda)=\iint_{\Sigma}(x^3yz^2+\lambda z)\,\dd S\quad(\lambda\in\mathbb R),\]
并给出 $I(1)$。说明能否仅因投影圆盘关于 $y=1$ 对称，就将含 $z$ 的积分判为零。
\end{exercise}
\begin{exercise}{522}{曲面上的恒等变换}
设 $\Sigma$ 为平面 $\frac x2+\frac y3+\frac z4=1$ 在第一卦限中的部分。先证明 $2x+\frac43y+z$ 在 $\Sigma$ 上为常数，再计算 $\iint_{\Sigma}(2x+\frac43y+z)\,\dd S$；利用这一结果及对称或重心性质，分别求 $\iint_{\Sigma}x\,\dd S$、$\iint_{\Sigma}y\,\dd S$。
\end{exercise}
\begin{exercise}{523}{平面三角形的二阶矩}
设 $t>0$，$\Sigma_t=\{(x,y,z):x+y+z=t,\ x,y,z\ge0\}$。求 $\iint_{\Sigma_t}y^2\,\dd S$，并给出 $t=1$ 时的值；利用变量轮换与 $(x+y+z)^2=t^2$，求 $\iint_{\Sigma_t}xy\,\dd S$。
\end{exercise}
\begin{exercise}{524}{球面约束与函数信息}
设 $R>0$，$f$ 在 $\mathbb R$ 上连续，$\Sigma_R$ 为上半球面 $z=\sqrt{R^2-x^2-y^2}$。求
\[I(R)=\iint_{\Sigma_R}\left[xy\sqrt{x^4+y^4}+zf(x^2+y^2+z^2+1)\right]\,\dd S.\]
若对所有 $R>0$ 都有 $I(R)=0$，确定 $f$ 在 $[1,\infty)$ 上的取值，并说明能否由此确定 $f$ 在 $(-\infty,1)$ 上的值。
\end{exercise}
\begin{exercise}{525}{球冠与截取参数}
设 $k>0$，$\Sigma_k$ 为单位球面上满足 $z\ge k\sqrt{x^2+y^2}$ 的顶部。求 $I(k)=\iint_{\Sigma_k}z\,\dd S$，给出原截面参数 $k=1$ 时的值；判断 $I(k)$ 的单调性，并用曲面包含关系解释结果。
\end{exercise}
\begin{exercise}{526}{八面体的分片与对称性}
设 $\Sigma:|x|+|y|+|z|=1$。计算曲面面积 $A$ 与 $\iint_{\Sigma}(x+|y|)\,\dd S$，进而求均匀曲面上 $|y|$ 的平均值；说明为何不能把含 $|y|$ 的积分也按奇函数消去。
\end{exercise}
\begin{exercise}{527}{球面的二阶矩}
设 $R>0$，$\Sigma:x^2+y^2+z^2=R^2$。利用变量对称性求 $\iint_{\Sigma}z^2\,\dd S$，再求 $\iint_{\Sigma}(a x^2+b y^2+c z^2)\,\dd S$，并给出该积分为零的参数条件，其中 $a,b,c\in\mathbb R$。
\end{exercise}
\begin{exercise}{528}{偏心投影域与奇偶性}
设 $a>0$，$\Sigma$ 为锥面 $z=\sqrt{x^2+y^2}$ 在柱体 $x^2+y^2\le2ax$ 内的部分。说明 $xy,yz,zx$ 三项中哪些积分由奇偶性为零，计算 $\iint_{\Sigma}(xy+yz+zx)\,\dd S$；再说明把柱体改为 $x^2+y^2\le-2ax$ 时积分怎样改变。
\end{exercise}
\begin{exercise}{529}{锥面积分与体积比较}
设 $\Sigma$ 为锥面 $z=\sqrt{x^2+y^2}$ 在柱体 $x^2+y^2\le2x$ 内的部分。求 $\iint_{\Sigma}z\,\dd S$；再设 $\Omega$ 为该曲面与其在 $xOy$ 平面上的投影之间的立体，求 $\Omega$ 的体积，说明二者的比例为何为常数。
\end{exercise}
\begin{exercise}{530}{半球的线性函数平均值}
设 $a>0$，$\Sigma$ 为上半球面 $z=\sqrt{a^2-x^2-y^2}$。求 $\iint_{\Sigma}(x+y+z)\,\dd S$ 及被积函数在该曲面上的平均值；若将上半球改为下半球，分别说明积分值与平均值怎样改变。
\end{exercise}
\begin{exercise}{531}{锥体边界的分片面积积分}
设 $\Sigma$ 为立体 $\sqrt{x^2+y^2}\le z\le1$ 的整个边界。分别求锥面与顶圆盘上的 $\iint(x^2+y^2)\,\dd S$ 并求总和；再将上界 $1$ 改为 $h>0$，用缩放关系求总积分，说明这类积分是否依赖法向选择。
\end{exercise}
\begin{exercise}{532}{平面曲面的三阶矩}
设 $\Sigma$ 为平面 $x+y+z=1$ 在第一卦限中的部分。计算 $\iint_{\Sigma}xyz\,\dd S$；再将平面常数 $1$ 改为 $t>0$，不重新积分，利用缩放关系求相应积分并说明缩放次数。
\end{exercise}
\begin{exercise}{533}{投影积分与截面换元}
设 $\Sigma$ 为平面 $x+y+z=1$ 在第一卦限中的部分。令 $c>0$，求
\[J(c)=\iint_{\Sigma}\frac{\dd S}{(c+x+y)^2},\]
并给出 $J(1)$；将二重积分化为以 $u=x+y$ 为变量的一元积分，进而判断 $J(c)$ 随 $c$ 的单调性。
\end{exercise}
\begin{exercise}{534}{圆柱面的参数与缩放}
设 $R,H>0$，$\Sigma$ 为圆柱面 $x^2+y^2=R^2$ 夹在 $z=0,z=H$ 之间的侧面。求 $\iint_{\Sigma}\frac1{x^2+y^2}\,\dd S$；进一步把被积函数改为 $(x^2+y^2)^{-p/2}$，求积分关于实参数 $p$ 的表达式，并讨论同时将 $R,H$ 放大 $k>0$ 倍时的变化。
\end{exercise}
\begin{exercise}{535}{锥面的参数化与尺度}
设 $\Sigma_h$ 是锥面 $z^2=3(x^2+y^2)$ 被 $z=0,z=h$ 截出的部分，$h>0$。写出其参数表示及面积元素，求 $I(h)=\iint_{\Sigma_h}(x^2+y^2)\,\dd S$；给出 $h=3$ 时的值，并说明高度扩大为原来两倍时积分变为原来的几倍。
\end{exercise}
\begin{exercise}{536}{抛物面面积元素与比较}
设 $\Sigma$ 为抛物面 $z=2-x^2-y^2$ 在 $xOy$ 平面上方的部分。将 $\iint_{\Sigma}(x^2+y^2)\,\dd S$ 化为一元定积分并计算；再求其在 $xOy$ 平面上的投影区域 $D$ 上的二重积分 $\iint_D(x^2+y^2)\,\dd x\,\dd y$，比较两者大小并解释原因。
\end{exercise}
\begin{exercise}{537}{曲面质量与平均高度}
设 $\Sigma$ 为抛物面 $z=2-x^2-y^2$ 在 $xOy$ 平面上方的部分。求曲面面积 $A$ 与 $I=\iint_{\Sigma}3z\,\dd S$，进而求该均匀曲面的平均高度 $\bar z=I/(3A)$；验证结果应满足的几何范围。
\end{exercise}
\begin{exercise}{538}{平面投影与积分线性性}
设 $\Sigma$ 为平面 $2x+2y+z=6$ 在第一卦限中的部分，$\lambda\in\mathbb R$。求
\[I(\lambda)=\iint_{\Sigma}(2xy-2x^2-\lambda x+z)\,\dd S,\]
给出 $I(1)$，并求使 $I(\lambda)=0$ 的参数。须写明投影区域和面积元素。
\end{exercise}
\begin{exercise}{539}{球冠的对称性与参数}
设 $a>0,0<h<a$，$\Sigma_h$ 为球面 $x^2+y^2+z^2=a^2$ 上满足 $z\ge h$ 的部分。说明哪些被积项可由对称性消去，求 $I(h)=\iint_{\Sigma_h}(x+y+z)\,\dd S$；再求满足 $I(h)=I(0)/2$ 的 $h$，其中 $I(0)$ 表示上半球面上的同一积分。
\end{exercise}
\begin{exercise}{540}{面积元素与投影元素辨析}
设 $S$ 为单位上半球面，取上侧。分别计算 $\iint_Sz\,\dd S$ 与 $\iint_Sz\,\dd x\,\dd y$；从两种面积元素的关系解释它们为何不同。反向后两个积分各怎样变化？在一般图像曲面 $z=g(x,y)$ 上写出二者相等所需的积分条件。
\end{exercise}
\begin{exercise}{541}{锥面奇异密度与重心}
设 $h>0$，锥面 $S_h:z=\sqrt{x^2+y^2},0<z\le h$ 的面密度为 $\rho=z^{-p}$，$p\in\mathbb R$。把质量定义为从 $z=\varepsilon$ 截起后令 $\varepsilon\to0^+$ 的极限。求质量有限的参数范围；在此范围内求总质量及重心高度，并讨论 $p$ 趋近临界值时重心高度的极限。
\end{exercise}
\subsection{第二类曲面积分}
\begin{exercise}{542}{球面内侧与投影定向}
设 $\Sigma$ 为单位球面在第一卦限的部分，取内侧。将 $\iint_{\Sigma}xyz\,\dd x\,\dd y$ 化为平面二重积分并计算；再说明取外侧时的结果。若将积分面扩为整个上半球面的内侧，该积分又是多少？分别说明定向与对称性的作用。
\end{exercise}
\begin{exercise}{543}{平面定向与反常积分}
设 $\Sigma$ 为单位圆盘 $z=0,\ x^2+y^2\le1$ 的下侧。先求 $\iint_{\Sigma}(x^2+y^2)\,\dd x\,\dd y$，再讨论
\[I(p)=\lim_{\varepsilon\to0^+}\iint_{\Sigma\cap\{x^2+y^2\ge\varepsilon^2\}}(x^2+y^2)^{p/2}\,\dd x\,\dd y\]
何时存在有限值，并在存在时求出 $I(p)$，其中 $p\in\mathbb R$。
\end{exercise}
\begin{exercise}{544}{球面投影与齐次缩放}
设 $R>0$，$\Sigma_-$ 为球面 $x^2+y^2+z^2=R^2$ 的下半部分，取下侧。计算 $\iint_{\Sigma_-}x^2y^2\,\dd x\,\dd y$；再求同一被积式在上半球面上侧及整球面外侧的积分，说明两者的符号关系。用缩放关系核验结果关于 $R$ 的次数。
\end{exercise}
\begin{exercise}{545}{偏心圆域上的球面通量}
设 $a>0$，$\Sigma$ 为上半球面 $z=\sqrt{a^2-x^2-y^2}$ 在柱体 $x^2+y^2\le ax$ 内的部分，取上侧。求 $\iint_{\Sigma}z^2\,\dd x\,\dd y$，并分别计算投影区域的面积与其关于原点的极惯性矩 $\iint_D(x^2+y^2)\,\dd x\,\dd y$，核验它们与所求积分的关系。
\end{exercise}
\begin{exercise}{546}{闭曲面与可积奇点}
设 $\Omega=\{(x,y,z):1\le y\le2,\ x^2+z^2\le y^2\}$，$\Sigma=\partial\Omega$ 取外侧。把
\[\iint_{\Sigma}\frac{2y}{\sqrt{x^2+z^2}}\,\dd z\,\dd x\]
解释为先从积分面中删去 $x^2+z^2<\varepsilon^2$ 的部分、再令 $\varepsilon\to0^+$ 的极限。分别求侧面、上底面和下底面的贡献，证明该极限存在并计算总值；说明不能未经处理就直接在整个 $\Omega$ 上套用高斯公式的原因。
\end{exercise}
\begin{exercise}{547}{柱面定向与切向分量}
设 $\Sigma$ 为圆柱面 $x^2+y^2=1$ 在第一卦限且 $0\le z\le3$ 的部分，取法向量的 $x$ 分量非负的一侧。计算 $\iint_{\Sigma}x\,\dd y\,\dd z+y\,\dd z\,\dd x+z\,\dd x\,\dd y$；再说明把末项系数 $z$ 换成任意连续函数 $R(x,y,z)$，结果是否改变，并从法向量解释原因。
\end{exercise}
\begin{exercise}{548}{两类曲面积分的转换}
设 $\Sigma$ 为平面 $3x+2y+2\sqrt3z=6$ 在第一卦限的部分，取上侧，$P,Q,R$ 在 $\Sigma$ 上连续。求单位法向量，将 $\iint_{\Sigma}P\,\dd y\,\dd z+Q\,\dd z\,\dd x+R\,\dd x\,\dd y$ 化为第一类曲面积分；应用所得公式计算 $P=3,Q=2,R=2\sqrt3$ 时的通量，并解释其与曲面面积的关系。
\end{exercise}
\begin{exercise}{549}{变法向量与转换公式}
设 $\Sigma$ 为抛物面 $z=8-x^2-y^2$ 在 $xOy$ 平面上方的部分，取上侧。将连续向量场 $(P,Q,R)$ 通过 $\Sigma$ 的通量写成第一类曲面积分及投影平面上的二重积分两种形式；再取 $(P,Q,R)=(2x,2y,1)$，计算通量并说明该向量场与曲面法向的关系。
\end{exercise}
\begin{exercise}{550}{法向正交与任意函数消去}
设 $\Sigma=\{x-y+z=1:x\ge0,\ y\le0,\ z\ge0\}$，取上侧，$f$ 在 $\Sigma$ 上连续。计算
\[\iint_{\Sigma}(f+x)\,\dd y\,\dd z+(2f+y)\,\dd z\,\dd x+(f+z)\,\dd x\,\dd y.\]
再将三个系数中的 $f,2f,f$ 分别改为 $af,bf,cf$，求使积分对任意连续函数 $f$ 都不变的实参数条件，并给出必要性证明。
\end{exercise}
\begin{exercise}{551}{闭柱面分片与奇点条件}
设 $R>0$，$\Sigma$ 为圆柱体 $x^2+y^2\le R^2,\ |z|\le R$ 的整个边界，取外侧。分片计算
\[\iint_{\Sigma}\frac{x\,\dd y\,\dd z+z^2\,\dd x\,\dd y}{x^2+y^2+z^2}.\]
说明两个底面上的积分为何抵消，并指出直接在整个柱体上使用高斯公式时缺少哪项条件。
\end{exercise}
\begin{exercise}{552}{柱面通量与无效分量}
设 $\Sigma$ 为圆柱面 $x^2+y^2=4$ 被 $z=0$ 与 $x+z=2$ 截出的侧面，取外侧。对实参数 $c$ 求
\[I(c)=\iint_{\Sigma}(-y)\,\dd z\,\dd x+(z+c)\,\dd x\,\dd y,\]
并给出 $c=1$ 时的值；解释结果与 $c$ 是否有关，再利用补上两个底面后的高斯公式核验。
\end{exercise}
\begin{exercise}{553}{抛物面定向与补面核验}
设 $\Sigma$ 为 $z=\frac12(x^2+y^2)$ 介于 $z=0,z=2$ 之间的曲面，取下侧。直接计算 $\iint_{\Sigma}(z^2+x)\,\dd y\,\dd z-z\,\dd x\,\dd y$；再补上 $z=2$ 处圆盘，用高斯公式核验结果，明确两块曲面的定向。
\end{exercise}
\begin{exercise}{554}{半球面的法向与两种算法}
设 $R>0$，$\Sigma$ 为上半球面 $z=\sqrt{R^2-x^2-y^2}$ 的上侧。求向量场 $(xz,yz,z^2)$ 通过 $\Sigma$ 的通量；分别利用球面的单位法向量和补面高斯公式计算，并说明补面上的通量为何为零。
\end{exercise}
\begin{exercise}{555}{柱面参数化与高度参数}
设 $h>0$，$\Sigma_h$ 为柱面 $x^2+y^2=1,\ 0\le z\le h$ 的外侧。求
\[I(h)=\iint_{\Sigma_h}x(y-z)\,\dd y\,\dd z+(x-y)\,\dd x\,\dd y,\]
给出 $h=2$ 时的值，并说明柱面上哪一个坐标面积元素恒为零。
\end{exercise}
\begin{exercise}{556}{抛物面的高度与方向}
设 $h>0$，$\Sigma_h$ 为抛物面 $z=x^2+y^2$ 介于 $z=0,z=h$ 之间的部分，取下侧。求 $I(h)=\iint_{\Sigma_h}x^2\,\dd y\,\dd z+z\,\dd x\,\dd y$，给出 $h=1$ 时的结果；说明其中一项为何因对称性消失，并判断反向后结果如何变化。
\end{exercise}
\begin{exercise}{557}{四分之一柱面的通量}
设 $a,h>0$，$\Sigma$ 为柱面 $x^2+y^2=a^2$ 在第一卦限且 $0\le z\le h$ 的部分，取外侧。计算 $\iint_{\Sigma}x\,\dd y\,\dd z+z\,\dd x\,\dd y$；再将首项系数 $x$ 改为 $y$，求新的积分，并解释两者为何不能仅凭 $x,y$ 对称就判为相等。
\end{exercise}
\begin{exercise}{558}{平面下侧与参数零点}
设 $\Sigma=\{x+y-z=1:x\ge0,\ y\ge0,\ z\le0\}$，取下侧。求 $I(c)=\iint_{\Sigma}c\,\dd y\,\dd z+z\,\dd x\,\dd y$，其中 $c\in\mathbb R$；给出 $I(1)$ 并求使 $I(c)=0$ 的参数，写明投影区域和定向符号。
\end{exercise}
\begin{exercise}{559}{无散场与参数扰动}
设 $h>0$，$\Sigma$ 为圆锥体 $\{0\le y\le h,\ x^2+z^2\le y^2\}$ 的边界外侧。先计算向量场 $(y-z,z-x,x-y)$ 通过 $\Sigma$ 的通量，再求向量场 $(y-z,z-x,x-y)+\lambda(x,y,z)$ 的通量，其中 $\lambda\in\mathbb R$；分别说明散度计算及体积计算的依据。
\end{exercise}
\begin{exercise}{560}{球面对称性与线性向量场}
设 $R>0$，$\Sigma$ 为球面 $x^2+y^2+z^2=R^2$ 的外侧。先利用球面法向量直接计算 $\iint_{\Sigma}x\,\dd y\,\dd z+y\,\dd z\,\dd x+z\,\dd x\,\dd y$，再推广到向量场 $(\lambda x,\mu y,\nu z)$，并求通量为零的参数条件。
\end{exercise}
\begin{exercise}{561}{无散场与四面体分面通量}
设 $a>0$，$\Omega=\{x,y,z\ge0:x+y+z\le a\}$，$S=\partial\Omega$ 取外侧。求向量场 $(yz,zx,xy)$ 的总通量，再分别计算三个坐标面上的通量，由此求斜面 $x+y+z=a$ 上的通量；说明总通量为零并不要求各面通量为零。
\end{exercise}
\begin{exercise}{562}{径向场与内部奇点}
设 $R>0$，$\Sigma_R$ 为球面 $x^2+y^2+z^2=R^2$ 的外侧，$p\in\mathbb R$。计算
\[I_p(R)=\iint_{\Sigma_R}\frac{x\,\dd y\,\dd z+y\,\dd z\,\dd x+z\,\dd x\,\dd y}{(x^2+y^2+z^2)^{p/2}},\]
求使通量与 $R$ 无关的 $p$，并给出 $p=3$ 时的值；解释此时在原点外散度为零却不能推出球面通量为零的原因。
\end{exercise}
\begin{exercise}{563}{下半球的奇偶消去}
设 $R>0$，$\Sigma_-$ 为下半球面 $z=-\sqrt{R^2-x^2-y^2}$，取下侧。计算 $\iint_{\Sigma_-}x^2\,\dd y\,\dd z+y^2\,\dd z\,\dd x+z^2\,\dd x\,\dd y$；再求整个球面外侧的同一通量，并用上、下半球分解核验结论。
\end{exercise}
\begin{exercise}{564}{柱面通量的算法选择}
设 $\Sigma$ 为柱面 $x^2+y^2=1,\ 0\le z\le2$ 的外侧。用补面高斯公式计算 $\iint_{\Sigma}x(y-z)\,\dd y\,\dd z+(x-y)\,\dd x\,\dd y$；分别求上下底面通量，并与直接参数化所得结果核对。
\end{exercise}
\begin{exercise}{565}{锥面通量与对称项}
设 $\Sigma$ 为锥面 $z=1-\sqrt{x^2+y^2}$（$0\le z\le1$）的上侧。求向量场 $(x^2-yz,y^2-xz,2z)$ 通过 $\Sigma$ 的通量，指出直接投影计算中哪些项因对称性消去；再用补面法计算并说明底面通量是否为零。
\end{exercise}
\begin{exercise}{566}{闭曲面通量与尺度}
设 $a>0$，$\Sigma_a$ 为立方体 $[0,a]^3$ 的边界，取外侧。利用高斯公式计算
\[\iint_{\Sigma_a}y(x-z)\,\dd y\,\dd z+x^2\,\dd z\,\dd x+(y^2+xz)\,\dd x\,\dd y.\]
分别计算面 $x=a$ 与 $x=0$ 对通量的贡献，核对二者之和与对应散度项的体积分一致，并说明该通量随 $a$ 的尺度变化规律。
\end{exercise}
\begin{exercise}{567}{定向面积与切向场}
设 $S$ 为锥面 $z=\sqrt{x^2+y^2},1\le z\le2$ 的下侧。证明向量场 $F=(x,y,z)$ 与 $S$ 相切，因而通量为零；再计算 $\iint_S|F|\,\dd S$。据此判断能否用向量场的模直接代替法向分量计算通量。
\end{exercise}
\begin{exercise}{568}{抛物面上的参数通量}
设 $S$ 为抛物面 $z=1-x^2-y^2,z\ge0$ 的上侧。对 $F_{a,b}=(ax,ay,bz)$ 求通量 $\Phi(a,b)$，确定使通量为零的全部实参数对；再计算 $\iint_SF_{a,b}\cdot n\,\dd S$ 中法向分量逐点为零的参数条件，比较两种条件。
\end{exercise}
\subsection{高斯公式}
\begin{exercise}{569}{球锥区域与通量分解}
设 $R>0$，$\Omega$ 为锥面 $z=\sqrt{x^2+y^2}$ 与上半球面 $z=\sqrt{R^2-x^2-y^2}$ 围成的区域。令 $\Sigma=\partial\Omega$ 取外侧，求
\[\iint_{\Sigma}x\,\dd y\,\dd z+y\,\dd z\,\dd x+z\,\dd x\,\dd y.\]
再直接判断锥面部分的通量，进而求球面部分的通量；说明高斯公式的光滑性条件在此处是否满足。
\end{exercise}
\begin{exercise}{570}{补面与定向}
设 $\Sigma$ 为锥面 $z=\sqrt{x^2+y^2}$（$0\le z\le1$）的下侧，$c\in\mathbb R$。通过补上 $z=1$ 处的圆盘求
\[I(c)=\iint_{\Sigma}x\,\dd y\,\dd z+2y\,\dd z\,\dd x+3(z-c)\,\dd x\,\dd y.\]
明确补面法向的方向，给出 $c=1$ 时的积分值，并求使锥面通量为零的 $c$。
\end{exercise}
\begin{exercise}{571}{半球补面与缩放}
设 $R>0$，$\Sigma_R$ 为上半球面 $z=\sqrt{R^2-x^2-y^2}$ 的上侧。求 $\iint_{\Sigma_R}xy\,\dd y\,\dd z+x\,\dd z\,\dd x+x^2\,\dd x\,\dd y$，给出 $R=2$ 时的值；指出闭半球的总通量及补面通量各是多少。
\end{exercise}
\begin{exercise}{572}{散度恒等式与区域对称}
设 $a>0$，$\Omega=\{0\le z\le a^2-x^2-y^2\}$，$S=\partial\Omega$ 取外侧，体积为 $V$。证明
\[\iint_Sx^2yz^2\,\dd y\,\dd z-xy^2z^2\,\dd z\,\dd x+z(1+xyz)\,\dd x\,\dd y=V,\]
并求 $V$。若将区域改为 $\Omega\cap\{x\ge0,y\ge0\}$，上述通量是否仍等于新区域体积？用散度积分说明。
\end{exercise}
\begin{exercise}{573}{径向散度与曲面分解}
设 $a>0$，$S$ 为上半球体 $x^2+y^2+z^2\le a^2,\ z\ge0$ 的整个边界外侧。求向量场 $(xz^2,x^2y-z^3,2xy+y^2z)$ 的总通量，分别说明平底面与弯曲球面的贡献；用球坐标写出散度体积分的计算过程。
\end{exercise}
\begin{exercise}{574}{全边界与侧面的区别}
设 $S$ 为柱体 $x^2+y^2\le1,\ 0\le z\le3$ 的全边界外侧。求 $\iint_Sx(y-z)\,\dd y\,\dd z+(x-y)\,\dd x\,\dd y$，再分别求上下底面贡献，并据此恢复侧面通量。说明仅改变积分面的范围可能怎样改变结果。
\end{exercise}
\begin{exercise}{575}{立方体平移与奇函数}
设 $a>0$，$S$ 为 $[0,a]^3$ 的边界外侧，求向量场 $(x^2,y^2,z^2)$ 的通量；再将立方体改为 $[-a,a]^3$，求同一向量场的通量，并解释差异来自散度的哪种性质。
\end{exercise}
\begin{exercise}{576}{通量与高阶球面矩}
设 $a>0$，$S$ 为半径 $a$、球心在原点的球面外侧。利用高斯公式求向量场 $(x^3,y^3,z^3)$ 的通量，再由单位法向量推得 $\iint_S(x^4+y^4+z^4)\,\dd S$，进而求 $\iint_Sz^4\,\dd S$。
\end{exercise}
\begin{exercise}{577}{径向通量与体积恢复}
设 $S$ 为柱体 $x^2+y^2\le9,\ 0\le z\le3$ 的整个边界外侧。求向量场 $(x,y,z)$ 的通量，并直接分侧面和两底面核验；若某个分片光滑闭曲面上该场的外通量为 $M$，求它围成区域的体积并说明适用条件。
\end{exercise}
\begin{exercise}{578}{单位立方体的参数通量}
设 $S=\partial[0,1]^3$ 取外侧。求 $I(c)=\iint_Scxz\,\dd y\,\dd z-y^2\,\dd z\,\dd x+yz\,\dd x\,\dd y$，给出 $I(4)$；求使总通量为零的实数 $c$，并解释总通量为零能否推出向量场处处为零。
\end{exercise}
\begin{exercise}{579}{斜顶柱体与散度对称}
设 $\Omega=\{x^2+y^2\le1,\ 0\le z\le1-x\}$，$S=\partial\Omega$ 取外侧。求向量场 $(2xz,xz\cos y,3yz\sin x)$ 的总通量，写明消失项对应的区域对称性；再计算 $\Omega$ 的体积及 $z$ 坐标的体积平均值。
\end{exercise}
\begin{exercise}{580}{旋转曲面与指定法向}
将 $x=0,z=\sqrt{y-1}$（$1\le y\le3$）绕 $y$ 轴旋转得到 $\Sigma$，取法向量与 $y$ 轴正向成钝角的一侧。先写出曲面方程，再求向量场 $(x(8y+1),2(1-y^2),-4yz)$ 的通量；用补上 $y=3$ 圆盘的方法核验定向与结果。
\end{exercise}
\begin{exercise}{581}{非零补面通量}
设 $\Sigma$ 为球面 $x^2+y^2+z^2=4,z\ge0$ 的外侧。求 $\iint_{\Sigma}yz\,\dd z\,\dd x+2\,\dd x\,\dd y$，并分别列出下底面通量与闭半球总通量。若把常数 $2$ 换为 $c$，求所得结果关于 $c$ 的表达式。
\end{exercise}
\begin{exercise}{582}{方向余弦与通量表达}
设 $\Sigma$ 为半径 $R>0$ 的原点球面外侧，$(\cos\alpha,\cos\beta,\cos\gamma)$ 为单位法向量。把 $\iint_{\Sigma}(x^3\cos\alpha+y^3\cos\beta+z^3\cos\gamma)\,\dd S$ 改写为第二类曲面积分并计算；再求其中单独一项 $\iint_{\Sigma}z^3\cos\gamma\,\dd S$。
\end{exercise}
\begin{exercise}{583}{锥面方向余弦与高度}
设 $h>0$，$\Sigma$ 为锥面 $x^2+y^2=z^2,\ 0\le z\le h$ 的下侧。求方向余弦，计算 $\iint_{\Sigma}(x^2\cos\alpha+y^2\cos\beta+z^2\cos\gamma)\,\dd S$；说明前两项分别为何为零，并用缩放关系核验高度次数。
\end{exercise}
\begin{exercise}{584}{抛物面补面与常数项}
设 $\Sigma$ 为 $z=1-x^2-y^2,z\ge0$ 的上侧。对实参数 $c$，求向量场 $(2x^3,2y^3,3(z^2-c))$ 通过 $\Sigma$ 的通量；给出 $c=1$ 时的值，并指出为何散度虽与 $c$ 无关，开曲面通量仍可能依赖 $c$。
\end{exercise}
\begin{exercise}{585}{椭圆投影与换元}
设 $\Sigma$ 为 $z=1-x^2-\frac{y^2}{4},\ 0\le z\le1$ 的上侧。计算向量场 $(xz,2yz,3xy)$ 的通量；将椭圆投影化为单位圆盘，明确雅可比因子，并解释底面通量是否为零。
\end{exercise}
\begin{exercise}{586}{平移多项式与区域奇偶}
设 $\Sigma$ 为抛物面 $z=x^2+y^2,z\le1$ 的上侧。求向量场 $((x-1)^3,(y-1)^3,z-1)$ 的通量，写出投影后的偶函数部分；再求只把第三分量改为 $z$ 时通量的增量。
\end{exercise}
\begin{exercise}{587}{四面体的散度矩}
设 $\Omega=\{x,y,z\ge0:2x+y+2z\le2\}$，$S=\partial\Omega$ 取外侧。求向量场 $(x^2+1,-2y,3z)$ 的通量，并计算 $\Omega$ 的体积及 $\iiint_{\Omega}x\,\dd V$，说明它们如何组成总通量。
\end{exercise}
\begin{exercise}{588}{开曲面通量与非零底面}
设 $\Sigma$ 为抛物面 $z=1-x^2-y^2,z\ge0$ 的上侧。计算向量场 $(y^2-x,z^2-y,x^2-z)$ 通过 $\Sigma$ 的通量；另求闭合抛物面体的总通量和底面通量，验证补面时的减号。
\end{exercise}
\begin{exercise}{589}{抛物面体的高阶矩通量}
设 $\Omega=\{x^2+y^2\le z\le4\}$，$S=\partial\Omega$ 取外侧。求向量场 $(x^3,y^3,z^3)$ 的通量，将散度积分化为柱坐标积分；再分别求顶面与抛物面部分的通量，核验二者之和。
\end{exercise}
\begin{exercise}{590}{四面体的坐标矩}
设 $S$ 为四面体 $x,y,z\ge0,\ x+y+z\le1$ 的边界外侧。计算 $\iint_Sxy\,\dd y\,\dd z+yz\,\dd z\,\dd x+xz\,\dd x\,\dd y$；再将右端常数 $1$ 改为 $t>0$，用缩放关系求相应通量。
\end{exercise}
\begin{exercise}{591}{球锥区域的柱坐标积分}
设 $S$ 为 $\sqrt{x^2+y^2}\le z\le\sqrt{2-x^2-y^2}$ 围成区域的边界外侧。求向量场 $(2xz,yz,-x^2)$ 的通量，写明柱坐标积分限；再单独求锥面部分的通量，并由总通量求球面部分的贡献。
\end{exercise}
\begin{exercise}{592}{锥顶的反常通量}
设 $\Sigma$ 为圆锥体 $0\le z\le1,\ x^2+y^2\le z^2$ 的侧面外侧。将 $\iint_{\Sigma}x^2\,\dd y\,\dd z+\dd z\,\dd x+\frac{e^z}{\sqrt{x^2+y^2}}\,\dd x\,\dd y$ 定义为删去 $z<\varepsilon$ 后的极限，证明极限存在并计算；说明面积元素中的因子如何抵消锥顶附近的奇性。
\end{exercise}
\begin{exercise}{593}{球壳锥域上的任意函数项}
设 $\Omega=\{1\le x^2+y^2+z^2\le4,\ x\ge\sqrt{y^2+z^2}\}$，$S=\partial\Omega$ 取外侧，$f\in C^1(\mathbb R)$。求向量场 $(x^3,y^3-zf(x),z^3+yf(x))$ 的总通量。证明结果与 $f$ 无关，并从球面与锥面的法向说明附加场 $(0,-zf(x),yf(x))$ 在每块边界上都相切。
\end{exercise}
\begin{exercise}{594}{高斯公式的内部光滑性}
在 $r=\sqrt{x^2+y^2+z^2}>0$ 上设 $F=(x,y,z)/r^3$。求其散度，分别计算它通过球壳 $1\le r\le2$ 全边界及外球面 $r=2$ 的外通量；写明内球面作为球壳边界时的方向，解释为何不能把外球面单独视为零通量面。
\end{exercise}
\begin{exercise}{595}{任意区域零通量的参数识别}
设 $F=(ax^2+yz,by^2+xz,cz^2+xy)$。若任意有界分片光滑区域边界的外通量均为零，求 $a,b,c$，并证明必要性。若只要求每个以原点为球心的球面外通量为零，参数条件如何改变？说明这种检验的信息差别。
\end{exercise}
\subsection{斯托克斯公式}
\begin{exercise}{596}{旋度与平面法向的正交}
设 $L$ 为第一卦限内平面三角形 $x+y+z=1$ 的边界，从 $z$ 轴正向看为逆时针。计算 $\oint_L(y^2+z^2)\,\dd x+(z^2+x^2)\,\dd y+(x^2+y^2)\,\dd z$；再将第三项系数增加 $\lambda x$，求环量的增量，并解释为何原积分无需逐边计算。
\end{exercise}
\begin{exercise}{597}{曲面选择与恰当微分}
设 $L$ 为 $y^2+z^2=1,x=y$ 的交线，从 $x$ 轴正向看为逆时针。用斯托克斯公式求 $\oint_Lx^2y^3\,\dd x+\dd y+z\,\dd z$；再利用 $x=y$ 将曲线上的微分式写为某个单变量函数的全微分，独立验证结果。
\end{exercise}
\begin{exercise}{598}{边界方向与逐边核验}
设 $L$ 为第一卦限平面三角形 $x+y+z=1$ 的边界，方向与上侧法向符合右手定则。用斯托克斯公式计算 $\oint_Lz\,\dd x+x\,\dd y+y\,\dd z$，写出三个顶点的正向顺序，再沿三条边直接积分核验。
\end{exercise}
\begin{exercise}{599}{六边形截面与旋度通量}
设 $L$ 为平面 $x+y+z=\frac32$ 与单位立方体 $[0,1]^3$ 表面的截痕，从 $x$ 轴正向看为逆时针。计算 $\oint_L(y^2-z^2)\,\dd x+(z^2-x^2)\,\dd y+(x^2-y^2)\,\dd z$；列出平面六边形的顶点及其面积，说明旋度与法向的内积为何在该截面上为常数。
\end{exercise}
\begin{exercise}{600}{球面截圆与参数}
设 $a>0,|c|<\sqrt3a$，$L_c$ 为 $x^2+y^2+z^2=a^2,x+y+z=c$ 的交圆，从 $x$ 轴正向看为逆时针。求 $\oint_{L_c}y\,\dd x+z\,\dd y+x\,\dd z$，给出 $c=0$ 时的值，并求其绝对值最大时的 $c$。
\end{exercise}
\begin{exercise}{601}{水平截圆的旋度积分}
设 $h>0$，$L_h$ 为 $x^2+y^2=2z,z=h$ 的交圆，从 $z$ 轴正向看为逆时针。求 $\oint_{L_h}3y\,\dd x-xz\,\dd y+yz^2\,\dd z$，给出 $h=2$ 时的值，并说明沿水平曲线为何第三项对环量没有贡献。
\end{exercise}
\begin{exercise}{602}{精确项与环量}
设 $L$ 为 $x^2+y^2=9,z=0$ 的圆周，从 $z$ 轴正向看为逆时针。计算 $\oint_L2y\,\dd x+3x\,\dd y-z^2\,\dd z$；若将 $-z^2\,\dd z$ 换成 $\dd\varphi$，其中 $\varphi\in C^1(\mathbb R^3)$，证明积分值不变，并指出闭曲线条件的作用。
\end{exercise}
\begin{exercise}{603}{旋度通量的换面}
设 $A=(y^2,xy,xz)$，$\Sigma$ 为单位上半球面上侧。将 $\iint_{\Sigma}(\nabla\times A)\cdot n\,\dd S$ 化为边界曲线积分并计算；再直接在单位圆盘上计算同一旋度通量，解释换面成立所需的方向与光滑性条件。
\end{exercise}
\begin{exercise}{604}{去底立方体的边界方向}
设 $A=(y-z,yz,-xz)$，$\Sigma$ 为 $[0,2]^3$ 的表面外侧去掉 $z=0$ 底面后的五个面。求 $\iint_{\Sigma}(\nabla\times A)\cdot n\,\dd S$，明确其边界从 $z$ 轴正向看是顺时针还是逆时针；利用闭曲面上旋度通量为零核验结果。
\end{exercise}
\begin{exercise}{605}{三角形环量与缩放}
设 $a>0$，$L$ 沿 $A(a,0,0)\to B(0,a,0)\to C(0,0,a)\to A$。求 $\oint_L(z-y)\,\dd x+(x-z)\,\dd y+(y-x)\,\dd z$；将三角形各点整体平移同一向量后，证明环量不变，并说明环量关于 $a$ 的次数。
\end{exercise}
\begin{exercise}{606}{斜椭圆的投影面积}
设 $a,b>0$，$L$ 为 $x^2+y^2=a^2,\ \frac xa+\frac zb=1$ 的交线，从 $x$ 轴正向看为逆时针。求 $\oint_L(y-z)\,\dd x+(z-x)\,\dd y+(x-y)\,\dd z$；写出与给定方向相容的单位法向量及椭圆面积，核验用投影面积计算的结果。
\end{exercise}
\begin{exercise}{607}{球柱交线与投影区域}
设 $L$ 为 $x^2+y^2+z^2=1,x^2+y^2=x,z\ge0$ 的闭交线，从 $x$ 轴正向看为逆时针。计算 $\oint_Ly^2\,\dd x+z^2\,\dd y+x^2\,\dd z$；以 $x=1-z^2$ 为所选曲面，写出其在 $yOz$ 面上的投影区域，说明该曲面与 $L$ 的边界方向相容。
\end{exercise}
\begin{exercise}{608}{力场做功与反向路径}
设 $F=(y,z,x)$，$L$ 为第一卦限内平面三角形 $x+y+z=1$ 的整个边界，从 $z$ 轴正向看为顺时针。求 $F$ 沿 $L$ 做的功；再沿反向边界计算，并说明若在 $F$ 上增加任意 $C^1$ 标量场的梯度，闭路功是否改变。
\end{exercise}
\begin{exercise}{609}{同边界曲面的定向与换面}
设 $A\in C^1(\mathbb R^3)$，$S_1$ 为单位上半球面上侧，$S_2$ 为单位圆盘 $z=0$ 上侧。证明 $\iint_{S_1}(\nabla\times A)\cdot n\,\dd S=\iint_{S_2}(\nabla\times A)\cdot n\,\dd S$；再说明把圆盘改取下侧时的等式，并用闭曲面边界方向核验符号。
\end{exercise}
\begin{exercise}{610}{截圆环量与参数优化}
设 $a>0,|h|<a$，$L_h:x^2+y^2+z^2=a^2,z=h$ 从 $z$ 轴正向看为逆时针。对 $A=(-yz,xz,0)$ 求环量 $I(h)$，求 $|I(h)|$ 的最大值及对应的 $h$；分别用圆盘旋度积分和曲线参数化核验。
\end{exercise}
\subsection{曲面积分的应用}
\begin{exercise}{611}{零散度与局部通量}
设 $a,h>0$，$A=(yz,xz,xy)$，$S$ 为圆柱体 $x^2+y^2\le a^2,0\le z\le h$ 的全表面外侧。求总通量及上底面通量；判断总通量为零能否推出每一块曲面上的通量均为零，并以该向量场在某个小平面片上的通量说明。
\end{exercise}
\begin{exercise}{612}{立方体通量与参数次数}
设 $a>0$，$A=(2x-z,x^2y,-xz^2)$。求它通过 $[0,a]^3$ 边界外侧的通量，整理为 $a$ 的多项式，并求使总通量为零的全部正数 $a$。
\end{exercise}
\begin{exercise}{613}{平移球面的常散度}
设 $A=(2x+3z,-xz-y,y^2+2z)$。求该场通过球心 $(3,-1,2)$、半径 $3$ 的球面外侧的通量；推广到任意球心、半径 $R>0$，说明哪些几何参数会影响结果。
\end{exercise}
\begin{exercise}{614}{半圆柱曲面的上侧}
设 $A=(0,yz,z^2)$，$\Sigma$ 为 $y^2+z^2=1,z\ge0,0\le x\le1$ 的上侧。直接求通量，并补齐半圆柱体的其余边界面核验；再将 $0\le x\le1$ 改为 $0\le x\le h$，求通量随 $h>0$ 的变化。
\end{exercise}
\begin{exercise}{615}{通量与边界分配}
设 $A=(x,y,z)$，求其通过单位立方体 $[0,1]^3$ 边界外侧的总通量及六个面的各自贡献；再对以原点为中心的单位边长立方体求各面通量，解释总量与分配的差别。
\end{exercise}
\begin{exercise}{616}{散度与闭曲面条件}
求 $A=(x^2+yz,y^2+xz,z^2+xy)$ 的散度，并求其散度为零的点集；若 $S$ 为以原点为球心的任意球面，判断 $A$ 的外通量是否为零，说明局部散度不恒为零为何不妨碍该结论。
\end{exercise}
\begin{exercise}{617}{复合函数偏导与散度}
求 $A=(e^{xy},\cos(xy),\cos(xz^2))$ 的散度及其在 $(0,1,0)$ 处的值；利用散度连续性，判断足够小且以该点为中心的球面外通量的符号，并说明理由。
\end{exercise}
\begin{exercise}{618}{散度与旋度的区分}
对 $A=(y^2,xy,xz)$ 分别求 $\nabla\cdot A$ 与 $\nabla\times A$，求二者同时为零的点集；再计算 $\nabla\cdot(\nabla\times A)$，解释这是标量还是向量。
\end{exercise}
\begin{exercise}{619}{散度与半球柱体通量}
对 $A=(0,yz,z^2)$ 求散度；利用所得结果求它通过半圆柱体 $0\le x\le1,y^2+z^2\le1,z\ge0$ 全边界的外通量，并说明为何该结果等于弯曲侧面的上侧通量。
\end{exercise}
\begin{exercise}{620}{点散度与小球通量极限}
设 $A=(xy^2,ye^z,x\ln(1+z^2))$，求其散度及在 $P(1,1,0)$ 处的值。以 $S_r$ 表示中心为 $P$、半径 $r$ 的球面外侧，求 $\lim_{r\to0^+}\frac{\iint_{S_r}A\cdot n\,\dd S}{4\pi r^3/3}$，并利用散度连续性证明。
\end{exercise}
\begin{exercise}{621}{旋度与环量的局部关系}
求 $A=(x^2-y,4z,x^2)$ 的旋度。设 $L$ 为 $z=0$ 平面上以原点为圆心、半径 $r>0$ 的圆周，从 $z$ 轴正向看为逆时针，利用旋度求 $\oint_L A\cdot\dd\boldsymbol r$，并直接参数化验证。
\end{exercise}
\begin{exercise}{622}{旋度消去与势函数}
求 $A=(z+\sin y,x\cos y-z,0)$ 的旋度，判断该场是否为梯度场；若增加第三分量 $x-y$，判断新场是否为梯度场，成立时求一个势函数。
\end{exercise}
\begin{exercise}{623}{复合函数的旋度}
求 $A=(x^2\sin y,y^2\sin(xz),xy\sin(\cos z))$ 的旋度及其在原点的值；再验证该场在整个平面 $z=0$ 上的旋度是否恒为零，须写出各分量。
\end{exercise}
\begin{exercise}{624}{点旋度与平面环量}
设 $A=(xy,-yz,xz)$。求 $\nabla\times A$ 在 $(1,1,0)$ 处的值；令 $L_r$ 为平面 $z=0$ 上以 $(1,1,0)$ 为圆心、半径 $r$ 的正向圆周，求 $\oint_{L_r}A\cdot\dd\boldsymbol r$，并解释其与该点旋度第三分量的关系。
\end{exercise}
\begin{exercise}{625}{径向函数的拉普拉斯算子}
在 $r=\sqrt{x^2+y^2+z^2}>0$ 上令 $u=\ln r$。直接计算 $\nabla u$ 与 $\nabla\cdot\nabla u$，求后者在 $(1,1,1)$ 处的值；判断 $u$ 在原点外是否为调和函数。
\end{exercise}
\begin{exercise}{626}{径向调和函数与边值}
设 $f\in C^2((0,\infty))$，$r=\sqrt{x^2+y^2+z^2}>0$。推导 $\nabla\cdot\nabla f(r)$ 的表达式，求满足其恒为零的全部 $f$；进一步确定满足 $f(1)=0,f(2)=1$ 的唯一解，并讨论它在 $r\to0^+$ 时是否有有限极限。
\end{exercise}
\begin{exercise}{627}{梯度场无旋的条件}
设 $u$ 在开集 $D\subset\mathbb R^3$ 上具有连续二阶偏导数，逐分量证明 $\nabla\times\nabla u=0$。若 $L$ 是 $D$ 内分段光滑闭曲线，证明 $\oint_L\nabla u\cdot\dd\boldsymbol r=0$，并说明此结论是否需要 $D$ 单连通。
\end{exercise}
\begin{exercise}{628}{散度旋度恒等式的正则条件}
设 $A=(P,Q,R)$ 在开集上具有连续二阶偏导数。逐项证明 $\nabla\cdot(\nabla\times A)=0$。若 $S$ 为一个有界区域的分片光滑边界，且闭区域也包含于该开集，证明旋度通过 $S$ 的外通量为零，指出此处需要哪些高斯公式条件。
\end{exercise}
\begin{exercise}{629}{径向调和函数与通量定常}
设 $u=f(r)$，$f\in C^2((1,2))\cap C^1([1,2])$，满足 $\nabla\cdot\nabla u=0$、$f(1)=1,f(2)=0$。求 $f$，计算 $\nabla u$ 通过任意球面 $r=\rho$（$1<\rho<2$）外侧的通量；解释该通量为何与 $\rho$ 无关，并核验球壳全边界的总通量。
\end{exercise}
\subsubsection{综合提高题}
\begin{exercise}{630}{切平面距离与投影约化}
设 $S$ 为椭球面 $\frac{x^2}2+\frac{y^2}2+z^2=1,z\ge0$。对 $P\in S$，记原点到该点切平面的距离为 $\rho(P)$。求 $\rho$ 的表达式，将 $\iint_S\frac z\rho\,\dd S$ 化为圆盘上的多项式二重积分并计算；说明赤道处用图像曲面面积公式时应怎样理解。
\end{exercise}
\begin{exercise}{631}{有限柱面与无界极限}
设 $R,H>0$，$\Sigma_H$ 为柱面 $x^2+y^2=R^2,0\le z\le H$ 的侧面。求 $I(H)=\iint_{\Sigma_H}\frac{\dd S}{x^2+y^2+z^2}$，再求 $H\to\infty$ 的极限；解释曲面面积趋于无穷是否必然导致该积分发散。
\end{exercise}
\begin{exercise}{632}{球面二次型与最小值}
设 $R>0$，$S$ 为原点球面 $x^2+y^2+z^2=R^2$。求 $\iint_S(ax+by+cz+d)^2\,\dd S$；当 $a^2+b^2+c^2=1$ 时，求该积分关于 $a,b,c,d$ 的最小值及全部等号条件。
\end{exercise}
\begin{exercise}{633}{偏心球面的坐标平移}
设 $a>0$，$S:x^2+y^2+z^2=2az$。确定球心及半径，利用平移和对称性求 $\iint_S(x^2+y^2+z^2)\,\dd S$；再求 $\iint_Sz\,\dd S$ 并用球面方程核验二者关系。
\end{exercise}
\begin{exercise}{634}{相交球面的球冠优化}
设 $a>0$，半径 $R>0$ 的球面 $S$ 的球心在定球面 $x^2+y^2+z^2=a^2$ 上。求 $S$ 位于定球内部部分的面积 $A(R)$，明确有交圆的 $R$ 范围及无交圆时的取值；求 $A(R)$ 的最大值及取到最大值的 $R$。
\end{exercise}
\begin{exercise}{635}{分段曲面函数与尺度}
设 $t>0$，$S_t:x^2+y^2+z^2=t^2$，在 $z\ge\sqrt{x^2+y^2}$ 时令 $f=x^2+y$，其余点令 $f=0$。求 $F(t)=\iint_{S_t}f\,\dd S$，并证明 $F(t)/t^4$ 为常数；说明分段边界对曲面积分是否有影响。
\end{exercise}
\begin{exercise}{636}{下半球取上侧的通量}
设 $a>0$，$S$ 为下半球面 $z=-\sqrt{a^2-x^2-y^2}$，取上侧。求 $\iint_S\frac{ax\,\dd y\,\dd z+(z+a)^2\,\dd x\,\dd y}{\sqrt{x^2+y^2+z^2}}$，分别写出两项的贡献；再说明该侧与球体外侧的关系，并给出外侧积分。
\end{exercise}
\begin{exercise}{637}{挖球法与径向通量}
设 $S$ 为椭球面 $2x^2+2y^2+z^2=4$ 的外侧。通过在内部挖去以原点为中心的小球，计算 $\iint_S\frac{x\,\dd y\,\dd z+y\,\dd z\,\dd x+z\,\dd x\,\dd y}{(x^2+y^2+z^2)^{3/2}}$；说明小球半径的允许范围，以及为何结果与椭球三个半轴的具体数值无关。
\end{exercise}
\begin{exercise}{638}{切平面垂直条件与曲面换元}
设 $S:x^2+y^2+z^2-yz=1$。求使切平面垂直于 $xOy$ 面的切点轨迹 $C$，将 $S$ 位于 $C$ 上方的部分明确写为 $z\ge y/2$，记为 $\Sigma$。计算
\[\iint_{\Sigma}\frac{(x+\sqrt3)|y-2z|}{\sqrt{4+y^2+z^2-4yz}}\,\dd S.\]
在投影计算中说明绝对值的处理及积分区域的椭圆换元。
\end{exercise}
\begin{exercise}{639}{抛物面补面与方向校验}
设 $\Sigma$ 为 $z=x^2+y^2,z\le1$ 的上侧。用高斯公式求向量场 $((x-1)^3,(y-1)^3,z-1)$ 的通量；明确曲面在闭合区域边界中应取何方向，再以散度积分与顶面通量解释最终符号。
\end{exercise}
\begin{exercise}{640}{任意闭曲面条件与微分方程}
设 $f\in C^1((0,\infty))$ 且 $\lim_{x\to0^+}f(x)=1$。已知半空间 $x>0$ 内任意有界分片光滑区域的边界外侧均满足
\[\iint_Sxf(x)\,\dd y\,\dd z-xyf(x)\,\dd z\,\dd x-ze^{2x}\,\dd x\,\dd y=0.\]
先由小球区域和散度连续性证明相应散度逐点为零，再求 $f$，并核验所求函数满足全部条件。
\end{exercise}
\begin{exercise}{641}{锥台侧面与切向场}
设 $S$ 为锥面 $z=\sqrt{x^2+y^2},1\le z\le2$ 的侧面外侧。求向量场 $(y,-x,z^2)$ 的通量；说明水平向量场 $(y,-x,0)$ 为何处处与锥面相切，再用补上两个圆盘的方法核验。
\end{exercise}
\begin{exercise}{642}{斜截柱面的环量}
设 $L$ 为 $x^2+y^2=1,x-y+z=2$ 的交线，从 $z$ 轴正向看为逆时针。求 $\oint_L(z-y)\,\dd x+(x-z)\,\dd y+(x-y)\,\dd z$；将平面常数 $2$ 改为 $c$，证明环量与 $c$ 无关并说明原因。
\end{exercise}
\begin{exercise}{643}{水平圆周与高度讨论}
设 $h>0$，$L_h$ 为 $z=x^2+y^2,z=h$ 的交圆，从 $z$ 轴正向看为逆时针。求 $\oint_{L_h}2y\,\dd x+xz\,\dd y-yz^2\,\dd z$，给出 $h=4$ 时的值，并求环量为零的全部正数 $h$。
\end{exercise}
\begin{exercise}{644}{平面约束化简微分式}
设 $L$ 为 $x^2+y^2=1,y+z=0$ 的交线，从 $z$ 轴正向看为逆时针。先利用 $z=-y$ 将 $\oint_Lz\,\dd x+y\,\dd z$ 化为平面曲线积分并求值，再以斯托克斯公式核验法向与符号。
\end{exercise}
\begin{exercise}{645}{斜圆盘上的旋度积分}
设 $L$ 为 $x^2+y^2=1,z=x+y$ 的交线，从 $z$ 轴正向看为逆时针。求 $\oint_Lxz\,\dd x+x\,\dd y+\frac{y^2}2\,\dd z$；把平面改为 $z=x+y+c$ 后，求环量关于 $c\in\mathbb R$ 的表达式。
\end{exercise}
\begin{exercise}{646}{菱形边界与分片光滑条件}
设 $L$ 为 $x+y+z=2,|x|+|y|=1$ 的交线，从 $z$ 轴正向看为逆时针。用斯托克斯公式求 $\oint_L(y^2-z^2)\,\dd x+(2z^2-x^2)\,\dd y+(3x^2-y^2)\,\dd z$，明确菱形投影区域；说明边界四个角点为何不妨碍公式使用。
\end{exercise}
\begin{exercise}{647}{开曲线补闭合与定向}
设 $L$ 为上半球面 $z=\sqrt{2-x^2-y^2}$ 与平面 $z=x$ 的交线，从 $A(0,\sqrt2,0)$ 沿 $x\ge0$ 的半椭圆到 $B(0,-\sqrt2,0)$。计算
\[\int_L(y+z)\,\dd x+(z^2-x^2+y)\,\dd y+x^2y^2\,\dd z.\]
补上 $B$ 到 $A$ 的直线段以使用斯托克斯公式，分别计算补线贡献，并明确闭路方向。
\end{exercise}
\begin{exercise}{648}{半椭球面与法向投影}
设 $\Sigma:x=\sqrt{1-3y^2-3z^2}$，取前侧，即法向的 $x$ 分量非负。求向量场 $(x,y^3+z,z^3)$ 的通量；利用补面法核验，并说明在 $x=0$ 圆盘上哪一项可能贡献通量。
\end{exercise}
\begin{exercise}{649}{零通量条件与旋转抛物面}
设 $f\in C^1((0,\infty))$ 且在 $0^+$ 有有限极限，定义
\[A=(xy-x^2y-xz^2,\ xy^2+2yf(z),\ zf(z)-yz).\]
已知上半空间内任意有界分片光滑区域边界的外通量为零。
\begin{enumerate}
\item 由此推导 $f$ 的微分方程，求 $f$ 并验证 $0^+$ 条件排除哪些解。
\item 将 $x=0,z=1+y^2,0\le y\le1$ 绕 $z$ 轴旋转所得曲面取上侧，求 $A$ 通过该曲面的通量，明确曲面方程与投影域。
\end{enumerate}
\end{exercise}
\begin{exercise}{650}{闭曲面是否包围奇点}
设 $0<a<1$，$S$ 为 $x^2+y^2-a\le z\le1$ 围成区域的整个边界外侧。计算 $\iint_S\frac{x\,\dd y\,\dd z+y\,\dd z\,\dd x+z\,\dd x\,\dd y}{(x^2+y^2+z^2)^{3/2}}$，说明原点的位置及挖球法的定向；再讨论将参数改为 $-1<a<0$ 时结果如何变化。
\end{exercise}
\begin{exercise}{651}{曲面约束下的分母替换}
将 $x=0,z=\sqrt{y+1},-1\le y\le1$ 绕 $y$ 轴旋转得到 $\Sigma$，取法向量与 $y$ 轴正向成钝角的一侧。先求曲面方程，利用该约束计算
\[\iint_{\Sigma}\frac{x^2\,\dd y\,\dd z+(1-y^3)\,\dd z\,\dd x+(2y^2z-z^2)\,\dd x\,\dd y}{y-x^2-z^2}.\]
说明能否在曲面上先用约束替换分母，再选取一个光滑向量场计算通量。
\end{exercise}
\begin{exercise}{652}{旋转双曲面与任意函数抵消}
将过 $A(1,0,0),B(1,1,1)$ 的直线绕 $z$ 轴旋转，截取 $0\le z\le1$ 的侧面并取外侧，记为 $\Sigma$。
\begin{enumerate}
\item 求曲面方程及一个外法向量。
\item 设 $f\in C^1(\mathbb R)$，求向量场 $(xf(xy)-2x,y^2-yf(xy),(z+1)^2)$ 通过 $\Sigma$ 的通量；指出含 $f$ 的部分为何不贡献总通量，并用补面计算核验。
\end{enumerate}
\end{exercise}
\begin{exercise}{653}{锥面上的径向相切项}
设 $\Sigma$ 为 $z=\sqrt{x^2+y^2},1\le x^2+y^2\le4$ 的下侧，$f\in C(\mathbb R)$。计算向量场 $(xf(xy)+2x-y,yf(xy)+2y+x,zf(xy)+z)$ 的通量；先从法向证明含 $f$ 的向量场与锥面相切，再解释连续而未必可导的条件为何足够。
\end{exercise}
\begin{exercise}{654}{椭球面边界与旋度消去}
设 $\Sigma$ 为 $4x^2+y^2+z^2=1$ 在第一卦限的部分，取上侧，边界 $L$ 按右手定则定向。求
\[\oint_L(yz^2-\cos z)\,\dd x+2xz^2\,\dd y+(2xyz+x\sin z)\,\dd z.\]
先把含三角函数的部分识别为恰当微分，计算其余部分的旋度，再将通量化为坐标平面上的积分。
\end{exercise}

\clearpage
\section{第13章\quad 无穷级数}
\subsection{常数项级数的概念与性质}
\begin{exercise}{655}{部分和与余项}
已知级数 $\sum_{n=1}^{\infty}u_n$ 的前 $n$ 项和为 $S_n=\frac{2n}{n+2}$。求 $u_1,u_2$ 及 $u_n$ 的表达式，按定义求级数的和 $S$，并求出使 $|S-S_n|<10^{-3}$ 的最小正整数 $n$。
\end{exercise}
\begin{exercise}{656}{裂项求和与误差}
设 $S_N=\sum_{n=1}^{N}\frac{3}{(3n-1)(3n+2)}$。求 $S_N$，据此证明相应无穷级数收敛并求其和；再求用 $S_N$ 近似该和时误差不超过 $10^{-3}$ 所需的最小 $N$。
\end{exercise}
\begin{exercise}{657}{对数级数的参数判别}
先求 $\sum_{n=1}^{N}\ln\left(1+\frac{2}{2n-1}\right)$ 并判定对应无穷级数的敛散性，再讨论实参数 $p$ 为何值时级数
\[\sum_{n=1}^{\infty}\frac{1}{n^p}\ln\left(1+\frac{2}{2n-1}\right)\]
收敛。说明临界参数处所用的判别依据。
\end{exercise}
\begin{exercise}{658}{定义与柯西准则}
对级数 $\sum_{n=0}^{\infty}(\sqrt{n+1}-\sqrt n)$，求前 $N+1$ 项和并判定敛散性；进一步计算 $\sum_{n=N}^{4N-1}(\sqrt{n+1}-\sqrt n)$，说明这一结果怎样反映该级数不满足柯西收敛准则。
\end{exercise}
\begin{exercise}{659}{根式裂项与绝对收敛}
对级数 $\sum_{n=1}^{\infty}(\sqrt[n+1]{2}-\sqrt[n]{2})$，求前 $N$ 项和并求无穷级数的和；利用通项符号判定其是否绝对收敛，再求使截断误差小于 $10^{-3}$ 的最小 $N$，可将答案保留为含对数的取整表达式。
\end{exercise}
\begin{exercise}{660}{三角数倒数与余项}
求 $\sum_{n=1}^{\infty}\frac1{1+2+\cdots+n}$ 的前 $N$ 项和及总和，给出精确余项；再讨论 $\sum_{n=1}^{\infty}\frac{n^p}{1+2+\cdots+n}$ 随实参数 $p$ 的敛散性。
\end{exercise}
\begin{exercise}{661}{有理化裂项与误差}
将 $\frac1{\sqrt{n+1}\sqrt n(\sqrt{n+1}+\sqrt n)}$ 写为两个相邻项之差，求相应级数的和及前 $N$ 项余项；求使误差不超过 $10^{-3}$ 的最小 $N$，并比较其与 $N^{-1/2}$ 的渐近关系。
\end{exercise}
\begin{exercise}{662}{等比级数的参数与和}
讨论实参数 $q$ 为何值时 $\sum_{n=1}^{\infty}(2-q)^n$ 收敛，求收敛时的和函数 $S(q)$；再求使 $S(q)=1$ 的全部 $q$，并验证所得参数属于收敛范围。
\end{exercise}
\begin{exercise}{663}{对数比值的几何级数}
先证明对任意 $a>1$ 都有 $0<\frac{\ln a}{a}<1$，再求 $S(a)=\sum_{n=1}^{\infty}(\frac{\ln a}{a})^n$，特别给出 $a=3$ 的值；求 $S(a)$ 在 $a>1$ 上的最大值及等号条件。
\end{exercise}
\begin{exercise}{664}{三角函数公比与临界角}
分别判定 $\sum_{n=0}^{\infty}\tan^n\theta$ 在 $\theta=\frac\pi6,\frac\pi4,\frac\pi3$ 时的敛散性，收敛时求和；再给出 $0<\theta<\frac\pi2$ 内的完整收敛范围，说明临界角处为何不能求和。
\end{exercise}
\begin{exercise}{665}{符号几何级数与绝对收敛}
对 $a>0$ 讨论 $\sum_{n=1}^{\infty}\frac{\cos n\pi}{a^n}$ 的敛散性与收敛类型，求收敛时的和；给出 $a=3$ 时的值及截取前 $N$ 项的精确余项。
\end{exercise}
\begin{exercise}{666}{两个幂次判据的交集}
求使 $\sum n^{-2p}$ 与 $\sum n^{-(3-p)}$ 同时收敛的全部实数 $p$，说明各临界点的情况；是否存在使两者同时发散的 $p$？证明结论。
\end{exercise}
\begin{exercise}{667}{必要条件与加括号的逆命题}
分别判断以下命题，成立者证明，不成立者给出通项明确的反例：
\begin{enumerate}
\item $\sum u_n$ 发散必推出 $u_n$ 不趋于零。
\item $\sum u_n,\sum v_n$ 均发散必推出 $\sum(2u_n-v_n)$ 发散。
\item $\sum u_n$ 发散必推出 $100+\sum_{n=1}^{\infty}u_{n+100}$ 发散。
\item $\sum u_n$ 发散必推出 $\sum_{n=1}^{\infty}(u_{2n-1}+u_{2n})$ 发散。
\end{enumerate}
对最后一项，补充 $u_n\to0$ 后结论是否改变？说明理由。
\end{exercise}
\begin{exercise}{668}{有限改变与逐项常数扰动}
已知 $\sum_{n=1}^{\infty}u_n=S$。逐一判定 $\sum1000u_n$、$\sum(u_n+\frac1{1000})$、$1000+\sum u_n$、$\sum u_{n+100}$ 的敛散性；收敛者用 $S$ 与必要的首项表示其和，发散者说明部分和趋向何处。
\end{exercise}
\begin{exercise}{669}{四类通项的判别依据}
分别判定下列级数的敛散性，给出各自最直接的判别依据，并求其中收敛级数的和：
\[\text{(1)}\ \sum_{n=1}^{\infty}\left(\frac1{2n}+\frac1{2^n}\right),\quad
\text{(2)}\ \sum_{n=1}^{\infty}\cos\frac1{n^2},\quad
\text{(3)}\ \sum_{n=1}^{\infty}\ln\left(1+\frac1n\right),\quad
\text{(4)}\ \sum_{n=1}^{\infty}(-1)^ne^{-n}.\]
对(3)另写出前 $N$ 项和，核验判定。
\end{exercise}
\begin{exercise}{670}{混合裂项的误差主项}
求 $\sum_{n=1}^{\infty}(3^{-n}+\frac3{(3n-1)(3n+2)})$ 的和及前 $N$ 项余项 $R_N$；求 $\lim_{N\to\infty}NR_N$，说明误差主要由哪一部分产生。
\end{exercise}
\begin{exercise}{671}{通项条件与参数}
讨论实参数 $a$ 为何值时级数
\[\sum_{n=1}^{\infty}\left(a\frac{n}{6n+1}+\frac{(-1)^n}{6^n}\right)\]
收敛，并在收敛时求和。对 $a=1$，再判断部分和趋于 $+\infty$、$-\infty$ 还是振荡，并说明理由。
\end{exercise}
\begin{exercise}{672}{级数线性运算}
已知实数项级数 $\sum_{n=1}^{\infty}u_n$ 收敛，$q\in\mathbb R\setminus\{0\}$。求使 $\sum_{n=1}^{\infty}(u_n+3q^{-n})$ 收敛的全部 $q$，并求两个级数的和之差。若原级数条件收敛，讨论新级数是否可能绝对收敛。
\end{exercise}
\begin{exercise}{673}{正项级数的临界指数}
讨论 $\sum_{n=1}^{\infty}(n^{1-p}+3\cdot2^{-n})$ 的敛散性，给出实参数 $p$ 的完整范围；分别说明在临界值处以及通项不趋于零的参数范围内，采用哪种判别方法即可完成判定。
\end{exercise}
\begin{exercise}{674}{线性组合与部分和}
已知 $\sum_{n=1}^{\infty}u_n=2$，且 $\sum_{k=1}^{n}v_k=\frac{2n}{3n+1}$。求 $v_n$ 及 $\sum_{n=1}^{\infty}(2u_n-3v_n)$；进一步求使 $\sum_{n=1}^{\infty}(a u_n-3v_n)$ 的和为零的实数 $a$，并说明逐项线性运算的依据。
\end{exercise}
\begin{exercise}{675}{奇偶分组与余项估计}
已知 $\sum_{n=1}^{\infty}n^{-2}=\frac{\pi^2}{6}$。求 $\sum_{n=1}^{\infty}(2n-1)^{-2}$ 及 $\sum_{n=1}^{\infty}(-1)^{n-1}n^{-2}$ 的和，并利用积分比较给出第一项级数截取前 $N$ 项后的余项上界。
\end{exercise}
\begin{exercise}{676}{收敛性质与反例}
设 $\sum u_n$ 收敛，$\sum v_n$ 发散。证明 $\sum(u_n+v_n)$ 必发散；分别构造满足上述条件且 $\sum u_nv_n$ 收敛、$\sum u_nv_n$ 发散的具体例子。各例均须写出通项并验证。
\end{exercise}
\begin{exercise}{677}{级数抵消与必要条件}
已知 $\sum(u_n+v_n)$ 收敛。证明 $\sum u_n$ 与 $\sum v_n$ 必同时收敛或同时发散；再构造一个 $u_n\to0,v_n\to0$ 且这两个级数均发散的例子，说明合并后的级数仍可收敛。
\end{exercise}
\begin{exercise}{678}{分组收敛与通项条件}
设 $\sum_{n=1}^{\infty}(u_{2n-1}+u_{2n})$ 收敛。证明原级数 $\sum u_n$ 收敛的充要条件为 $u_{2n}\to0$；举例说明若删去这一条件，分组级数可以收敛而原级数发散。
\end{exercise}
\begin{exercise}{679}{参数裂项与截断精度}
设 $a>0$，求 $\sum_{n=1}^{\infty}\frac{a}{(an+1)(a(n+1)+1)}$ 的和及前 $N$ 项余项，求使余项不超过 $\varepsilon>0$ 的最小正整数 $N$；讨论固定 $N$ 时余项随 $a$ 的变化。
\end{exercise}
\subsection{正项级数}
\begin{exercise}{680}{正项估计与拆分求和}
对 $\sum_{n=1}^{\infty}\frac{3+(-1)^{n+1}}{5^n}$，先给出两侧几何级数估计以证明收敛，再拆分求和；写出截取前 $N$ 项的精确余项，并验证它为正。
\end{exercise}
\begin{exercise}{681}{乘积项与几何上界}
利用几何级数比较判定下列级数的敛散性，并为各自从第 $N+1$ 项起的尾和给出明确上界：
\[\text{(1)}\ \sum_{n=1}^{\infty}\frac{3^{n-1}}{\prod_{k=1}^{n}(3k-2)},\qquad
\text{(2)}\ \sum_{n=1}^{\infty}\frac{2^n}{(2n-1)3^n}.\]
在(1)中明确给出相邻项比值及一个对所有 $n\ge1$ 成立的小于 $1$ 的上界。
\end{exercise}
\begin{exercise}{682}{不同增长阶的比较}
判定下列级数的敛散性，分别写出所选 $p$ 级数及通项比值的极限；对(3)进一步裂项求和：
\[\text{(1)}\ \sum_{n=1}^{\infty}\frac1{3n-1},\quad
\text{(2)}\ \sum_{n=1}^{\infty}\frac3{\sqrt n+\sqrt[3]n+\sqrt[4]n},\quad
\text{(3)}\ \sum_{n=1}^{\infty}\frac{n+3}{n(n+1)(n+2)}.\]
\end{exercise}
\begin{exercise}{683}{等价无穷小与尾和界}
利用极限比较判别下列级数，写出非零有限的比较极限：
\[\text{(1)}\ \sum_{n=1}^{\infty}\tan\frac\pi{3^n},\quad
\text{(2)}\ \sum_{n=1}^{\infty}\frac{3n^2+4}{n^4+1},\quad
\text{(3)}\ \sum_{n=1}^{\infty}\frac1{n\sqrt n}.\]
对(3)用积分比较给出前 $N$ 项余项的上下界，并说明误差随 $N$ 的阶。
\end{exercise}
\begin{exercise}{684}{三角函数中指数位置辨析}
分别判定 $\sum\sin\frac\pi6$、$\sum\sin\frac\pi{6n}$、$\sum\sin\frac\pi{6^n}$、$\sum\sin^n\frac\pi6$ 的敛散性，求第四个级数的和；对第三个级数给出几何级数尾和上界，解释这四种写法的通项极限或衰减速度。
\end{exercise}
\begin{exercise}{685}{对数等价与慢衰减}
判定 $\sum\ln(1+n^{-3})$、$\sum\ln(1+n^{-1/2})$、$\sum\frac1{\ln(3n+1)}$ 的敛散性，明确写出前两个的比较极限；再讨论 $\sum\ln(1+n^{-p})$ 随实参数 $p>0$ 的敛散性，处理临界值。
\end{exercise}
\begin{exercise}{686}{参数幂次与失效判据}
对 $p\in\mathbb R$ 讨论 $\sum_{n=1}^{\infty}n^{-p}\tan\frac1n$ 的敛散性，给出临界值处的比较极限；再求相邻项比值的极限，说明为什么比值判别法在本题不能区分参数范围。
\end{exercise}
\begin{exercise}{687}{已知上界能保证的收敛性}
设 $0\le u_n\le n^{-3}$。分别判定 $\sum u_n^3$、$\sum(u_n+n^{-3})$、$\sum(u_n+3^{-n})$、$\sum(u_n+\frac1{3n})$ 的敛散性；再求保证对所有满足条件的 $u_n$ 都有 $\sum n^p u_n$ 收敛的全部实数 $p$，并用具体例子证明范围不能扩大。
\end{exercise}
\begin{exercise}{688}{下界与渐近等价的作用}
设 $u_n>0$。
\begin{enumerate}
\item 若 $u_n\ge n^{-2}$，判断 $\sum nu_n$ 的敛散性，并说明能否据此判断 $\sum u_n$，不能时给出两种例子。
\item 若 $n^2u_n\to k>0$，判定 $\sum u_n$，并完整讨论 $\sum n^q u_n$ 随实参数 $q$ 的敛散性。
\end{enumerate}
\end{exercise}
\begin{exercise}{689}{平方改变收敛性的阈值}
设 $u_n=1-\cos\frac1{\sqrt n}$。求 $nu_n$ 的极限，判定 $\sum u_n$ 与 $\sum u_n^2$ 的敛散性；进一步讨论 $\sum u_n^p$ 对 $p>0$ 的收敛范围，说明取平方为什么改变判定。
\end{exercise}
\begin{exercise}{690}{阶乘与指数的比值判别}
分别用比值判别法处理 $\sum\frac1{(2n+1)!}$ 与 $\sum\frac{n!}{2^n}$，写出相邻项比值；对第二个级数，再将常数 $2$ 换为任意 $a>0$，说明是否可能通过选择 $a$ 使其收敛。
\end{exercise}
\begin{exercise}{691}{指数与阶乘的竞争}
判定 $\sum\frac{2^n}{n^2}$ 与 $\sum\frac{2^n n!}{n^n}$ 的敛散性，并求各自相邻项比值的极限；对收敛的级数，给出一个 $q<1$ 使充分大时尾项可由公比 $q$ 的几何级数控制。
\end{exercise}
\begin{exercise}{692}{奇数连乘与参数临界}
先用比值法判定 $\sum_{n=1}^{\infty}\frac{1\cdot3\cdots(2n-1)}{5^n n!}$；再将 $5$ 换为 $a>0$，求比值法能直接判定收敛或发散的参数范围。对临界值 $a=2$，利用
\[\frac{1\cdot3\cdots(2n-1)}{2^n n!}=\prod_{k=1}^{n}\left(1-\frac1{2k}\right)\ge\frac1{n+1}\]
判定敛散性，并以数学归纳法证明所用不等式。
\end{exercise}
\begin{exercise}{693}{指数衰减与和函数求和}
证明 $\sum_{n=1}^{\infty}\frac1{(n+1)e^n}$ 收敛；由 $-\ln(1-x)=\sum_{n=1}^{\infty}x^n/n$ 求其和，并给出用几何级数比较得到的一个前 $N$ 项余项上界。
\end{exercise}
\begin{exercise}{694}{两个指数的渐近比较}
对 $a>0$ 讨论 $\sum_{n=1}^{\infty}\frac{a^n}{7^n-2^n}$ 的敛散性，特别给出 $a=5$ 时的判定；说明在临界参数处为何必须检验通项，不能只引用比值极限为 $1$。
\end{exercise}
\begin{exercise}{695}{含参数系数的几何型级数}
设 $a>0$，讨论 $\sum_{n=1}^{\infty}\frac{an+1}{a^n}$ 的敛散性；在收敛范围内利用几何级数及其导数求和，并单独解释 $a=1$ 时通项的行为。
\end{exercise}
\begin{exercise}{696}{周期系数的拆分求和}
证明 $\sum_{n=1}^{\infty}\frac{3n+\sin^2(n\pi/3)}{3^n}$ 收敛；利用 $\sin^2(n\pi/3)$ 的周期性拆分为几何级数，求总和，并说明分组合法的依据。
\end{exercise}
\begin{exercise}{697}{递推通项与零初值例外}
设 $0\le a_n\le b_n$，且 $a_{n+1}=\frac12(1+\frac1n)^n a_n$。当 $a_1>0$ 时，判定 $\sum a_n$ 与 $\sum b_n$ 的敛散性；当 $a_1=0$ 时，说明是否仍可唯一判断 $\sum b_n$，并给出必要的反例。
\end{exercise}
\begin{exercise}{698}{混合级数与尾和估计}
对从 $n=1$ 开始的级数 $\sum_{n=1}^{\infty}(n^{-4}+\frac{4^n}{(n+1)!})$，分别判定两部分并证明总级数收敛；当 $N\ge7$ 时，利用后一部分相邻项比值不超过 $1/2$ 及积分比较，为总余项给出明确上界。
\end{exercise}
\begin{exercise}{699}{用级数收敛推出通项极限}
令 $u_n=\frac{(n!)^2}{(2n+1)!}$。先证明 $\sum_{n=1}^{\infty}u_n$ 收敛，由收敛必要条件求 $\lim u_n$；再找一个对所有 $n\ge1$ 成立的几何上界，以定量说明该极限。
\end{exercise}
\begin{exercise}{700}{根值法与统一几何控制}
对 $\sum_{n=1}^{\infty}(\frac{2n}{4n+1})^n$ 求根值极限并判定敛散性；进一步给出一个对全部 $n$ 成立的几何级数上界，并据此估计从第 $N+1$ 项起的尾和。
\end{exercise}
\begin{exercise}{701}{指数项的参数阈值}
讨论 $\sum_{n=1}^{\infty}\frac{a^n}{2^n(3n+1)}$ 随 $a>0$ 的敛散性，给出 $a=3$ 时的判定；对临界参数用比较法处理，并说明分母中的一次多项式为何不能压过指数增长。
\end{exercise}
\begin{exercise}{702}{根值判别与比较估计}
对下列两个级数，分别求通项的 $n$ 次方根的极限并判定敛散性。对收敛的级数，给出一个 $0<q<1$ 及整数 $N$，使其第 $n$ 项在 $n\ge N$ 时不超过 $q^n$。
\[\text{(1)}\quad\sum_{n=1}^{\infty}\left(\frac{4n-3}{3n+2}\right)^{2n};\qquad
\text{(2)}\quad\sum_{n=1}^{\infty}\left(1-\frac1n\right)^{n^2}.\]
\end{exercise}
\begin{exercise}{703}{指数项与多项式因子}
判定下列级数的敛散性，并说明乘上 $n^{-2}$ 这样的因子是否一定能使指数增长的通项趋于零。
\[\text{(1)}\quad\sum_{n=1}^{\infty}\frac{(\arctan n)^n}{2^n};\qquad
\text{(2)}\quad\sum_{n=1}^{\infty}\frac1{n^2}\left(1+\frac1n\right)^{n^2}.\]
对(1)给出一个可用于比较判别的几何级数。
\end{exercise}
\begin{exercise}{704}{根值判别的临界情形}
设 $a>0,p\in\mathbb R$。先判定 $\sum_{n=1}^{\infty}(\frac{na}{n+1})^n$ 的敛散性，再讨论
\[\sum_{n=1}^{\infty}\frac1{n^p}\left(\frac{na}{n+1}\right)^n\]
的敛散性。须单独处理根值判别法不能直接判定的参数值。
\end{exercise}
\begin{exercise}{705}{判别法失效与反例}
分别构造正项级数，使其相邻项比值极限均为 $1$，而一个收敛、一个发散；再证明若 $u_n>0$ 且 $n^2u_n\to L\in(0,\infty)$，则 $\sum u_n$ 收敛。说明极限比较与比值判别法所用的信息有何区别。
\end{exercise}
\begin{exercise}{706}{相消后的参数阈值}
对实参数 $p$ 讨论正项级数 $\sum_{n=1}^{\infty}n^p\left(\frac1n-\ln(1+\frac1n)\right)$ 的敛散性。先用定积分或二阶泰勒公式求括号内量与 $n^{-2}$ 的比值极限，证明正项性，并单独处理临界参数。
\end{exercise}
\subsection{任意项级数}
\begin{exercise}{707}{交错条件与绝对值比较}
分别判定 $\sum_{n=1}^{\infty}\frac{(-1)^n}{(2n)!}$ 与 $\sum_{n=1}^{\infty}\frac{(-1)^n}{\sqrt{n^2+a}+1}$（$a>0$）的收敛类型；对两者给出莱布尼茨余项界，说明参数 $a$ 是否改变第二个级数的收敛类型。
\end{exercise}
\begin{exercise}{708}{交错判别与通项条件}
对下列级数分别判定绝对收敛、条件收敛或发散，并指出莱布尼茨判别法的条件是否满足。
\[\text{(1)}\quad\sum_{n=1}^{\infty}(-1)^n\sin\frac1{\sqrt n};\qquad
\text{(2)}\quad\sum_{n=1}^{\infty}(-1)^n\cos\frac1{\sqrt n}.\]
对收敛的级数，估计取前 $N$ 项时的误差，并给出误差小于 $10^{-2}$ 的一个充分条件。
\end{exercise}
\begin{exercise}{709}{交错级数与绝对收敛}
判定下列级数的收敛类型，并明确验证绝对值通项的单调性及极限：
\[\text{(1)}\quad\sum_{n=1}^{\infty}(-1)^n\ln\left(1+\frac1n\right);\qquad
\text{(2)}\quad\sum_{n=2}^{\infty}\frac{(-1)^n}{\ln n}.\]
比较用莱布尼茨余项估计使误差小于 $10^{-3}$ 时，两者所需项数的数量级。
\end{exercise}
\begin{exercise}{710}{绝对收敛与符号拆分}
证明 $\sum_{n=1}^{\infty}(-1)^n\frac{2+(-1)^n}{4^n}$ 绝对收敛；把它拆为两个几何级数求和，并写出前 $N$ 项余项。说明本题按奇偶项分组求和的合法性。
\end{exercise}
\begin{exercise}{711}{交错形式与单调性缺失}
按所列顺序讨论级数 $\frac1{\sqrt2-1}-\frac1{\sqrt2+1}+\frac1{\sqrt3-1}-\frac1{\sqrt3+1}+\cdots$ 的敛散性。写出前 $2N$ 项和，说明通项虽趋于零且符号交替，为何仍不能用莱布尼茨判别法推出收敛；再判断奇数项部分和的极限行为。
\end{exercise}
\begin{exercise}{712}{幂指数与交错收敛类型}
分别判定 $\sum n^{-1/2}$、$\sum(-1)^nn^{-1/2}$、$\sum n^{-3/2}$、$\sum(-1)^nn^{-3/2}$ 的敛散性与收敛类型。证明对一般 $p>0$，正项级数 $\sum n^{-p}$ 与交错级数 $\sum(-1)^nn^{-p}$ 的收敛条件分别是什么，明确两个临界点的作用。
\end{exercise}
\begin{exercise}{713}{周期符号与稀疏交错项}
判定 $\sum_{n=1}^{\infty}\frac1n\sin\frac{n\pi}2$ 的收敛类型，写成去掉零项后的级数并求和；给出截取前 $4N$ 项的误差上界，说明插入零项为何不影响和。
\end{exercise}
\begin{exercise}{714}{交错判别的单调性证明}
判定 $\sum_{n=1}^{\infty}\frac{(-1)^n}{n-\ln n}$ 的收敛类型。证明分母始终为正、绝对值通项最终单调趋零，并给出前 $N$ 项误差界；绝对收敛部分须写出与调和级数的比较极限。
\end{exercise}
\begin{exercise}{715}{柯西不等式与最优上界}
设实数列满足 $\sum u_n^2<\infty$。证明 $\sum\frac{u_n}{n}$ 绝对收敛，并证明 $\sum\frac{|u_n|}{n}\le\frac\pi{\sqrt6}(\sum u_n^2)^{1/2}$；给出非零数列说明常数 $\pi/\sqrt6$ 不能减小。
\end{exercise}
\begin{exercise}{716}{平移交错级数的统一误差}
设 $x>0$，判定 $\sum_{n=1}^{\infty}\frac{(-1)^n}{n+x}$ 的收敛类型；证明对所有 $x>0$，前 $N$ 项误差均不超过 $1/(N+1)$，并据此给出不依赖 $x$ 的截断项数，使误差小于 $10^{-3}$。
\end{exercise}
\begin{exercise}{717}{有界振荡与几何控制}
判定 $\sum_{n=1}^{\infty}\frac{(-1)^n}{\pi^{n+1}}\cos\frac{n\pi}{n+1}$ 的收敛类型，给出绝对值级数的几何上界；再估计从第 $N+1$ 项起的尾和绝对值，说明无需检验余弦因子的单调性。
\end{exercise}
\begin{exercise}{718}{参数谐波项的分离}
求使 $\sum_{n=1}^{\infty}\frac{(-1)^n+kn}{n^2}$ 收敛的全部实数 $k$，说明收敛时的类型并求和，可用 $\sum n^{-2}=\pi^2/6$；对其余 $k$，判断部分和趋于何处。
\end{exercise}
\begin{exercise}{719}{四个交错级数的分类}
分别判定 $\sum(-1)^{n-1}n^{-2/3}$、$\sum(-1)^{n-1}n^{-3/2}$、$\sum(-1)^{n-1}(2/3)^n$、$\sum(-1)^{n-1}(3/2)^n$ 的收敛类型；对条件收敛者给出截断误差界，对几何型收敛者求和。
\end{exercise}
\begin{exercise}{720}{有界系数与绝对收敛}
分别判定下列级数的收敛类型，须给出绝对值比较或通项不趋零的依据：
\[\text{(1)}\ \sum\frac{(-1)^n(3+(-1)^n)}{(3n+1)\sqrt{n+3}},\qquad
\text{(2)}\ \sum\frac{(-1)^n2^n}{n^2+1},\]
\[\text{(3)}\ \sum\frac{(-1)^n\sqrt{n^2+1}}{n+1+\sin n},\qquad
\text{(4)}\ \sum\frac{(-1)^n}{3n-1}.\]
对(1)进一步说明若把分子中的 $3+(-1)^n$ 换成任意有界实数列，是否仍绝对收敛。
\end{exercise}
\begin{exercise}{721}{交错符号并非充分条件}
分别分类以下四个级数，并为每个给出关键极限或单调性论证：
\[\text{(1)}\ \sum(-1)^n\left(\frac n{n+1}\right)^n,\quad
\text{(2)}\ \sum(-1)^n3^n\sin\frac\pi{5^n},\]
\[\text{(3)}\ \sum(-1)^n\frac{n+2}{(n+1)\sqrt n},\quad
\text{(4)}\ \sum(-1)^n\frac{5^n-3^n}{5^n}.\]
对(3)给出前 $N$ 项的莱布尼茨误差界。
\end{exercise}
\begin{exercise}{722}{振荡系数与绝对值比值}
证明对每个 $x>0$，$\sum_{n=1}^{\infty}\frac{n^n}{3^nn!}\cos nx$ 均绝对收敛；对不含余弦因子的正项级数计算相邻项比值极限，并用同一尾和界同时控制全部 $x>0$ 时的截断误差。
\end{exercise}
\begin{exercise}{723}{平方可和与绝对收敛}
设 $k>0$，实数项级数 $\sum_{n=1}^{\infty}u_n^2$ 收敛。利用柯西不等式证明 $\sum_{n=1}^{\infty}\frac{u_n}{\sqrt{n^2+k}}$ 绝对收敛，并给出其绝对值级数的和关于 $\sum u_n^2$ 的一个上界。讨论把 $k>0$ 改为 $k=0$ 后结论是否仍成立。
\end{exercise}
\begin{exercise}{724}{交错级数与平方级数}
设 $u_n=(-1)^n\tan\frac1{\sqrt n}$。判定 $\sum u_n$ 的收敛类型及 $\sum u_n^2$ 的敛散性；进一步讨论正项级数 $\sum |u_n|^p$ 随实参数 $p>0$ 变化的敛散性，并处理临界值。
\end{exercise}
\begin{exercise}{725}{条件收敛与参数分解}
对实参数 $k$，讨论 $\sum_{n=1}^{\infty}\frac{k+n}{n^2}\cos n\pi$ 的收敛类型，并利用 $\sum n^{-2}=\frac{\pi^2}{6}$ 与 $\sum_{n=1}^{\infty}(-1)^{n-1}/n=\ln2$ 求其和。说明改变有限多个通项的符号是否影响收敛类型。
\end{exercise}
\begin{exercise}{726}{级数夹逼的证明}
设实数列满足 $u_n\le a_n\le v_n$，且 $\sum u_n,\sum v_n$ 都收敛。以非负级数 $\sum(a_n-u_n)$ 为出发点证明 $\sum a_n$ 收敛；再判断能否推出 $\sum|a_n|$ 收敛，若不能，给出满足全部条件的反例。
\end{exercise}
\begin{exercise}{727}{交错收敛的条件辨析}
设 $a_n>0,a_n\to0$。判断 $\sum(-1)^na_n$ 是否必收敛；构造明确数列使其发散，并指出该数列不满足的莱布尼茨条件。再证明若 $a_n$ 最终单调递减，则改变有限个首项不影响收敛性。
\end{exercise}
\begin{exercise}{728}{带正项扰动的交错级数}
对实参数 $p>0,q>0$ 讨论 $\sum_{n=1}^{\infty}((-1)^{n-1}n^{-p}+n^{-q})$ 的收敛、绝对收敛或条件收敛范围。须区分 $p=q$ 与 $p\ne q$，并证明绝对值级数的临界情况。
\end{exercise}
\subsection{幂级数}
\begin{exercise}{729}{收敛半径与边界通项}
求 $\sum_{n=0}^{\infty}\frac{(-1)^n(3n+1)}{3^{n+1}}x^n$ 的收敛半径及收敛域；在收敛区间内拆成几何级数及其导数求和，并说明两个边界点为何不能直接代入所得有理函数作为级数和。
\end{exercise}
\begin{exercise}{730}{两种指数系数的主导项}
求 $\sum_{n=1}^{\infty}\frac{3^n+(-2)^n}{n^2}x^n$ 的收敛半径、收敛域及端点处的收敛类型；再将 $n^2$ 改为 $n^p$（$p>0$），讨论两个端点的收敛性随 $p$ 的变化。
\end{exercise}
\begin{exercise}{731}{稀疏幂次的半径}
设 $a_n$ 最终非零且 $|a_{n+1}/a_n|\to4$。求 $\sum_{n=0}^{\infty}a_nx^{2n+1}$ 的收敛半径；说明能否仅凭该条件确定两个端点的收敛性，并分别给出端点都绝对收敛和端点都发散的系数例子。
\end{exercise}
\begin{exercise}{732}{条件收敛点决定半径}
已知 $\sum_{n=1}^{\infty}a_n(x+1)^n$ 在 $x=-2$ 条件收敛。证明其收敛半径等于 $1$，判断能否确定另一个端点的敛散性；若不能，分别给出另一个端点收敛与发散的系数例子，并写出相应完整收敛域。
\end{exercise}
\begin{exercise}{733}{阶乘系数的整函数级数}
求 $\sum_{n=1}^{\infty}\frac{n^2}{(2n)!}x^n$ 的收敛域；当 $x>0$ 时，令 $t=\sqrt x$，从 $\cosh t=\sum_{n=0}^{\infty}t^{2n}/(2n)!$ 出发求其和函数，并求在 $x=0$ 的连续取值。
\end{exercise}
\begin{exercise}{734}{临界端点与通项增长}
求 $\sum_{n=1}^{\infty}\frac{n!}{n^n}x^n$ 的收敛区间及收敛域。处理端点时，可利用 $(1+1/n)^n<e$ 证明 $e^nn!/n^n$ 随 $n$ 增加，进而检验通项是否趋零；无须使用斯特林公式。
\end{exercise}
\begin{exercise}{735}{端点收敛类型与求导}
求 $\sum_{n=1}^{\infty}\frac{(-1)^n}{2^n\sqrt n}x^n$ 的收敛域，分别判断端点的收敛类型；再求逐项求导后的幂级数的收敛域，说明求导保持半径是否也保持端点性质。
\end{exercise}
\begin{exercise}{736}{平移中心与导数级数}
求 $\sum_{n=1}^{\infty}\frac{(x-2)^n}{n^2+1}$ 的收敛域及端点收敛类型；求其导数级数的收敛半径，分别判定导数级数在两个端点的敛散性。
\end{exercise}
\begin{exercise}{737}{内点绝对收敛的定量证明}
已知 $\sum a_nx^n$ 在 $x=-3$ 收敛。证明它在 $x=2$ 绝对收敛；进一步，对任意 $0<r<3$，用 $|a_n(-3)^n|\le M$ 给出 $|x|\le r$ 上从第 $N+1$ 项起的绝对值尾和上界。
\end{exercise}
\begin{exercise}{738}{中心与距离判定}
设 $\sum_{n=1}^{\infty}a_n(x-1)^n$ 的收敛半径为 $1$。判定它在 $x=3$ 的敛散性，并说明在 $x=0,2$ 处能否确定；分别给出半径恰为 $1$ 且端点均收敛、均发散的系数例子。
\end{exercise}
\begin{exercise}{739}{两个等距端点确定收敛域}
已知 $\sum_{n=1}^{\infty}a_n(x+\frac12)^n$ 在 $x=1$ 收敛、在 $x=-2$ 发散。求收敛半径及完整收敛域，证明 $x=1$ 处必为条件收敛；给出满足要求的一个系数序列。
\end{exercise}
\begin{exercise}{740}{远处点的不可确定性}
已知 $\sum_{n=1}^{\infty}a_n(x-2)^n$ 在 $x=-1$ 收敛。说明能否唯一判断 $x=6$ 处的收敛性；分别构造在 $x=6$ 绝对收敛、条件收敛、发散的系数，并核验每个例子在 $x=-1$ 的条件。
\end{exercise}
\begin{exercise}{741}{逐项求导与端点}
已知 $\sum_{n=0}^{\infty}a_nx^n$ 的收敛半径为 $3$。求 $\sum_{n=1}^{\infty}na_n(x-2)^n$ 的收敛半径与收敛区间；说明能否仅凭已知条件判断其在 $x=-1,5$ 处的收敛性，并分别举例说明端点收敛与发散均可能发生。
\end{exercise}
\begin{exercise}{742}{几何级数与逐项积分}
求 $S(x)=\sum_{n=1}^{\infty}\frac{x^{n+1}}{3^n}$ 的收敛域及和函数，并说明在任意 $0<r<3$ 上逐项积分成立的理由，进而计算 $\int_0^r S(x)\,\dd x$。
\end{exercise}
\begin{exercise}{743}{望远镜求和与一致误差}
求 $S_N(x)=\sum_{n=1}^{N}(\frac{x^n}{n}-\frac{x^{n+1}}{n+1})$。据此求该级数在 $[-1,1]$ 上的和函数，并直接估计 $\sup_{|x|\le1}|S(x)-S_N(x)|$，判断是否一致收敛。
\end{exercise}
\begin{exercise}{744}{对数和函数与边界}
求 $\sum_{n=1}^{\infty}\frac{5^n}{n}x^n$ 的收敛域及和函数；单独判定两个边界点的收敛类型，并求在 $x=-\frac15$ 处级数的和，说明该值怎样与区间内部的和函数衔接。
\end{exercise}
\begin{exercise}{745}{幂级数求和与积分}
求 $S(x)=\sum_{n=1}^{\infty}\frac{(-1)^{n-1}}n x^{n+1}$ 的收敛域及和函数，分别处理两个端点；再利用和函数计算 $\sum_{n=1}^{\infty}\frac{(-1)^{n-1}}{n(n+2)}$，说明逐项积分所采用的区间与依据。
\end{exercise}
\begin{exercise}{746}{反正切级数与可去奇点}
求 $\sum_{n=0}^{\infty}\frac{(-1)^n}{2n+1}x^{2n}$ 的收敛域及和函数 $S(x)$，特别写明 $S(0)$；求 $S''(0)$，并利用端点处级数得到 $\sum_{n=0}^{\infty}\frac{(-1)^n}{2n+1}$ 的和。
\end{exercise}
\begin{exercise}{747}{逐项求导与积分应用}
从几何级数出发求 $S(x)=\sum_{n=1}^{\infty}nx^{2n-1}$ 的收敛域及和函数；求 $\int_0^{1/2}S(x)\,\dd x$，并说明为何可以在该积分区间内逐项积分。
\end{exercise}
\begin{exercise}{748}{展开中心与和函数}
求 $S(x)=\sum_{n=1}^{\infty}n(x+2)^n$ 的收敛域及和函数；在收敛区间内求 $S$ 的零点与单调区间，并判定级数在 $x=-3,-1$ 处是否收敛。
\end{exercise}
\begin{exercise}{749}{求导保持半径但改变边界}
对 $F(x)=\sum_{n=1}^{\infty}x^n/n^2$，分别求原级数、一次导数级数、二次导数级数的收敛域，证明三者收敛半径相同；说明为何不能把开区间内的逐项求导结论无条件推广到端点。
\end{exercise}
\begin{exercise}{750}{参数幂级数与可导区间}
设 $p>0$，$F_p(x)=\sum_{n=1}^{\infty}(x-1)^n/n^p$。求其收敛域及导数级数的收敛域，完整讨论两个端点随 $p$ 的变化；对任意 $0<r<1$，给出 $|x-1|\le r$ 上的几何控制，说明和函数在开区间内连续可导。
\end{exercise}
\subsection{函数展开成幂级数}
\begin{exercise}{751}{复合指数函数的系数}
将 $e^{-x^2}$ 展开成 $x$ 的幂级数并求收敛域，求其在原点的一般奇数阶与偶数阶导数；用展开至 $x^4$ 的多项式近似 $e^{-1/4}$，写出交错级数误差上界。
\end{exercise}
\begin{exercise}{752}{正弦展开与一般系数}
求 $\sin2x$ 的麦克劳林展开及 $x^7$ 的系数；对 $a\in\mathbb R$ 写出 $\sin(ax)$ 中 $x^{2m+1}$ 的系数，说明偶数次项为何全部为零，并给出展开的收敛域。
\end{exercise}
\begin{exercise}{753}{几何展开的常数项核验}
将 $1/(3-x)$ 展开成 $x$ 的幂级数，写出通项系数、收敛域及截断到 $x^N$ 的精确余项；通过 $x=0$ 的取值说明为何不能遗漏首项或把系数写成 $3^{-n}$。
\end{exercise}
\begin{exercise}{754}{三角恒等式与幂级数展开}
利用倍角公式将 $\sin^2x$ 展开成 $x$ 的幂级数，写出前四个非零项、通项及收敛域；求该函数在原点的 $2m$ 阶导数，并解释其奇数阶导数的取值。
\end{exercise}
\begin{exercise}{755}{平方根代换与实变量延拓}
求 $S(x)=\sum_{n=0}^{\infty}\frac{(-1)^nx^{n+1}}{(2n+1)!}$ 的收敛域。分别在 $x>0,x=0,x<0$ 写出实值和函数，可使用 $\sinh t=(e^t-e^{-t})/2$；说明原级数为何消除了平方根表达式在原点的表面限制。
\end{exercise}
\begin{exercise}{756}{逐项积分与端点求和}
从 $1/(1+x)=\sum_{n=0}^{\infty}(-1)^nx^n$ 出发推导 $\ln(1+x)$ 的展开，求完整收敛域；说明在 $x=1$ 与 $x=-1$ 两端的情况，再利用展开给出 $\ln\frac32$ 的三项近似及误差界。
\end{exercise}
\begin{exercise}{757}{泰勒展开与交错余项}
从 $\frac1{1+x^2}$ 的几何级数出发，推导 $\arctan x$ 的麦克劳林展开并给出收敛域；用前三个非零项近似 $\arctan\frac12$，写出严格误差上界。
\end{exercise}
\begin{exercise}{758}{有理函数的幂级数}
将 $f(x)=\frac1{(3-x)^2}$ 展开成 $x$ 的幂级数，给出通项系数及收敛域；利用系数求 $f^{(m)}(0)$，其中 $m$ 为非负整数，并说明两个端点均不在收敛域内的理由。
\end{exercise}
\begin{exercise}{759}{部分分式与收敛半径}
将 $f(x)=\frac1{x^2-6x+5}$ 展开成 $x$ 的幂级数，写出通项系数并求收敛域；说明半径为何由距离展开中心最近的实奇点决定，再求 $f^{(6)}(0)$。
\end{exercise}
\begin{exercise}{760}{平移中心的指数展开}
将 $e^{3x}$ 展开成 $(x-1)$ 的幂级数，说明其对全部实数成立；用该展开构造 $e^{3x}$ 在 $x=1$ 附近的三次泰勒多项式，并用拉格朗日余项估计其在 $|x-1|\le\frac1{10}$ 上的最大误差。
\end{exercise}
\begin{exercise}{761}{平移几何级数与余项}
将 $f(x)=\frac1{3-x}$ 展开成 $(x+1)$ 的幂级数并求收敛域；写出截断至 $(x+1)^N$ 的精确余项，并据此给出 $|x+1|\le2$ 上的统一误差上界。
\end{exercise}
\begin{exercise}{762}{对数的平移展开与端点}
利用 $\ln(1+t)$ 的展开推导 $\ln x$ 关于 $(x-2)$ 的幂级数，并求收敛域；说明 $x=0,4$ 两个端点的不同情形，再用展开的前两项近似 $\ln\frac52$ 并估计误差。
\end{exercise}
\begin{exercise}{763}{负中心展开与导数}
将 $x^{-2}$ 展开成 $(x+1)$ 的幂级数，明确收敛域及通项系数；由此求 $f^{(n)}(-1)$，并说明将 $x=-2$ 或 $x=0$ 代入该级数时为何不能得到有限的级数和。
\end{exercise}
\begin{exercise}{764}{部分分式与中心平移}
将 $f(x)=\frac{x}{x^2-3x+2}$ 展开成 $(x+2)$ 的幂级数，求其收敛域及通项系数；利用所得系数求 $f^{(4)}(-2)$，并单独验证收敛区间的两个端点。
\end{exercise}
\begin{exercise}{765}{泰勒系数与函数相等的条件}
设 $f$ 在原点邻域具有各阶导数。说明写出形式泰勒级数是否足以证明它等于 $f$。对 $f(x)=e^x$，用拉格朗日余项证明其麦克劳林展开在任意固定闭区间 $[-M,M]$ 上成立，并给出余项的统一上界。
\end{exercise}
\begin{exercise}{766}{有理函数的展开中心选择}
将 $f(x)=1/((1-x)(2-x))$ 分别展开为 $x$ 的幂级数与 $(x+1)$ 的幂级数，写出通项系数和收敛域；求两展开在 $x=0$ 的取值，说明展开中心改变为何会改变收敛区间却不改变函数值。
\end{exercise}
\subsection{幂级数的应用}
\begin{exercise}{767}{系数提取与一般导数}
设 $f(x)=2x^2/(1+x)$。用幂级数求 $f^{(5)}(0)$，再写出 $f^{(n)}(0)$ 对全部非负整数 $n$ 的表达式，分别处理 $n=0,1$；说明所用展开在哪个邻域成立。
\end{exercise}
\begin{exercise}{768}{乘积展开的奇偶阶导数}
设 $f(x)=(1+x)\arctan x$。求 $f^{(100)}(0),f^{(101)}(0)$，并给出一般偶数阶与奇数阶导数的公式；写明幂级数乘法中各阶系数来自哪些项。
\end{exercise}
\begin{exercise}{769}{指数乘积的两种求导方法}
设 $f(x)=xe^{-x}$。利用麦克劳林系数求 $f^{(n)}(0)$（$n\ge1$），再用莱布尼茨乘积求导公式核验；求 $\sum_{n=1}^{\infty}f^{(n)}(0)x^n/n!$ 的收敛域与和函数。
\end{exercise}
\begin{exercise}{770}{平移对数的高阶导数}
对 $f(x)=\ln(x^2+4x+5)$，写出在 $x=-2$ 附近的幂级数，求全部 $n\ge1$ 的 $f^{(n)}(-2)$；说明奇数阶导数为零的对称性依据，并给出展开的收敛域。
\end{exercise}
\begin{exercise}{771}{分式展开与导数恢复}
将 $f(x)=(x-1)/(2x+1)$ 展开成 $(x-1)$ 的幂级数，求其收敛域及 $f^{(n)}(1)$（$n\ge1$）；再写出截断至 $(x-1)^N$ 的精确余项，说明最接近展开中心的奇点在哪里。
\end{exercise}
\begin{exercise}{772}{积分几何级数求数项和}
构造 $F(x)=\sum_{n=1}^{\infty}x^n/(n+1)$，先求其收敛域与和函数，再求 $\sum_{n=1}^{\infty}1/(2^n(n+1))$ 的和；在 $x=0$ 单独写出和函数的连续值。
\end{exercise}
\begin{exercise}{773}{线性加权几何级数}
求 $F(x)=\sum_{n=0}^{\infty}(n+1)x^n$ 的收敛域及和函数，进而求 $\sum_{n=0}^{\infty}(n+1)/3^n$；通过对有限几何和求导，给出后一数项级数前 $N+1$ 项的精确余项。
\end{exercise}
\begin{exercise}{774}{积分求和与端点极限}
设 $F(x)=\sum_{n=0}^{\infty}\frac{(-1)^n x^{3n+1}}{3n+1}$。先在 $|x|<1$ 内求 $F'(x)$ 并表示 $F(x)$；再说明能否令 $x\to1^-$ 求出 $\sum_{n=0}^{\infty}\frac{(-1)^n}{3n+1}$，并将其和写成初等常数。
\end{exercise}
\begin{exercise}{775}{三角级数与加权求和}
从 $\cos x$ 的麦克劳林展开出发，求 $F(x)=\sum_{n=0}^{\infty}\frac{(-1)^n(2n+1)}{(2n)!}x^{2n}$ 的和函数及收敛域，进而求 $F(1)$；说明使用逐项求导的依据。
\end{exercise}
\begin{exercise}{776}{二次加权幂级数}
求 $S(x)=\sum_{n=1}^{\infty}\frac{n(n+1)}{2^n}x^n$ 的收敛域及和函数，进而求 $\sum_{n=1}^{\infty}\frac{(-1)^n n(n+1)}{2^n}$ 的和；再用同一和函数求 $\sum_{n=1}^{\infty}\frac{n(n+1)}{4^n}$。
\end{exercise}
\begin{exercise}{777}{端点代入与逐项求导辨析}
对 $F(x)=\sum_{n=1}^{\infty}(-1)^{n-1}x^n/n$，求收敛域与和函数。判断能否在 $x=1$ 直接把导数级数求和得到 $F$ 在该点的左导数；写出导数级数的通项行为，并说明它与 $\ln(1+x)$ 在该点可导之间为何不矛盾。
\end{exercise}
\begin{exercise}{778}{高阶系数与交错数项和}
设 $f(x)=\ln(1+x^2)$。求全部 $f^{(n)}(0)$，再构造幂级数求 $\sum_{n=1}^{\infty}(-1)^{n-1}/(n2^n)$ 及 $\sum_{n=1}^{\infty}(-1)^{n-1}/(n(n+1)2^{n+1})$ 的和。明确逐项积分区间及合法性。
\end{exercise}
\subsection{傅里叶级数}
\begin{exercise}{779}{非对称周期段与常数项}
设 $f$ 为 $2\pi$ 周期函数，在 $0<x\le2\pi$ 上 $f(x)=x^2$。求傅里叶系数 $a_0$ 与 $a_n$（$n\ge1$），并求傅里叶级数在 $x=0,2\pi$ 处的和；说明积分区间可否从 $[0,2\pi]$ 改为任意长度 $2\pi$ 的区间。
\end{exercise}
\begin{exercise}{780}{周期归约与跳跃平均}
将 $f(x)=x^3+1$（$-\pi\le x<\pi$）作 $2\pi$ 周期延拓。求其傅里叶级数在 $x=7\pi/2,\pi,0$ 处的和，并写出全体间断点处的和；解释在间断点改变 $f$ 的单点取值为何不改变傅里叶系数。
\end{exercise}
\begin{exercise}{781}{分段周期函数的和函数}
设 $f$ 为 $2\pi$ 周期函数，在 $[-\pi,\pi)$ 上为 $f(x)=x$（$x<0$），$f(x)=0$（$x\ge0$）。求其傅里叶级数的和函数在 $[-\pi,\pi]$ 上的表达式，给出 $S(3\pi)$，并求常数项 $a_0/2$。
\end{exercise}
\begin{exercise}{782}{正弦展开的奇延拓}
设 $f(x)=|x-\pi/2|$（$0\le x\le\pi$），$b_n=\frac2\pi\int_0^\pi f(x)\sin(nx)\,\dd x$。记 $S(x)=\sum_{n=1}^{\infty}b_n\sin(nx)$。说明对应奇延拓的周期与间断点，求 $S(-9\pi/4),S(0),S(\pi)$；解释端点处级数和与原函数端点值的区别。
\end{exercise}
\begin{exercise}{783}{分段线性函数的完整展开}
将 $f(x)=2x$（$-\pi\le x<0$）、$f(x)=4x$（$0\le x<\pi$）作 $2\pi$ 周期延拓。先分解为奇函数与偶函数之和，求全部傅里叶系数及级数的和函数 $S(x)$；明确 $x=(2k+1)\pi$ 与 $x=2k\pi$ 处的和，$k\in\mathbb Z$。
\end{exercise}
\begin{exercise}{784}{周期延拓与间断点取值}
设 $f$ 为周期 $2$ 的函数，在 $[-1,1)$ 上为
\[f(x)=\begin{cases}1-3x,&-1\le x<0,\\x^2+1,&0\le x<1.\end{cases}\]
记其傅里叶级数的和函数为 $S$。不计算全部傅里叶系数，求 $S(-7),S(0),S(\frac12)$，并写出所有满足 $S(x)\ne f(x)$ 的点。说明所用收敛定理的条件。
\end{exercise}
\begin{exercise}{785}{余弦延拓与跳跃}
设 $f(x)=2x$（$0\le x\le\frac12$），$f(x)=x^2+1$（$\frac12<x\le1$），并令
\[a_n=2\int_0^1f(x)\cos(n\pi x)\,\dd x\quad(n=0,1,\ldots),\qquad
S(x)=\frac{a_0}2+\sum_{n=1}^{\infty}a_n\cos(n\pi x).\]
说明该余弦级数对应的偶延拓与周期，求 $S(-\frac94),S(\frac12),S(1)$，并求常数项 $a_0/2$。
\end{exercise}
\begin{exercise}{786}{正弦延拓的边界值}
设 $f(x)=2x^2$（$0\le x<1$），$b_n=2\int_0^1f(x)\sin(n\pi x)\,\dd x$，$S(x)=\sum_{n=1}^{\infty}b_n\sin(n\pi x)$。说明相应奇延拓的周期与跳跃点，求 $S(-\frac12)+S(9)+S(10)$；再求 $b_n$ 并用它验证所有整数点处的和。
\end{exercise}
\begin{exercise}{787}{傅里叶恒等式的取点求和}
已知在 $0\le x\le2$ 上成立
\[(x-1)^2=\frac13+\frac4{\pi^2}\sum_{n=1}^{\infty}\frac{\cos(n\pi x)}{n^2}.\]
选择适当的 $x$ 求 $\sum n^{-2}$ 与 $\sum(-1)^{n-1}n^{-2}$，再求 $\sum_{n=1}^{\infty}(2n-1)^{-2}$。每次代入均须写明所用的点及级数收敛的理由。
\end{exercise}
\begin{exercise}{788}{余弦展开与周期端点}
将 $f(x)=x^2-x$（$0\le x\le1$）展开成周期为 $2$ 的余弦级数，写出全部傅里叶系数及在端点处的和；利用 $x=\frac12$ 处的展开求 $\sum_{m=1}^{\infty}\frac{(-1)^m}{m^2}$。
\end{exercise}
\begin{exercise}{789}{单点赋值与傅里叶和}
设 $f$ 为周期 $2$ 的函数，在 $(-1,0)$ 上等于 $-1$，在 $(0,1)$ 上等于 $1$，并令 $f(0)=c,f(1)=f(-1)=d$。证明傅里叶系数与 $c,d$ 无关，求傅里叶级数在全部整数点处的和，并确定使该和等于函数值的 $c,d$。
\end{exercise}
\begin{exercise}{790}{余弦展开与参数级数恒等式}
将 $f_a(x)=x^2+ax$（$0\le x\le1$）作偶延拓后再作周期 $2$ 延拓，求其余弦级数。确定使所有奇数阶余弦系数消失的 $a$，再利用 $x=0,1/2$ 的取值求出相应的两个数项级数恒等式，明确端点收敛依据。
\end{exercise}
\subsubsection{综合提高题}
\begin{exercise}{791}{通项趋零与复合极限}
已知 $\sum_{n=1}^{\infty}(u_n-\frac{3n}{n+1})$ 收敛。求 $u_n$ 的极限及 $\lim_{n\to\infty}(u_n^2-2u_n+3)$；判断能否据此断言 $\sum(u_n-3)$ 收敛，并证明你的结论。
\end{exercise}
\begin{exercise}{792}{错位级数与首项修正}
已知 $\sum_{k=1}^{n}u_k=\frac{n}{2n+1}$。求 $u_n$ 及 $\sum_{n=1}^{\infty}(u_n+u_{n+1}-u_{n+2})$ 的和；进一步求其前 $N$ 项和与总和之差，说明错位相加时哪些首项需要单独保留。
\end{exercise}
\begin{exercise}{793}{三重分母裂项与精确余项}
按级数收敛定义求 $\sum_{n=1}^{\infty}\frac1{n(n+1)(n+2)}$ 的和，写出前 $N$ 项和及精确余项；证明余项是 $N^{-2}$ 阶，并求 $N^2R_N$ 的极限。
\end{exercise}
\begin{exercise}{794}{幂和恒等式与级数求和}
先用数学归纳法证明 $1^3+2^3+\cdots+n^3=(n(n+1)/2)^2$，再求 $\sum_{n=1}^{\infty}\frac1{\sqrt{1^3+2^3+\cdots+n^3}}$ 的和；写出用前 $N$ 项近似时的精确误差。
\end{exercise}
\begin{exercise}{795}{通项收敛与部分和有界的区别}
设 $u_n>0$，$S_n=\sum_{k=1}^{n}u_k$。证明 $\{S_n\}$ 有界是 $\{u_n\}$ 收敛的充分条件，说明该条件是否必要并给出反例；再证明 $\{S_n\}$ 有界与正项级数 $\sum u_n$ 收敛等价，区分两种收敛对象。
\end{exercise}
\begin{exercise}{796}{单调部分和的倒数比较}
设 $u_n>0$，$S_n=\sum_{k=1}^{n}u_k\to\infty$。证明对 $n\ge2$ 有 $\frac{u_n}{S_n^2}\le\frac1{S_{n-1}}-\frac1{S_n}$，据此证明 $\sum\frac{u_n}{S_n^2}$ 收敛并给出关于 $u_1$ 的上界；讨论证明是否实际需要原级数发散。
\end{exercise}
\begin{exercise}{797}{抵消精度与参数辨析}
讨论 $\sum_{n=1}^{\infty}(\frac1{\sqrt n+1}-\frac1{\sqrt n}+\frac cn)$ 随 $c\in\mathbb R$ 的敛散性，给出 $c=1$ 时的收敛类型；在收敛情形下将通项化为单个正分式并求其与 $n^{-3/2}$ 的比值极限。
\end{exercise}
\begin{exercise}{798}{积分定义通项与尾和控制}
设 $u_n=\int_n^{n+1}\frac{e^{-x}}{1+\sin^2n}\,\dd x$。证明 $\sum u_n$ 收敛，给出级数和的两个明确常数界及前 $N$ 项余项的指数上界；说明分母中的振荡为何不改变收敛性。
\end{exercise}
\begin{exercise}{799}{一阶与二阶无穷小的级数差别}
分别判定 $\sum(e^{1/n}-1)$、$\sum\ln(1+1/n)$、$\sum\sin(1/n)$、$\sum(1-\cos(1/n))$ 的敛散性，写出对应等价无穷小；再求使 $\sum(1-\cos(1/n))^p$ 收敛的全部 $p>0$。
\end{exercise}
\begin{exercise}{800}{极限指数与幂次阈值}
先判定 $\sum_{n=1}^{\infty}\frac{(n+2)^n}{n^{n+2}}$，再对 $c>-1,p\in\mathbb R$ 讨论 $\sum_{n=1}^{\infty}\frac{(n+c)^n}{n^{n+p}}$ 的敛散性；写出与 $p$ 级数的比较极限，单独解释临界值。
\end{exercise}
\begin{exercise}{801}{交错分母产生的非交错余项}
判定 $\sum_{n=1}^{\infty}\frac{(-1)^{n+1}}{\sqrt n+(-1)^{n+1}}$ 的敛散性。对 $n\ge2$ 有理化分母，将通项分为一个交错部分与一个同号部分，分别判定；若发散，说明部分和趋于何处，不能仅凭符号交替作答。
\end{exercise}
\begin{exercise}{802}{参数分母与隐藏调和项}
设 $k>1$。把 $\frac{(-1)^n}{\sqrt{kn}-(-1)^n}$ 有理化，判定相应级数的敛散性并说明部分和的极限行为；解释参数 $k$ 是否改变结论，以及 $k>1$ 如何保证全部通项有定义。
\end{exercise}
\begin{exercise}{803}{整数相位与小量振荡}
设 $k>0$，判定 $\sum_{n=1}^{\infty}\sin(\pi\sqrt{n^2+k})$ 的收敛类型。令 $\delta_n=\pi(\sqrt{n^2+k}-n)$，证明 $\delta_n$ 单调趋于零并求 $n\delta_n$ 的极限，再判断 $k=0$ 时的情形。
\end{exercise}
\begin{exercise}{804}{奇偶子级数的联立关系}
设 $u_n>0$，$\sum u_n$ 发散而 $\sum(-1)^{n-1}u_n$ 收敛。证明奇数项子级数与偶数项子级数都发散，判定 $\sum(u_{2n-1}+u_{2n})$ 与 $\sum(u_{2n-1}-u_{2n})$ 的敛散性；用各自前 $N$ 项和说明所有结论。
\end{exercise}
\begin{exercise}{805}{分式变量代换与收敛域}
求 $\sum_{n=1}^{\infty}\frac1{2n+1}(\frac{1-x}{1+x})^n$ 的实收敛域，明确排除无定义点并检验全部边界。对任意 $0<a<b<\infty$，给出 $[a,b]$ 上的几何尾和界，说明其和函数在 $(0,\infty)$ 上连续。
\end{exercise}
\begin{exercise}{806}{收敛域的平移反推}
已知 $\sum_{n=1}^{\infty}\frac{(-1)^{n-1}}n(x-a)^n$ 对每个 $x>0$ 都发散，而在 $x=0$ 收敛。求 $a$，写出完整收敛域及和函数，并说明 $x=0$ 处是绝对收敛还是条件收敛。
\end{exercise}
\begin{exercise}{807}{单调系数与端点信息}
设实数列 $a_n$ 单调递减趋于零，且 $S_n=\sum_{k=1}^{n}a_k$ 无界。求 $\sum_{n=1}^{\infty}a_n(x-2)^n$ 的收敛半径与收敛域；证明各 $a_n$ 非负，分别解释内点、左端点和右端点的收敛类型。
\end{exercise}
\begin{exercise}{808}{二次分母拆分求和}
求 $S(x)=\sum_{n=2}^{\infty}\frac{x^n}{n^2-1}$ 的完整收敛域及和函数，单独写明 $x=0$ 的连续取值；求 $S(1/2)$，并由端点级数或极限求 $S(1)$。
\end{exercise}
\begin{exercise}{809}{阶乘加权和与可去点}
求 $S(x)=\sum_{n=1}^{\infty}\frac n{(n+1)!}x^n$ 的收敛域及和函数，明确 $x=0$ 的取值；求 $S(1)$，再由展开系数求 $S'(0)$，核验所得分式表达式在原点的可去性。
\end{exercise}
\begin{exercise}{810}{微分方程的多项式解与求和}
对每个正整数 $n$，求满足 $xy_n'-2ny_n=0$ 且 $y_n(1)=\frac1{n(2n-1)}$ 的多项式 $y_n(x)$，并证明在多项式类中的唯一性。求 $\sum_{n=1}^{\infty}(-1)^{n-1}y_n(x)$ 的收敛域及和函数，分别处理 $x=0$ 和两个端点。
\end{exercise}
\begin{exercise}{811}{反常面积与对数求和}
设 $D_n$ 为 $x\ge e^3,0\le y\le\frac1{x(\ln x)^{n+1}}$ 的无界区域，$n\ge1$。求其面积 $A_n$，证明 $\sum(-1)^{n-1}A_n$ 收敛并求和；再将 $e^3$ 改为 $e^b$（$b>0$），讨论所得面积级数随 $b$ 的敛散性。
\end{exercise}
\begin{exercise}{812}{两个对数的平移展开}
利用 $\ln(1+t)$ 的展开将 $\ln\frac{x}{1+x}$ 展开成 $(x-1)$ 的幂级数，写出通项系数与收敛域；单独判断两个端点，并说明它们是否属于原函数的定义域。
\end{exercise}
\begin{exercise}{813}{反正切恒等式的分支与端点}
设 $f(x)=x\arctan\frac{x+2}{x-2}$，其中 $\arctan$ 取值于 $(-\pi/2,\pi/2)$。先在 $-2<x<2$ 推导适用的反正切恒等式，再展开成 $x$ 的幂级数；分别求级数收敛域及展开等于原函数的范围，特别解释 $x=2$ 处级数收敛与原函数无定义的关系。
\end{exercise}
\begin{exercise}{814}{乘积展开的有限阶系数}
设 $f(x)=xe^{-x^2}\sin2x$。不反复求导，求 $f^{(4)}(0)$ 与 $f^{(6)}(0)$；说明提取这两阶系数各需将两个因子展开到什么阶，并证明全部奇数阶原点导数为零。
\end{exercise}
\begin{exercise}{815}{柯西乘积与一般导数}
设 $f(x)=(\sum_{n=1}^{\infty}x^n)^2$，$|x|<1$。通过柯西乘积求 $x^n$ 的系数及 $f^{(n)}(0)$（$n\ge2$），写出 $n=0,1$ 时的导数值；再用有理函数形式验证，并说明级数乘法合法的依据。
\end{exercise}
\begin{exercise}{816}{方波的跳跃平均与系数}
设 $f(x)=\operatorname{sgn}(\cos x)$，并规定 $\operatorname{sgn}(0)=0$。求其傅里叶级数和函数 $S$，计算 $S(-\pi/2)+S(3\pi)$；利用偶对称性求全部傅里叶系数，说明在跳跃点的取值依据。
\end{exercise}
\begin{exercise}{817}{余弦系数的奇偶分离}
将 $f(x)=1-x$（$0\le x\le1$）作偶延拓并以 $2$ 为周期延拓，记余弦系数为 $a_n$。求全部 $a_n$，求 $\sum_{n=1}^{\infty}a_{2n}$ 与 $\sum_{n=1}^{\infty}a_{2n-1}$；再用 $x=0,1$ 处的傅里叶和核验这两个数项和。
\end{exercise}
\begin{exercise}{818}{分段周期函数与系数总和}
设 $f$ 为周期 $2$ 的函数，在 $-1<x\le0$ 上 $f(x)=x^2$，在 $0<x\le1$ 上 $f(x)=x-1$。令 $a_n=\int_{-1}^{1}f(x)\cos(n\pi x)\,\dd x$（$n\ge0$）。求傅里叶级数的和函数在 $[-1,1]$ 上的完整表达式，并求 $\sum_{n=0}^{\infty}a_n$ 与 $\sum_{n=0}^{\infty}(-1)^na_n$；说明含 $a_0$ 的数项和与傅里叶常数项 $a_0/2$ 之间的差别。
\end{exercise}

\end{document}
